\pdfoutput=1
\documentclass[11pt]{amsart}
\usepackage{amsmath,amssymb,amsthm,mathtools,booktabs,microtype}
\usepackage[margin=1.2in]{geometry}
\usepackage{xcolor}
\usepackage[colorlinks=true,linkcolor=blue,citecolor=blue,urlcolor=blue]{hyperref}
\hypersetup{pdfauthor={Bernd Johannes Wuebben},pdftitle={Triple curves and unbounded denominators in threefold curve counting}}
\setlength{\emergencystretch}{2em}
\newtheorem{theorem}{Theorem}[section]
\newtheorem{proposition}[theorem]{Proposition}
\newtheorem{lemma}[theorem]{Lemma}
\newtheorem{corollary}[theorem]{Corollary}
\theoremstyle{definition}
\newtheorem{definition}[theorem]{Definition}
\newtheorem{example}[theorem]{Example}
\theoremstyle{remark}
\newtheorem{remark}[theorem]{Remark}
\newcommand{\ZZ}{\mathbb Z}
\newcommand{\QQ}{\mathbb Q}
\newcommand{\CC}{\mathbb C}
\newcommand{\PP}{\mathbb P}
\newcommand{\RR}{\mathbb R}
\newcommand{\FF}{\mathbb F}
\newcommand{\cO}{\mathcal O}
\DeclareMathOperator{\Tot}{Tot}
\DeclareMathOperator{\Pic}{Pic}
\DeclareMathOperator{\Sym}{Sym}
\DeclareMathOperator{\tr}{tr}
\DeclareMathOperator{\Bl}{Bl}
\DeclareMathOperator{\conv}{conv}
\title[Triple curves and unbounded denominators]{Triple curves and unbounded denominators in threefold curve counting}
\author{Bernd Johannes Wuebben}
\subjclass[2020]{14N35 (primary); 14C25, 14H30, 14H60 (secondary)}
\keywords{Local curves, universal curve counting, Gromov--Witten invariants, triple covers, integral cycle algebra}
\date{November 24, 2026\\This version: November 24, 2026}
\begin{document}
\begin{abstract}
We exhibit a smooth projective Calabi--Yau threefold on which the
exponents of the scalar local factors in the universal series of
effective curves have unbounded denominators, even among homology
classes of divisibility three. These factors are defined by scaling
both normal directions to a curve with the same weight. For a curve of
odd genus $2k+1$, the scalar class of its tripled zero section has
exact denominator $3^k$ in the integral cycle algebra in Chern number
and virtual dimension zero. We also determine the denominator in even
genus. From even genus four onward it is smaller than the least
product of weights needed to isolate the triple by single-fiber
evaluations, even allowing the normal bundle to vary.
A deformation separates the moving curves. Counting their embedded
images in the cycle space and their contributions in every source
genus on the same family determines the difference of the universal classes.
\end{abstract}
\maketitle
\raggedbottom
\setcounter{tocdepth}{1}
\newpage
\tableofcontents
\newpage
\section{Introduction}

On a single smooth projective Calabi--Yau threefold, the scalar
local factors in the universal series of effective curves can occur
with exponents of arbitrarily large denominator. This already happens
in classes of divisibility three. To see what is being factored, assign
to each effective integral homology class the unit of its compact
space of cycles. Form the generating series of these units. A choice
of local model for the multiples of an embedded curve gives a factor
for each class and genus. Scaling the two normal directions with
opposite weights in a balanced bundle gives integral exponents
\cite{WueProduct,WueGenerators}. Scaling both directions with weight one gives
the scalar factors introduced by Pardon \cite{Pardon}.
The same spaces of cycles enter both products.

The compact example of \cite[Theorem~1.1]{WueCounter} concerns a
threefold containing a Hirzebruch surface. A primitive class on the surface has genus-three multiplicity
$-14$. At three times this class the genus-ten scalar exponent has
fractional part $2/3$. There are two
questions here. What denominator does the local class of a tripled
curve carry? Can a single compact target realize increasingly large
denominators once the covering degree has been fixed?

For the local question let $C$ be a smooth projective curve of genus
$g$ and $N$ a rank-two normal bundle of degree $2g-2$. The tripled zero
section is the effective cycle $3[C]$ in the total space of $N$.
An evaluation at a point of $C$ applies a fiber covector to the three
sheets of a cubic cycle and takes a homogeneous symmetric polynomial
in the three values. Its degree is its weight under scalar dilation.
A tuple with compact common zero locus gives an integral supported
Thom class. Dividing it by the product of the weights gives Pardon's
scalar class $x_{g,3}$ \cite[Definition~4.7]{Pardon}. Its denominator
is the least positive integer that takes it into $H_{\ZZ}$, the part
of the integral algebra of semi-Fano cycles in first Chern number and
virtual dimension zero. Here semi-Fano means that each irreducible
component has nonnegative first Chern number. We use the integral
generators of \cite[Theorem~1.1]{WueGenerators}, obtained by giving opposite
weights to two normal lines of degree $g-1$ whose product is $K_C$.

\begin{theorem}\label{thm:cubic}
In genera zero and one the scalar and opposite-weight classes coincide
in every degree. For $g\ge2$ the denominator of the scalar triple is
\[
 \operatorname{den}(x_{g,3})=
 \begin{cases}
  3^k, &g=2k+1,\\[2pt]
  3^{\max\{0,\,k-2-v_3(2k+1)\}}, &g=2k+2,
 \end{cases}
\]
where $v_3$ is the valuation at three. For a fixed curve of genus
$g\ge2$, the least product of weights, over all rank-two bundles of
degree $2g-2$ and all tuples of homogeneous single-fiber evaluations
with compact common zero locus on cubic cycles, is
$3^{\lfloor(g-1)/2\rfloor}$. A general stable bundle with determinant
$K_C$ attains the minimum.
\end{theorem}

For example, on a genus-four curve the best choice of normal bundle
allows a single evaluation of weight three. The resulting supported
class is divisible by three in $H_{\ZZ}$, so the normalized class is
integral. In genus five the minimum is nine and the class itself has
denominator nine. Computing a minimum weight product thus gives a
sharp geometric bound on the normalization, while computing the class
denominator requires the divisibility of the numerator as well.
For use below we write $s_{g-1}$ for the exponent of three in
$\operatorname{den}(x_{g,3})$ and put $s_0=0$.

Let $e_{\beta,h}(Y)$ denote the exponent of the scalar genus-$h$
factor at an effective class $\beta$. On our compact threefold we
find infinitely many of these individual exponents to which the
odd-genus formula applies exactly.

\begin{theorem}\label{thm:arith-compact}
There is a smooth projective Calabi--Yau hypersurface $Y$ containing
a first Hirzebruch surface with negative section $s$ and ruling $f$.
For each $a=12j+3$, $j\ge0$, put $\alpha_a=as+(a+1)f$ and
$k=(a(a-1)/2-1)/2$. The class $\alpha_a$ is primitive and
$e_{3\alpha_a,9k+1}(Y)$ has exact denominator $3^k$. The first of
these exponents, in genus ten, has fractional part $2/3$.
\end{theorem}

The local Gromov--Witten theory of Bryan and Pandharipande \cite{BP}
computes numerical residues for two normal weights. A numerical
residue alone does not identify a class in a universal cycle algebra.
To make the comparison, we vary a balanced sum of two normal lines
in a family of algebraic targets and impose conditions on the
characteristic polynomials of the sheets. At prime covering degree
the nonconstant coefficients in the ratio of the two weights can be
computed on embedded images, after deformation at each prescribed
bound on source genus. Intersection on the cycle space counts each
of these images with its sign. On the same family, stable maps with
that image include maps with contracted components, whose degree-one
series is $(2\sin(u/2))^{2h-2}$ for an image of genus $h$.
Theorem~\ref{thm:main} consequently identifies the difference of the
scalar and opposite-weight logarithms in prime degree. For triples,
the lower-degree terms are integral, so the cubic trace in
\S\ref{sec:arithmetic} determines the denominator of the full class.
A stable-pair calculation on the Hirzebruch surface singles out a
primitive genus with multiplicity prime to three. Its contribution
survives in the scalar exponent of Theorem~\ref{thm:arith-compact}.
The general formula in \S\ref{sec:compact} sums over all effective
integral roots of the triple class.

The genus-ten example distinguishes three measurements of the triple.
For a genus-three core the scalar-minus-opposite-weight cubic
logarithm has coefficient $4720/3$ at $(2\sin(u/2))^{18}$. Changing
variables and taking the coefficient of $u^{18}$ instead gives the
fixed source-genus-ten residue $23593/5$. It is the former coefficient
that multiplies $-14$ in the change of the compact exponents.
A different comparison uses the point class of an isolated triple in a
nonlinear neighborhood of the curve. For a general stable normal bundle of
canonical determinant there are eight one-parameter families of
embedded cyclic triples in genus three. A general quadratic term
obstructs all eight, and the isolated triple exceeds the scalar class
by $8x_{10,1}/3$, where $x_{10,1}$ is the class of a simple genus-ten
curve (Theorems~\ref{thm:classification} and~\ref{thm:isolated}).
This describes the moving covers behind the failure of triple
rigidity studied in \cite{BP06}.

In composite degree the characteristic polynomials can be proper
powers. We prove the comparison outside the weights where these
powers occur and compute the correction from doubled images. This
determines the scalar-to-opposite-weight conversion in degree four
for every core genus and in degree six for a genus-two core
(\S\ref{sec:composite}). Appendices~\ref{app:split}
and~\ref{app:pencils} treat individual split cubic weights and
the contribution of cuspidal cubics over a genus-five curve.

\subsection*{Acknowledgements}
I thank John Pardon for a helpful discussion of equivariant local curve classes.

\section{Changing the normal weights}\label{sec:statement}
Let $C$ be a smooth connected projective curve of genus $g\ge2$, and
let $L\oplus M$ be a balanced bundle, meaning
$LM=K_C$ and $\deg L=\deg M=g-1$. Write $H$ for the rational
Chern-number-zero, virtual-dimension-zero summand of Pardon's semi-Fano universal cycle algebra \cite[Definition~1.1 and Theorem~1.2]{Pardon} and
$x_{h,1}\in H$ for its multiplicity-one scalar generator in genus $h$.
For two nonzero integral fiber weights with ratio $z$, let
$X_{g,m}(z)$ denote the normalized class of the multiple zero section
of covering degree $m$. Section~\ref{sec:weights} constructs its
Laurent polynomial in $z$. We use the logarithm for disjoint union
and write
\[
 P_{m,j}=[z^j q^m]\log\Big(1+\sum_{l>0}X_{g,l}(z)q^l\Big).
\]
The core genus $g$ is fixed in this notation. Let $N_{k,m,j}$ be the
coefficient of $z^j$ in the connected genus-$k$, degree-$m$ local
Gromov--Witten residue of $L\oplus M$, with the two independent fiber
weights as in \cite{BP}. Set $n_s=(g-1)\binom{s}{2}$ and
$w=(2\sin(u/2))^2$.

\begin{theorem}\label{thm:main}
Let $d$ be the largest proper divisor of $m\ge2$. If
$n_d<|j|\le n_m$, then
\[
 P_{m,j}=\sum_{h\ge0}a_{h,j}x_{h,1},\qquad
 \sum_{h\ge0}a_{h,j}w^{h-1}
       =\sum_{k\ge0}N_{k,m,j}u^{2k-2}.
\]
The left sum has finite support. Its coefficients vanish for
$h<1+m(g-1)+|j|$. For each genus bound $K$, the construction in
Theorem~\ref{thm:separate} gives $a_{h,j}$ for $h\le K$ as the
signed number of constrained embedded curves of genus $h$ divided by
$\prod_{i=2}^m i^{2(i-1)(g-1)}$.
\end{theorem}

For a prime $p$ one has $d=1$ and $n_d=0$. For the degree-$p$ logarithm $P_p(z)=\sum_jP_{p,j}z^j$,
symmetry gives
\[
        P_p(1)-P_p(-1)=4\sum_{\substack{j>0\\j\ {\rm odd}}}P_{p,j}.
\]
This determines the connected scalar--antidiagonal discrepancy at $p$.
At composite $m$ the boundary of the stated
range has a geometric meaning. If $d>1$, a linear system of spectral
polynomials at $j=n_d$ necessarily meets the locus of powers of
degree-$d$ polynomials. Division of cycles then detects a nonzero
lower-degree primitive. We give the intersection calculation in
\S\ref{sec:weights}.

Let $X$ be a smooth projective complex threefold with
$K_X\simeq\mathcal O_X$ and $H^1(X,\mathcal O_X)=H^2(X,\mathcal O_X)=0$.
For an effective integral class $\beta$, let $U_{X,\beta}\in H$ be
the class of its compact cycle space. Write $F_h^{\rm sc}$ and
$F_h^{\rm ad}$ for the scalar and balanced antidiagonal universal
local series of a genus-$h$ core, both with constant term one.
The exponents $e_{\beta,h}$ and $c_{\beta,h}$ are defined by
\[
 \sum_\beta U_{X,\beta}Q^\beta
    =\prod_{\beta\ne0,h}F_h^{\rm sc}(Q^\beta)^{e_{\beta,h}}
    =\prod_{\beta\ne0,h}F_h^{\rm ad}(Q^\beta)^{c_{\beta,h}}
\]
in the effective integral class monoid completed by ample degree.
The $e_{\beta,h}$ are rational, and $c_{\beta,h}$ are the integral
elementary-cluster coefficients by \cite[Theorem~1.1]{WueProduct}.
If the image of $\alpha$ in the free
quotient of $H_2(X;\ZZ)$ is primitive, write
$D_{g,p}(w)$ for the difference of the scalar and antidiagonal connected
degree-$p$ local residues of a core of genus $g$. Theorem~\ref{thm:main}
gives the coefficient identity
\begin{equation}\label{eq:compact-prime}
 \sum_h(c_{p\alpha,h}-e_{p\alpha,h})w^{h-1}
    = \sum_{p\gamma=p\alpha}\sum_{g\ge2}c_{\gamma,g}D_{g,p}(w).
\end{equation}
Indeed the only covering degrees at $p\alpha$ are one and $p$, and at
every effective integral root $\gamma$ the degree-one generators of both
products agree. For core genera zero and one the two local classes
agree in all degrees. Finite support at each class makes~\eqref{eq:compact-prime}
a polynomial identity after taking the difference.

The first possible change is in genus $p+2$. The use of the rational
basis $x_{h,1}$ in this argument can also be made before assigning any
numerical invariant to a complex family: proper pullback to connected
maps is injective on the multiplicity-one span when all source genera
are used \cite[Sections 5--7]{WueProduct}. For integral denominators we
use the stronger result that the $a_{h,d}:=X_{h,d}(-1)$ freely generate
$H_{\ZZ}$, with $a_{h,1}=x_{h,1}$. This is the integral basis theorem
of \cite[Theorem~1.1]{WueGenerators}. In particular the intersection of
this integral algebra with the rational span of the $x_{h,1}$ is
their free integral span: they are distinct degree-one variables
in that polynomial presentation.

The proof of Theorem~\ref{thm:main} occupies
Sections~\ref{sec:weights}--\ref{sec:comparison}. The argument needs
a family of targets to give the same count in two different theories.
For a prescribed bound on source genus we vary the target and impose
the spectral coefficient conditions until the constrained connected
images are regular embedded curves. The cohomological insertion
counts their signs. The relative virtual class on maps also sees all
contracted components, and each embedded image of genus $h$ contributes
the entire degree-one series $w^{h-1}$. This construction uses an
algebraic relative obstruction theory on the chosen family and
proper genus detection on the cycle spaces.

\section{The cubic trace and its exact denominator}\label{sec:arithmetic}
For the numerical calculation let $F_g$ and $E_g$ be the disconnected
degree-generating series at scalar and opposite weights. Set $D_n=D_{n+1,3}$
and $n=g-1$. The transition to the integral lattice will use the full
row of coefficients of $D_n$, including coefficients in genera greater
than the genus of the base.
\subsection{The quadratic and cubic corrections}

Write $a_{g,m}=X_{g,m}(-1)$ and $x_{g,m}=X_{g,m}(1)$.
With $s=2\sin(u/2)$ the scalar and antidiagonal degree-two residues
are, respectively,
\[
 s^{4n}\big((4-2s)^n+(4+2s)^n\big),\qquad
 2s^{4n}(4-s^2)^n,\qquad n=g-1
\]
by \cite[\S8]{BP}. The linear terms of the two series agree. Thus
the difference $J_n=[q^2](\log F_{n+1}-\log E_{n+1})$ is
\begin{equation}\label{eq:arith-two}
 J_n(w)=2\sum_{l=1}^n4^{n-l}
 \left(\binom n{2l}-(-1)^l\binom nl\right)w^{2n+l}.
\end{equation}
In particular $J_1=2w^3$, $J_2=24w^5-2w^6$. The first formula and
Theorem~\ref{thm:main} give the universal identity
\[
 a_{2,2}=x_{2,2}-2x_{4,1}.
\]
At this genus the only nonconstant modes of $X_{2,2}(z)$ have
weights $\pm1$ and are equal. Each is $x_{4,1}/2$, as in the
geometric double-cover calculation of \cite{WueProduct}.

At degree three let $w^3\lambda_i$ be the three eigenvalues of the
balanced scalar handle operator. Its characteristic equation after
division by $w^3$ in the eigenvalue variable is
\begin{equation}\label{eq:arith-char}
 X^3-AX^2+\tfrac13 BX-\tfrac13 C=0,
\end{equation}
where, with $L_0=w^2+54w-243$,
\[
 A=w^2+50w+81,\quad B=4(w-6)L_0,\quad
 C=4(w-4)(w-3)^2L_0.
\]
This is the scalar specialization of the partition multiplication and
cap operators of \cite[\S6 and Appendix]{BP}, derived in Appendix~\ref{app:handle}. The antidiagonal eigenvalues are
$U,U,V$, where $U=w^3(4-w)(3-w)^2$ and $V=w^3(3-w)^2$ by the
hook-length formula. Writing $T_n=\sum_i\lambda_i^n$ gives
\begin{equation}\label{eq:arith-difference}
 D_n=w^{3n}T_n-2U^n-V^n-w^n J_n.
\end{equation}
Here $[q]F_{n+1}=[q]E_{n+1}=w^n$. The final term removes the
contribution of a degree-one and a degree-two connected component.
For example $T_1=A$ and $T_2=A^2-2B/3$. Substitution gives
\begin{align}
 D_1&=120w^4-20w^5+2w^6,\notag\\
 D_2&=14448w^7-1136w^8+\tfrac{4720}{3}w^9
             -332w^{10}+40w^{11}-2w^{12}.\label{eq:arith-row}
\end{align}
Let $\iota(w^{h-1})=x_{h,1}$, extended linearly. Recovering the
disconnected universal generators from their logarithms gives
\[
 a_{3,3}=x_{3,3}-\iota(D_2)
          +x_{3,1}(-24x_{6,1}+2x_{7,1}).
\]
Indeed the quadratic difference here is $-\iota(J_2)$, and the
linear generators agree. In particular the coefficient of $x_{10,1}$
in the scalar-minus-antidiagonal cubic logarithm is $4720/3$. The corresponding
coefficient in the disconnected scalar residue before the logarithm
and the change of weights is $[w^9]w^6T_2=292/3$.

The two positive odd weights contributing to $[w^9]D_2$ can also
be distinguished. Write $b_1,b_3$ for their respective coefficients
at $x_{10,1}$. All higher odd weights vanish at this genus by
Theorem~\ref{thm:main}. Let $B_d(z)$ denote the balanced handle at
weights $(z,1)$ in degree $d$, with the normalizations in
Appendix~\ref{app:handle}. The genus-three connected coefficient is
\[
 \operatorname{tr}(B_3(z)^2)-w^2\operatorname{tr}(B_2(z)^2)+w^6/3.
\]
For clarity, its values at the two additional weight ratios needed to
separate the modes are
\[
\begin{array}{c|cc}
 z&2&-2\\ \hline
 [w^9]&549&-19891/6 .
\end{array}
\]
They follow by substituting $(t_1,t_2)=(\pm2,1)$ in the multiplication,
metric and cap in Appendix~\ref{app:handle} and reducing with
$r^2-(2-w)r+1=0$, $r=e^{iu}$. Equality of the opposite weights and
subtraction at $z=1,-1$ and $z=2,-2$ give
\[
 4(b_1+b_3)=4720/3,\qquad 5b_1+\tfrac{65}{4}b_3=23185/6.
\]
Thus $b_1=674/3$ and $b_3=506/3$ in the universal algebra.

\subsection{The valuation of the cubic trace}

Write $v_3$ for the usual valuation on rational numbers, with
$v_3(3)=1$. The common denominator of a finite list of rational
numbers is its least positive integral multiplier making the entire
list integral.

\begin{theorem}\label{thm:arith-rows}
The common denominator of the coefficients of $D_n$ is $3^{s_n}$,
where
\[
 s_{2k}=k,\qquad s_{2k+1}=\max\{0,k-2-v_3(2k+1)\}.
\]
For $s_n>0$ the polynomial $3^{s_n}D_n$ has the following reduction
in $\mathbb F_3[w]$:
\begin{align*}
 n=2k:&\qquad 2(-1)^kw^{9k},\\
 n=2k+1:&\qquad
 2(-1)^k\frac{2k+1}{3^{v_3(2k+1)}}w^{9k+2}(w-1).
\end{align*}
\end{theorem}
\begin{proof}
Equation~\eqref{eq:arith-char} gives a recurrence over
$\ZZ[1/3][w]$ for $T_n$, with $T_0=3$, $T_1=A$ and
$T_2=A^2-2B/3$. In \eqref{eq:arith-difference} all terms other than
$w^{3n}T_n$ are integral polynomials. Work in the completion of
$\QQ(w)$ for the Gauss valuation at three, with residue field
$\mathbb F_3(w)$. Its valuation on a polynomial is the least valuation
of a coefficient. The polynomial
\[
 P(X)=3X^3-3AX^2+BX-C
\]
reduces to a linear polynomial with derivative $\bar B=w^3$.
Hensel's lemma gives a unit root $z_0$ with
$\bar z_0=\bar C/\bar B=\bar A$. For the two remaining roots let
$S=A-z_0$ be their sum and $T=C/(3z_0)$ their product. Since
\[
 P(A)=AB-C=72L_0(3w^2-14w-25),
\]
the divided difference $(P(A)-P(z_0))/(A-z_0)$ reduces to $\bar B$,
and hence
\[
 v_3(S)=2,\quad \overline{S/9}=2(w-1)/w,\qquad
 v_3(T)=-1,\quad \overline{3T}=w^3.
\]
The remaining roots have equal valuation $-1/2$ in a splitting
field, since their sum has positive valuation and their product
has valuation $-1$. Their ratio reduces to $-1$. For $k>0$ their
$2k$th power sum divided by $(-T)^k$ therefore reduces to $2$.
Adding the integral term $z_0^{2k}$ and multiplying by
$3^kw^{6k}$ proves the even case. Directly $D_0=0$.

For $n=2k+1$ the power sum of the remaining pair is
\[
 \sum_{l=0}^k\frac{n}{2l+1}\binom{k+l}{2l}(-T)^{k-l}S^{2l+1}.
\]
This is the coefficient of $t^n$ in
$-t\frac{d}{dt}\log(1-St+Tt^2)$. The $l=0$ term has valuation
$v_3(n)-k+2$. Each subsequent term has valuation greater by at least
$5l-v_3(2l+1)>0$. If the first valuation is nonnegative, the entire
trace is integral at three. Otherwise clearing its denominator leaves
only the first term in the reduction. Inserting the reductions of
$S/9$ and $3T$ and the factor $w^{3(2k+1)}$ gives the second formula.
\end{proof}

\begin{proposition}
Only primes at most $m$ divide $\operatorname{den}(x_{g,m})$.
For a prime $p$ and $g\ge2$ the $p$-part of
$\operatorname{den}(x_{g,p})$ is the $p$-part of the common
denominator of the coefficients of $D_{g,p}(w)$.
\end{proposition}

\begin{proof}
In the balanced split bundle use scalar weights $(1,1)$. The
elementary symmetric functions of the $d$ projected fiber coordinates on each
of the two normal lines are sections of their tensor powers of
orders at most $d$. Their evaluations at finitely many points
detect every coefficient, and simultaneous vanishing means that
the cycle is the $d$-fold zero section. The supported integral
Thom class is divided by a product of weights from $1$ through $d$
in Pardon's construction \cite[Definition~4.7]{Pardon}.
Hence $x_{g,d}$ is integral at every prime $p>d$.

In the difference of the degree-$p$ logarithms, the only term
with denominator $p$ from the logarithm uses $p$ copies of the common
linear coefficient. It cancels. Every term involving between two
and $p-1$ factors belongs to
$H_{(p)}=H_{\ZZ}\otimes\ZZ_{(p)}$: the denominator from the logarithm is a
unit at $p$ and the indices of its factors are smaller than $p$.
Also $a_{g,p}$ is integral. Theorem~\ref{thm:main} gives
$x_{g,p}-\iota(D_{g,p})\in H_{(p)}$. Because the $x_{h,1}$ belong
to the polynomial basis, clearing the denominator at $p$ on their
span is equivalent to clearing it on their coefficients.
\end{proof}

\begin{proof}[Proof of the denominator formula in Theorem~\ref{thm:cubic}]
For genera zero and one the local models of Section~\ref{sec:statement}
have empty separating tuple: the two lines of $\cO(-1)^{\oplus2}$
have no nonzero sections after any positive-degree cover, and in genus
one one can take a nontorsion degree-zero line and its inverse. A
section of a pulled-back degree-zero line would trivialize it and
taking the norm would make a positive power of the chosen line
trivial. Thus the whole compact cycle space in each multiple is the
point, independent of weights. For $g\ge2$ take $n=g-1$.
Theorem~\ref{thm:main} at degree two and equation~\eqref{eq:arith-two}
give $x_{g,2}-a_{g,2}=\iota(J_n)$ with $J_n\in\ZZ[w]$.
Thus every scalar double is integral. Expanding
the cubic logarithms gives
\[
 x_{g,3}=a_{g,3}+x_{g,1}\iota(J_n)+\iota(D_{g,3}).
\]
The first two terms on the right are integral, so the denominator
over $\ZZ$ is exactly the common coefficient denominator of
$D_{g,3}$. Its value is $3^{s_n}$ by the valuation in Theorem~\ref{thm:arith-rows}. In particular no factor of two
survives.
\end{proof}


Theorem~\ref{thm:arith-rows} also identifies the entire leading residue: in $3^{s_n}x_{n+1,3}\bmod3H_{\ZZ}$,
replace each $w^{h-1}$ in that theorem by $x_{h,1}$ whenever $s_n>0$.
For $g=2,3,4,5,6,7,8$ the resulting denominators are
$1,3,1,9,1,27,3$.

For arbitrary nonzero normal weights $(r,s)$ in the genus-two
quadratic case, the symmetry and the Laurent bound give
$X_{2,2}(r/s)=Y+(r/s+s/r)Z$. The difference at equal and opposite
weights is $2x_{4,1}$, whence
$X_{2,2}(r/s)=x_{2,2}+(r-s)^2x_{4,1}/(2rs)$.

For comparison, at a
genus-three core, the scalar-minus-opposite-weight connected logarithm
has coefficient $4720/3$ at $(2\sin(u/2))^{18}$. At fixed arithmetic
source genus ten its coefficient at $u^{18}$ is $23593/5$. Finally,
for a general stable normal bundle a neighborhood with a general second
jet isolates the triple and its point class exceeds the scalar class
by $(8/3)x_{10,1}$. Here eight counts cyclic covers and three is the
weight of their amplitude.

\section{Scalar exponents on a compact threefold}\label{sec:compact}

\begin{corollary}\label{cor:arith-roots}
Let the image of $\alpha$ in $H_2(X;\ZZ)/\mathrm{tors}$ be primitive,
where $X$ is a smooth projective Calabi--Yau threefold. For each $g$
let $C_g=\sum_{3\gamma=3\alpha}c_{\gamma,g}$, summing over the
effective integral roots. The common denominator of the scalar product
exponents $e_{3\alpha,h}$ is $3^M$, where
\[
 M=\max\bigl(\{0\}\cup
       \{s_n-v_3(C_{n+1}): n\ge0,\ C_{n+1}\ne0\}\bigr).
\]
In particular all these exponents are integral if and only if
$3^{s_n}\mid C_{n+1}$ for every $n\ge0$.
\end{corollary}
\begin{proof}
By \eqref{eq:compact-prime}, the polynomial of differences at this
class is $\sum_{n\ge0}C_{n+1}D_n$. There are finitely many occupied
genera, and $c_{3\alpha,h}$ is an integer for each $h$ by
\cite[Theorem~1.1]{WueProduct}. For $M>0$ scale by $3^M$ and reduce
modulo three. Each maximizing index contributes a nonzero unit times
the corresponding reduced polynomial of Theorem~\ref{thm:arith-rows}.
Their monomial supports, $\{9k\}$ or $\{9k+2,9k+3\}$, are disjoint
for distinct indices. This proves equality of the denominator; the
sum over torsion-distinct roots of a single genus has already been
taken in $C_g$. For $M=0$ all rows in the sum are integral.
\end{proof}

In particular \eqref{eq:arith-row} gives
\begin{equation}\label{eq:arith-ten}
 e_{3\alpha,10}=c_{3\alpha,10}-\tfrac{4720}{3}C_3.
\end{equation}
Indeed $D_1$ has degree six, and for $n\ge3$ the smallest possible
power of $w$ is $3n+1$ by Theorem~\ref{thm:main}.

Let $Y$ be the smooth anticanonical threefold in the toric resolution
of $X_9$ whose fan is specified in Appendix~\ref{sec:geometry}. Its second homology is
torsion-free and it contains a divisor $F\simeq\mathbb F_1$ with
normal line $K_F$. Let $s,f$ be the negative section and a ruling fiber.
For $a=12j+3$, $j\ge0$, set
\[
 \alpha_a=as+(a+1)f,\quad G=a(a-1)/2,\quad k=(G-1)/2.
\]
We use the same notation for curve classes on $F$ and their images in
$Y$. Two integral nef divisors on $Y$ recover the coefficients of
$s,f$ by intersection. Their difference takes value one on $\alpha_a$.
The four other nef degrees of \eqref{eq:degrees} vanish on both
$s$ and $f$. Their common zero face confines every effective curve of
class $\alpha_a$, and every effective summand, to $F$. Thus $\alpha_a$
has primitive free image and it is the sole integral root at
$3\alpha_a$.

These are the classes in Theorem~\ref{thm:arith-compact}.

\subsection{The highest primitive genus}

Put $D=as+(a+1)f$. In the presentation $F=\operatorname{Bl}_p\PP^2$
one has $D=(a+1)H-E$. Its arithmetic genus is $G$, and
\[
 N=h^0(F,\mathcal O_F(D))=\binom{a+3}{2}-1
                          =\tfrac12(a+1)(a+4).
\]
We show that the compact elementary exponents satisfy
\begin{equation}\label{eq:arith-minimum}
 c_{\alpha_a,h}=0\quad(h>G),\qquad
 c_{\alpha_a,G}=(-1)^{N-1}N.
\end{equation}

First consider stable pairs on $Z=\operatorname{Tot}(K_F)$. All
positive compact curves are supported on $F$: affine fibers contain
no complete curves, and on a normalization projecting nontrivially
to $F$ the tautological section of the negative canonical line
vanishes. For an effective decomposition $D=\sum D_i$, with at
least two nonzero classes $D_i=a_i s+b_i f$, $a_i,b_i\ge0$, we claim
\begin{equation}\label{eq:arith-split}
 \sum_i\chi(\mathcal O_{D_i})-\chi(\mathcal O_D)
 =\sum_{i<j}D_iD_j
 =\sum_i b_i(a-a_i)-\sum_{i<j}a_i a_j>0.
\end{equation}
These numbers depend only on the divisor classes by Riemann--Roch.
Choose $A=\max a_i$ and $t=a-A$. For $t>0$ the first sum in the
last expression is at least $(a+1)t$, and the second is at most
$At+t(t-1)/2$. Their difference is at least $t(t+3)/2$. For $t=0$
some positive multiple of a fiber lies in a different summand, giving
an intersection of at least $a$.

For a pure curve $W\subset Z$ filter $\mathcal O_W$ by powers of
the ideal of $F$. The graded quotients are of the form
$K_F^{-i}\otimes\mathcal O_F/J_i$. The divisorial hull of $J_i$ is
$\mathcal O_F(-D_i)$, and its quotient by $J_i$ has finite length
$l_i$. The sum of the divisor classes $[D_i]$ is $[W]$. The Euler
characteristic of the $i$th layer is
$\chi(\mathcal O_{D_i})-iK_FD_i+l_i$, with both corrections
nonnegative. By \eqref{eq:arith-split}, two nonzero divisor layers
in class $D$ give $\chi(\mathcal O_W)>1-G$. If there is only one,
it is the zeroth layer, since a nonzero higher layer along a curve
requires a nonzero zeroth layer along that curve. Then
$I_F\mathcal O_W$ has finite support and vanishes by purity. Such a
curve is an effective Cartier divisor in $|D|$. The section of a
stable pair has image $\mathcal O_W$ and finite-length cokernel.
It follows that the minimum pair Euler characteristic is $1-G$ and
all pairs at the minimum are quotient pairs in $|D|$. For a product
of pairs in a nontrivial additive splitting of $D$, apply
\eqref{eq:arith-split} to their combined divisor layers. Its total
Euler characteristic is strictly greater than $1-G$.

The minimum pair space is $\PP^{N-1}$ as a scheme. Indeed
surjectivity of the section in a family follows from Nakayama's
lemma, identifying it with the corresponding Hilbert functor.
For every $W\in|D|$, the normal sequence for $W\subset F\subset Z$
has terms $\mathcal O_W(D)$, $N_{W/Z}$ and $K_F|_W$. The last has
no sections since $H^0(F,K_F)=0$ and $H^1(F,K_F-D)=0$. The latter
follows from duality and Kodaira vanishing, as $D-K_F$ is ample.
Thus $h^0(N_{W/Z})=h^0(\mathcal O_W(D))=N-1$ by rationality of
$F$. The closed immersion of $|D|$ exhausts the closed points, so the
local dimension and the embedding dimension of the Hilbert space
both equal $N-1$ there. Its local rings are regular, proving the
scheme assertion. The symmetric pair obstruction theory on this smooth
proper scheme has degree $(-1)^{N-1}N$ \cite{PT}. The strict product
inequality implies that the coefficient of the formal logarithm of
the stable-pair series at $Q^D$ has the same lowest term,
$(-1)^{N-1}N z^{1-G}$ in the pair variable $z$.

To identify this term with a compact elementary exponent, deform $Y$
to the normal line of $F$. Blow up $Y\times\mathbb A^1$ at
$F\times0$ and remove the strict transform of $Y\times0$. The resulting
smooth family has fibers $Y$ off zero and $Z$ over zero and contains
a Cartier divisor $F\times\mathbb A^1$. In each positive face class
all connected maps have image in that divisor by the nef degrees on
$Y$ and negativity on $Z$. This is true as a statement about families
of maps: on every geometric source the Cartier line restricts to the
pullback of $K_F$. Sections vanish on every positive component by
negativity, then on all contracted components by connectedness. In a
two-term presentation of the derived direct image the first differential
is therefore fiberwise injective and locally split injective. The
canonical section vanishes on the universal domain, including over
nonreduced bases. Hence the relative moduli is the product of the
proper map stack of $F$ and the parameter line. Deformation invariance
of its ordinary relative stable-map class gives equality of the
connected GW degrees of $Y$ and $Z$ in each face class.

The fixed-target toric correspondence \cite[Theorem~1]{PP} applies
to $Z$. The pair spaces in these classes are proper: in the projective
line completion of $K_F$ their geometric supports are on the zero
section, and a flat specialization of a pair supported there remains
set-theoretically on that section. For instance a power of its ideal
annihilating the generic pair sheaf annihilates the flat family over
a discrete valuation ring. Thus the equivariant degrees have their
ordinary specialization. Since $K_Z$ is trivial there is no prefactor
in the empty-insertion correspondence under $z=-e^{iu}$. Taking the
formal logarithm through class $D$ uses only the finitely many
positive summands on this face, with disconnected maps nonconstant on
each connected component. The result identifies the connected fixed-target
GW series for $Y$ in $\alpha_a$ with the rational function of $z$
whose lowest term was computed above.

On the compact projective $Y$, the algebraic and symplectic GW degrees
agree \cite{Siebert}. At a class with primitive free image the finite
elementary-cluster expansion of \cite{IP,DIW} is
$\sum_h c_{\alpha_a,h}w^{h-1}$, with these same compact exponents by
\cite[Theorem~1.1]{WueProduct}. Since $w=(1+z)^2/z$, this is a rational
function of $z$. Equality of the expansions at $z=-1$ gives equality
of rational functions. At $z=0$ the summand of largest genus $h$ has
lowest Laurent power $1-h$ with coefficient $c_{\alpha_a,h}$.
Comparison with the lowest pair term proves~\eqref{eq:arith-minimum}.

\begin{proof}[Proof of Theorem~\ref{thm:arith-compact}]
For $a=12j+3$ the integer $G=2k+1$ is odd and $N\equiv2\pmod3$.
In \eqref{eq:compact-prime} the top occupied row is $D_{2k}$.
Theorem~\ref{thm:arith-rows} shows that its denominator is $3^k$ and
that $3^kD_{2k}$ reduces to $2(-1)^kw^{9k}$. Every lower row has
denominator dividing $3^{k-1}$. The coefficient multiplying the top
row is a unit at three by~\eqref{eq:arith-minimum}, while the
compact elementary exponents are integral. This proves the exact
denominator at output genus $9k+1$. For $a=3$ one has $G=3$, $N=14$
and $c_{\alpha_3,3}=-14$. Equation~\eqref{eq:arith-ten} gives
$e_{3\alpha_3,10}=c_{3\alpha_3,10}+66080/3$. For $a=15$ the
corresponding genus is $469$ and the denominator is $3^{52}$.
\end{proof}

\subsection{The first possible fractional exponents}
If the free image of a class is primitive its scalar and opposite-weight
exponents agree. At free divisibility two, \eqref{eq:compact-prime}
and the integral coefficients of \eqref{eq:arith-two} imply
integrality in all genera, summing over every integral double root.
There is also a uniform statement in small genus. For all effective
$\beta$ on a smooth projective Calabi--Yau threefold as above,
\[
 e_{\beta,h}=c_{\beta,h}\quad(0\le h\le3),\qquad
 e_{\beta,4}=c_{\beta,4}-2\sum_{2\gamma=\beta}c_{\gamma,2}.
\]
Indeed pull the logarithms of the two compact products back to
connected maps in source genus at most four. A positive-degree
map to a core of genus $g\ge2$ in projection degree $d$ has source
genus at least $1+d(g-1)$ by normalization and Riemann--Hurwitz.
The zero and one core genera agree in every degree. At the remaining
small degrees Theorem~\ref{thm:main} and \eqref{eq:arith-two} leave
precisely the term $2x_{4,1}$ from $(g,d)=(2,2)$ through source genus
four. This assertion also holds for the equivariant cycle insertions:
their supported pullbacks are over the zero-section multiple, hence
empty below that same genus. Induct on ample degree, reading source
genera in increasing order by the proper detection of
\cite[Sections~5--6]{WueProduct}. Proper multiples of a genus-four change
can first be detected in genus seven. This gives the displayed identities.

One may also try to deduce integrality from the scalar stable-pair
series. If all scalar local pair coefficients were integral Laurent
series in the pair variable $z$, division of a compact integral pair
series by already determined lower factors would leave an integral
coefficient $\sum_g e_{\beta,g}z^{1-g}(1+z)^{2g-2}$ in the next
class. The largest occupied genus has a unique lowest Laurent term
with unit coefficient; descending induction would give integral
exponents. The local hypothesis fails. In the genus-three scalar
cubic the disconnected trace in \eqref{eq:scalar-traces}, under
$w=(1+z)^2/z$, has coefficient $20+292/3=352/3$ at $z^{-9}$.
Here one uses the fixed-target equivariant local GW/pairs correspondence
\cite[Section~4.2]{Pardon}; the lowest compact proper pair count used
in Theorem~\ref{thm:arith-compact} is integral.

\section{Fiber weights and spectral polynomials}\label{sec:weights}

\subsection{The class over a product of projective spaces}
Pardon's algebra $H$ is obtained from compactly supported cohomology on
spaces of cycles in families of threefolds. A closed oriented real
parameter space of dimension $b$, together with a compactly supported
$b$-cocycle on its relative cycle space, gives a class of relative
degree zero. A cocycle over an interval of such families equates its
two restrictions if its support stays compact. One may also restrict
the target to an open neighborhood containing that support.
Multiplication uses the disjoint union of targets. We use the complex
families of \cite[\S3]{Pardon} and their map $\Phi$ to the corresponding
algebra of almost-complex families \cite[\S2]{WueProduct}. The latter
allows smooth oriented parameter chains.

Take the two lines general in their Picard varieties subject to $LM=K_C$.
Thus $H^0(C,L)=H^0(C,M)=0$ and $L^2\ne K_C$.
For $A=L,M$, the vector space $H^0(C,A^i)$ has dimension
$(i-1)(g-1)$ for $i\ge2$, by Riemann--Roch and Serre duality.
A compact curve in the rank-two bundle projects finitely onto $C$,
since a fiber is affine. With multiplicities included, a cycle of
covering degree $m$ gives on each axis a monic polynomial of degree
$m$. Its coefficient of order $i$ is a holomorphic section of $A^i$.
In families these coefficients are continuous functions of the image
cycle. Near a transverse base fiber the local sheets give elementary
symmetric functions. Their extension at a branch value is the
characteristic polynomial of multiplication in the finite divisor on
the line axis. Equivalently one can push the cycle to the line and use
the coordinates on its space of monic spectral divisors
\cite[\S2]{WueBalanced}.

Fix a basis of the dual of each coefficient space using linear
combinations of evaluations at finitely many points of $C$.
Write $F\subset C$ for their union. These are $n_m$ tests on each
axis. They all vanish precisely on $m[C]$.
Let $P,Q$ denote the hyperplane lines on the factors of
$B_i=\PP^i\times\PP^{2n_m-i}$. Twist $L,M$ by $P,Q$ respectively.
Each test now takes values in $P^s$ or $Q^s$, where $s$ is its
order. Put
\[
 W_{g,m}=\prod_{s=2}^m s^{2(s-1)(g-1)}.
\]
On the relative cycle space take the pullback of the Thom class of
the direct sum of all test lines and divide by $W_{g,m}$. We denote
its image in $H$ by $Y_i$. Its cohomological degree equals
$\dim_{\RR}B_i$, and its support is the compact family of core cycles.
The choice of a general splitting computes the same normalized
one-circle classes as any splitting of the specified degrees.
Along a path in the Picard variety use a fixed sufficiently large
number of fiber evaluations and frames on small intervals. They
separate all coefficients throughout the path. Appending a redundant
test multiplies the insertion by its normalized Euler class and is
exactly the hyperplane stabilization of the Borel construction
\cite[\S2]{WueBalanced}.

\begin{lemma}\label{lem:weights}
For nonzero integral fiber weights $r,s$ the corresponding class is
\[
               X_{g,m}(r/s)=\sum_{i=0}^{2n_m}(r/s)^{i-n_m}Y_i .
\]
It is invariant under inversion of the weight ratio. The logarithmic
coefficient $P_{m,j}$ is primitive for the coproduct induced by addition
of cycles. Pullback $\rho_e$ under multiplication of cycles by $e$
sends $P_{m,j}$ to $P_{m/e,j}$ if $e\mid m$, and to zero otherwise.
\end{lemma}
\begin{proof}
Over two large projective spaces form the same supported tuple.
Pull it back by the pair of degree-$|r|,|s|$ Veronese maps from
$\PP^{2n_m}$, conjugating the corresponding source when a weight is
negative. On that oriented real-analytic base the resulting lines are
the lines of Chern classes $rH,sH$, where $H$ is the hyperplane
class. The extra scalar needed for the one-circle normalization is
$r^{-n_m}s^{-n_m}$. The fundamental class of the graph of the pair
maps to the sum of the coordinate linear products
$[\PP^i\times\PP^{2n_m-i}]$ with coefficients $r^is^{2n_m-i}$,
as is seen by pairing with the two hyperplane classes. Pairing a base
chain with the supported test cocycle commutes with its boundary:
the cocycle restrictions on common faces agree in the shifted
cochain complex for the dualizing sheaf of the base
\cite[\S2.4 and Remark~3.6]{Pardon}. These chain relations can be
taken real analytic on each simplex. Indeed embed the product of
projective spaces as Hermitian projectors, approximate a subdivided
bounding chain by affine interpolation of vertices followed by the
analytic tubular retraction, and use the straight interpolation and
retraction on the boundary prism. The support over each simplex is
the compact core family. This proves the polynomial formula. Swapping
the lines and moving the splitting along the Picard variety gives
its symmetry. For each pair of weights the cycle addition and
multiplication identities for the normalized tests are those of
\cite[\S3]{WueBalanced}; their proofs use only the orders and the
nonvanishing of the weights. Taking the logarithm and comparing
coefficients as Laurent polynomials gives the asserted operations.
\end{proof}

The degree in $z$ of the degree-$l$ logarithm is at most $n_l$ by
superadditivity of $n_l$. If $|j|>n_d$ for the largest proper divisor
$d$ of $m$, the primitive $P_{m,j}$ is therefore killed by every
proper division. By \cite[Lemma~7.1]{WueProduct} it is a finite linear
combination of $x_{h,1}$. This use of the primitive-division criterion
is independent of any evaluation by numerical invariants.

We recall the geometric detection that will determine the coefficients.
On the space of connected stable maps of a fixed arithmetic genus $k$
to fibers of a smooth family, the map to the image cycle is proper
over compactly supported insertions \cite[\S5]{WueProduct}. Pullback of
the insertion, with the oriented parameter chain, gives a class
$T_k$ in the corresponding colimit of map-space cohomology groups.
The construction annihilates products of positive-degree classes
represented on disjoint targets. In an acyclic balanced normal bundle
the only degree-one image is its zero section, so
$T_k(x_{h,1})=0$ if $k<h$. At $k=h$ the only map to this cycle is
the regular embedding, and $T_h(x_{h,1})\ne0$ by the fixed-genus
detection theorem of \cite[\S\S5--6]{WueProduct}. We shall use only
these properties of $T_k$. They require no numerical functional on a
family of singular maps.

\subsection{The algebraic pairing}
Let $\mathcal F=\overline{\mathcal M}_k(C,m)$ be the proper stack of
connected unmarked stable maps, with universal source $\pi$ and map
$f$. Write $V_A=R\pi_*f^*A$ for $A=L,M$ and
$a=k-1-m(g-1)$. Riemann--Hurwitz on the positive components gives
$a\ge0$ if $\mathcal F$ is nonempty. Both complexes have virtual
rank $-a$, and $\mathcal F$ has virtual complex dimension $2a$.
Expanding the two inverse moving Euler classes in the fiber
characters gives
\begin{equation}\label{eq:chern}
 N_{k,m,j}=
  \int_{[\mathcal F]^{\rm vir}}
    c_{a-j}(-V_L)c_{a+j}(-V_M),\qquad |j|\le a ,
\end{equation}
and $N_{k,m,j}=0$ outside this range
\cite[\S2]{BP}. The ranks refer to Euler characteristics and are
unchanged by jumps of the spaces of sections.

The pairing of the normalized tuple over $B_{n_m+j}$ with the
fiberwise algebraic stable-map virtual class is precisely
$N_{k,m,j}$. To verify it, first use the compactification
$\PP_C(\mathcal O\oplus L\oplus M)$ over a point. The tuple with
two independent torus characters $x,y$ is supported on maps into
the core and extends by zero to the projective stable-map space.
The only fixed component meeting this support is $\mathcal F$:
invariant finite multisets in the affine fiber lie at its origin.
The fixed obstruction theory there is the one for maps to $C$. The
moving theory is $V_L\otimes x\oplus V_M\otimes y$, where the
characters also denote their weight lines. Virtual localization
\cite{GP} identifies the equivariant pushforward with
\[
       x^{n_m} y^{n_m}\sum_{l}N_{k,m,l}(x/y)^l .
\]
The monomial comes from the normalized Euler class of the tests on the
fixed locus. The pushforward is an ordinary equivariant polynomial
because the support is proper. In the associated projective bundle
over $B_{n_m+j}$ the fiberwise obstruction theory is the descent of
the fixed-target theory. Keeping the residual two-torus action thus
substitutes $x+c_1(P),y+c_1(Q)$ in this polynomial. Integrating over
the base at $x=y=0$ gives the claimed coefficient, including when
$a>n_m$. This calculation uses the stable-map theory for the family
of threefolds over $B_{n_m+j}$. Maps to the total projective variety
considered as a single target have a different obstruction theory.

\subsection{Avoidance of multiples}
Assume $j>0$ by symmetry. Keep the $L$ tests equal to zero and
translate the $M$ tests to homogeneous polynomials of their respective
orders on $\PP^{n_m-j}$. Along the linear homotopy the image cycle is
confined to the $M$ axis. The prescribed coefficients determine one
relative monic divisor in each fiber, with multiplicities. Its graph
over the parameter base and the interval is compact.

Let $e>1$ divide $m$ and put $s=m/e$. The $e$th powers of monic
spectral polynomials of degree $s$ form a closed subvariety of the
affine coefficient space of codimension $n_m-n_s$. Indeed the first
$s$ coefficients of a power successively recover those of its monic
root by division by $e$. The remaining coefficients are polynomials
in the recovered coefficients, of the required orders under scaling.
This argument works also for polynomials whose divisors are reducible.
A section whose values are general homogeneous polynomials avoids the
union of the power loci if $j>n_d$: evaluation is surjective at every
point of the projective base and its dimension is less than every
such codimension. A neighborhood of the compact translated insertion
then misses all cycle multiples, since a convergent sequence of
multiples has a convergent subsequence of roots by compactness of
bounded cycles. In particular this is possible for every nonzero Laurent exponent in the stated range
when $m$ is prime. A divisor in the permitted neighborhood can still
have individual repeated components.

At the first excluded positive weight for composite $m$, take
$s=d$, $e=m/d$ and $j=n_d$. The equations on the remaining
coefficients, of orders $i>d$, are sections of $\mathcal O(i)$ with
$(i-1)(g-1)$ copies for each $i$. Their top Chern number on
$\PP^{n_m-n_d}$ is
\[
                \prod_{i=d+1}^m i^{(i-1)(g-1)} >0 .
\]
Thus the entire power locus cannot be avoided by choosing other
homogeneous values. In fact $\rho_eP_{m,n_d}=P_{d,n_d}\ne0$:
the latter has nonzero connected minimum-genus detection at its
outermost weight, as we verify next.

\begin{lemma}\label{lem:outer}
For $s\ge2$ the coefficients at the two outermost weights of
$P_s(z)=\sum_jP_{s,j}z^j$ are both
\[
 \frac{x_{H_s,1}}{A_s},\qquad H_s=1+(g-1)\binom{s+1}{2},
 \qquad A_s=\prod_{i=2}^s i^{(i-1)(g-1)} .
\]
\end{lemma}
\begin{proof}
At $j=n_s$ the second projective factor is a point.
Translate its $M$-tuple to a smooth integral spectral divisor $D$,
for example $v^s=b$ with $b$ having a simple zero and all zeros
simple. For general $M$ such a $b$ exists: the Abel map for degree
$s(g-1)\ge g$ is dominant and the reduced divisors are dense.
Eisenstein's criterion gives integrality; adjunction gives genus
$H_s$. The other tests still vanish, so the unique supported
cycle is $D$ with zero $L$ field. For the projection $f:D\to C$ the
monic equation gives $f_*\mathcal O_D=\bigoplus_{i=0}^{s-1}M^{-i}$.
Consequently $R=f^*L$ has $H^0(D,R)=0$. In the surface the normal
line $J=f^*M^s$ satisfies $JR=K_D$, so $H^1(J)=0$. Coefficient
variation gives all $H^0(J)$ of dimension $n_s$. Their tests cut
these surface deformations with degree one.
On the smooth spectral locus these are all the curves: the $M$ projection is
birational to a smooth divisor and the lift is a
section of $R$, which vanishes. This describes its topology as well,
by the family of monic divisors and the continuous inverse given by
their coefficients. The normalized $L$ tests there give the ordinary
Euler class $c_1(\mathcal O(1))^{n_s}$. Base Gysin replaces
$\PP^{2n_s}$ by a linear $\PP^{n_s}$. What remains is the point class
of the fixed smooth divisor in the coefficient space, divided by $A_s$.

We identify that point class in the universal algebra. Choose a basis
$c_1,\ldots,c_{n_s}$ of $H^1(D,R)$ and a class
$\delta\in H^1(D,RJ^{-1})$ inducing an isomorphism
$H^0(J)\to H^1(R)$ by multiplication. Indeed $H^1(RJ^{-1})$ is
dual to $H^0(J^2)$. Choose $n_s$ points whose evaluations in fixed
frames form a basis of $H^0(J)^*$, and take the sum of evaluations in
the squared frames. The pairing induced by this functional is
nondegenerate. On a cover of $C$ by two affine opens represent the
$c_i$ and $\delta$ by cocycles on their preimages in $D$. The monic
basis lifts them to polynomials in $v$ on the overlap with the
appropriate line twists. Write $P_*(v)$ for the fixed monic equation.
In the $L$ line change the transition by adding the translation
\[
  t\Big(\sum_{i=1}^{n_s}\sigma_i\widetilde c_i(v)
     +\sigma_0 P_*(v)\widetilde\delta(v)\Big),
       \qquad 0\le t\le1 ,
\]
where the $\sigma_i$ are the homogeneous coordinates on
$\PP^{n_s}$. The point test is still the equation of the fixed $M$ spectral divisor. On $D$ the translation splits for every parameter at $t=0$,
and for $t>0$ precisely at $[1:0:\cdots:0]$, with unique lift.
Thus the point test has compact support throughout this deformation.
Near the final point the lifting obstruction is a holomorphic section
of $H^1(R)\otimes\mathcal O(1)$, using reduction modulo the variable
monic polynomial to identify the cohomology in families.
Its derivative in the projective coordinates is the basis $(c_i)$
and its derivative in the divisor coordinates is multiplication by
$\delta$, up to a common sign. Both are isomorphisms. Its zero
scheme is a graph over either set of coordinates. The lifts vary
holomorphically on this graph by cohomology and base change and
$H^0(R)=0$. The divisor point test on the graph is therefore a
positive parameter point class. In that fixed target the normal
sequence of the lifted curve has subline $R$, quotient $J$, and
connecting isomorphism on $H^0(J)$ just computed. Its normal bundle
is acyclic. Deformation to the normal bundle identifies the point
class with $x_{H_s,1}$. Finally every product in the logarithmic
subtraction has strictly smaller weight than $n_s$. This proves the
formula, including its normalization.
\end{proof}

\section{The common algebraic family for the residue comparison}\label{sec:deformations}

Let $C$ be a smooth projective curve of genus at least two and $B$ a
smooth projective variety. We start with
\[
 X_0=\operatorname{Tot}\big((L\boxtimes P)\oplus(M\boxtimes Q)\big)
       \longrightarrow C\times B,
 \qquad L\otimes M=K_C ,
\]
where $P,Q$ are line bundles on $B$. Fix a degree $m>0$ over $C$.
By the support of a cycle we mean the union of its distinct image
curves. If $A=\sum_i e_i A_i$ is a cycle, its reduced cycle is
$\sum_i A_i$. A \emph{proper multiple} is a cycle $eA'$ with
$e>1$ and $A'$ effective. Necessarily $e\mid m$.

Let $\tau$ be a tuple of the coefficient evaluations of
\S\ref{sec:weights}, translated by a section of its test bundle
$T\to B$. For this section it suffices that the evaluations be
linear combinations of symmetric functions of the intersection with
finitely many framed fibers over $F\subset C$, that the translation
be smooth, and that $\operatorname{rk}_{\mathbb C}T=\dim_{\mathbb C}B$.
We allow the supported Thom class to be localized in an open set of the
degree-$m$ cycle space, provided its zero set there is compact. All
motions below refer to that localized class. They preserve any scalar
factor in its normalization.

\begin{theorem}\label{thm:separate}
Suppose the compact zero set of $\tau$ contains no proper multiple.
For every $K$ there is a smooth quasi-projective variety $S$ smooth
over $B$, an algebraic section $i:B\to S$, and a smooth family of
Zariski locally trivial fibrations $X\to C\times S$ with fiber $\mathbb A^2$ restricting along $i$ to $X_0$.
There is a homotopy of the parameter section and the tests, with compact
joint zero set in the chosen cycle neighborhoods, with the following
properties.
\begin{enumerate}
\item At its endpoint, the zeros with connected support and sum of
normalization genera at most $K$ are finitely many smooth embedded
curves of multiplicity one. The relative space of embedded curves is
smooth of dimension $\dim S$ there. The target variations surject onto the normal obstruction space
there. The relative space meets the parameter graph transversely, and
the tests have transverse zero on this intersection.
\item On the single algebraic family $X/S$, the products of the test
Thom class with the normal Thom classes of the parameter graphs have
equal compactly supported classes at the two ends. At the initial end,
capping the relative stable-map virtual class with the graph class
restricts it to the relative stable-map class of $X_0/B$.
\end{enumerate}
In particular every connected stable map of arithmetic genus at most
$K$ on the endpoint zero set has one of the indicated image cycles.
\end{theorem}

The tests in this statement may have zeros of disconnected support.
The compact support in the second assertion includes such zeros. We
impose connectedness only in determining the image of a connected
stable map. We first describe the target deformations used in the
proof.

\subsection{Deformations at finitely many fibers}

Consider a family already constructed over a smooth quasi-projective
$S\to B$. Off a finite subset $P_0\subset C$, disjoint from $F$,
keep its identification with the pullback of $X_0$. Initially $S=B$
and $P_0$ is empty. Fix $p\notin F\cup P_0$. On a Zariski neighborhood
$U$ of $p$ choose frames for $L,M$ and a uniformizer $t$ regular
there with its only zero at $p$. Shrink $U$ to avoid the old punctures.
In these frames the fiber variable $y$ takes values in the pullback
of the plane bundle $E=P\oplus Q$ over $B$.

Let $b,r$ be positive integers. On $U\setminus\{p\}$ use the change
of coordinates
\begin{equation}\label{eq:shear-gluing}
 y\longmapsto y+(\iota_\ell\Omega)
           \sum_{a=1}^b\sum_{v=0}^{r-1} c_{av}t^{-a}\ell(y)^v,
 \qquad
 \ell\in E^\vee,\quad \Omega\in\bigwedge^2 E,\quad c_{av}\in\mathbb C .
\end{equation}
All the displayed parameters are coordinates on a vector bundle over
$S$. Since $\ell(\iota_\ell\Omega)=0$, this automorphism preserves
$\ell(y)$, has inverse obtained by negating the $c_{av}$, and has
Jacobian determinant one in the fiber. At $p$ we glue a new copy of
the original affine chart by a composition of two such automorphisms
with independent parameters. For a finite set of new punctures use
the old target on their complement. Each new chart is chosen to miss
every other puncture. On the intersection of two new charts use their
transitions through the old chart. This gives an algebraic family
$X'\to C\times S'$ with $S'$ the total space of the new parameter
bundle. It restricts to the old family whenever all $c_{av}=0$,
including at the algebraic zero section $S\to S'$. Off the union of
the old and new punctures the original identification is preserved.
Since the inserted determinants are one, the relative tangent
determinant over $C$ in each threefold is still pulled back from
$L\otimes M$. Every compact curve in it has threefold Chern number
zero. The affine-plane fibers contain no compact curves.

Here is the infinitesimal effect. Let
$h_i:D_i\to X_s$ be the smooth normalizations of the distinct
positive components of a cycle, with finite projections
$f_i:D_i\to C$. Set $D=\coprod_iD_i$, $f=\coprod_i f_i$, and let
$V|_{D_i}=h_i^*T_{X_s/C}$. Over a general $p$ the covers are
unramified and all $r=\sum_i\deg f_i$ sheet values in the affine
plane are distinct. Indeed the normalized maps are birational onto
different curves. The number $r$ counts the sheets of the reduced
cycle. Choose two independent forms $\ell$ which each separate their
values at $p$, and use a nonzero $\Omega$. These choices exist since
only finitely many lines in $E^\vee$ are forbidden.

\begin{lemma}\label{lem:shear-jets}
Let $Z$ be a finite subscheme of the smooth domains, disjoint from
$f^{-1}(p)$. For sufficiently large $b$, variations of the
coefficients $c_{av}$ in the two factors of \eqref{eq:shear-gluing},
at $c=0$, span the obstruction space $H^1(D,V(-Z))$ for deforming
the graphs with fixed covers and prescribed restriction to $Z$.
\end{lemma}

\begin{proof}
Since $f^*p$ has positive degree on each component, we can arrange
\[
 H^1(D,V(bf^*p)(-Z))=0 .
\]
The connecting morphism from
principal parts of order $b$ over $p$ then surjects onto
$H^1(D,V(-Z))$. On each sheet, multiplication by $t^b$ identifies
these principal parts with two copies of
$R_b=\mathbb C[t]/(t^b)$. Write $y_1(t),\ldots,y_r(t)$ for the
formal sheet values. For either $\ell$ the matrix
$(\ell(y_j(t))^v)_{1\le j\le r,\ 0\le v<r}$ has Vandermonde
determinant a unit in $R_b$. For each $v$ the coefficients
$\sum_{a=1}^b c_{av}t^{b-a}$ run through $R_b$. Thus one shear
prescribes arbitrary principal parts on every sheet in the direction
$\iota_\ell\Omega$. The two directions span the plane. Differentiation
of the transition on the punctured overlap gives exactly the resulting
connecting class, up to the choice of sign for the \v Cech
differential. Prescribing the graph on $Z$ replaces $V$ by $V(-Z)$
because the difference of two local lifts vanishes on $Z$.
\end{proof}

We will need this lemma on a compact set of smooth normalization
tuples, simultaneously for $Z$ of lengths zero, one and two.
The corresponding relative symmetric powers of the disjoint smooth
source curves are proper over the tuple space. For each tuple
one can take three good punctures, so that at least one avoids the
base images of any such $Z$. At a given puncture $b$ can be chosen
uniformly over these subschemes by relative Serre vanishing for the
ideal of the universal divisor. Surjectivity of the graph deformation
problem with the jets imposed is open by the rank condition on a two-term
resolution of the sections of $V(-Z)$. Its differential augmented by
the parameter variations is surjective. This remains true where the dimensions of
its kernel or of $H^1(D,V(-Z))$ change. Compactness gives finitely
many choices of directions on neighborhoods of the tuples and their
marked subschemes. Distinct choices of $\ell,\Omega$ at the same
Borel parameter can be placed in separate factors of the parameter
bundle. Extend each choice locally over $B$ and then smoothly by a
cutoff. With the $c_{av}$ still zero, moving the parameter section to
these reference values leaves the target family unchanged. Coincident
choices of new punctures can also be accommodated by composing their
shears in the same transition. The linearized variations add at $c=0$.

\subsection{Specialization of the weighted images}

We give the compactness argument for the tuples to which the lemma
will be applied. Each $X'\to C\times S'$ is affine of finite type by
Zariski descent, and $C\times S'$ is quasi-projective. Hence $X'$ is
quasi-projective, and admits a projective closure over $S'$. In each
threefold fiber the projection to $C$ has contractible fibers. Thus
on degree-$m$ curves any chosen relative polarization has fixed
degree locally in the parameter. On compact sets of parameters these
degrees are bounded. The degree of a positive summand is bounded by
that of the whole cycle for the polarization.

Choose a relatively compact open neighborhood of the original compact
zero set, with closure in the given open locus of cycles. Such
neighborhoods can be taken in the fixed-degree relative Chow space
of a projective closure, on the open set of cycles with support in
the affine target. They can also be obtained from a compact bound
on supports and masses \cite[Theorem~2.11]{DIW}. For $e\mid m$ the
image of multiplication by $e$ is closed in this locus. Along a
convergent sequence of multiples the roots have the same supports
and their masses are divided by $e$. A subsequential limit of the
roots is therefore a root of the limiting cycle. Choose the closure
disjoint from these images for all $e>1$. Its boundary misses the
zero set. Whenever the parameter family is enlarged, take such a
neighborhood of the whole current compact zero set in the enlarged
relative cycle space, with closure on the old graph inside the
preceding neighborhood. The compact boundary misses the simultaneous
zero of the old graph and tests. It continues to miss this zero for
all sufficiently small motions of both. The supported Thom product
in the new neighborhood thus restricts to the previous product by
excision. Applied successively, this construction keeps compact
support also for zeros whose normalizations exceed the genus cutoff.

For a cycle $A=\sum_{i=1}^r e_iA_i$ write $d_i$ for the projection
degree and $h_i$ for the normalization genus of $A_i$. Set
\[
       \lambda(A)=\left(\sum_i d_i,\ \sum_i h_i,\ -r\right)
\]
with the lexicographic order. The first entry uses the reduced cycle.
Only finitely many values occur for $\sum_i h_i\le K$.

\begin{lemma}\label{lem:weighted-specialization}
For a convergent sequence in the chosen neighborhood with fixed
$\lambda$, the limiting value of $\lambda$ is at most that of the
sequence, and the sum of normalization genera cannot increase. If
$\lambda$ stays equal, then, after relabeling and taking a
subsequence, the normalization maps converge to the smooth
normalizations of the limiting cycle, with their original weights.
Connected supports have connected limit.
\end{lemma}

\begin{proof}
Keep the degrees, genera and weights of the individual labels constant
along a subsequence. Stable-map compactness applied to their
unweighted normalization maps gives a stable limit of each label.
One can take these limits in the projective closure. The bounds on the
component degrees give properness there, and the limiting total support
lies in the affine open. The weighted sum of their image cycles is
the given limit. Write $A'_j$ for its distinct image curves and
$b_{ij}$ for the sum of the covering degrees onto $A'_j$ in the
$i$th stable limit. With $d'_j=\deg(A'_j/C)$ one has
\[
       d_i=\sum_j b_{ij}d'_j,\qquad e'_j=\sum_i e_i b_{ij}.
\]
Every column sum of $(b_{ij})$ is positive. It follows that the
reduced degree drops unless every column sum is one. Each $A'_j$
has normalization genus at least two since it covers $C$. By
Riemann--Hurwitz any component covering it has at least its genus.
The genus formula for the nodal sources now bounds the sum of the
distinct normalization genera by $\sum_i h_i$.

At equal reduced degree every positive component maps with degree
one onto a different image, and each old label has at least one
positive component. Thus the number of distinct images cannot
decrease. If the entire triple is unchanged, each label has only one
positive component, with its full genus. The dual graph is a tree and
all contracted vertices are rational. Any nontrivial such tree has a
contracted leaf, contrary to stability without markings. The
limiting maps are therefore smooth and birational onto their images,
and retain the corresponding weights. Finally, a Hausdorff limit of
nonempty compact connected sets is connected.
\end{proof}

\subsection{Separation through a fixed genus}

\begin{proof}[Proof of Theorem~\ref{thm:separate}]
Consider the current constrained cycles of connected support whose
normalization genera sum to at most $K$. Suppose the values of
$\lambda$ below a given value have been treated, with finitely many
transverse multiplicity-one embedded zeros retained. Around each
of these the embedded smooth maps form an open set of the full cycle
space \cite[Proposition~2.21]{DIW}. After any one enlargement of
the family the implicit function theorem preserves this isolated zero
with its sign under sufficiently small motions of the graph and tests.
Take smaller disjoint neighborhoods of the retained zeros.

On the complement, take the least remaining value of $\lambda$. The
normalization tuples with this value over the current constrained set
form a compact set up to finite relabeling. Indeed their total cycles
have compact closure as above. Lemma~\ref{lem:weighted-specialization}
places the limit in the same bounded connected locus. Minimality
forces equality of the triples and hence convergence in the smooth
normalization charts. The same argument along sequences of
perturbations tending to the current graph has two consequences. An
earlier absent value cannot occur on this complement under sufficiently
small perturbations. Every new zero of the present value must have
its tuple in the chosen neighborhoods of the compact original tuples.
If there is no remaining value, the induction ends.

Enlarge the family using Lemma~\ref{lem:shear-jets} on this compact
set. The space of ordered smooth graph maps over the enlarged
parameter $S'$ is smooth near the tuples. To see this, first deform
their finite covers of $C$ in the smooth Hurwitz stacks. For a smooth
finite source cover $f_i$, the cokernel of
$T_{D_i}\to f_i^*T_C$ has finite length and the kernel vanishes.
After these cover deformations and the parameter in $B$ have been
prescribed, the remaining obstruction for the graphs lies in
$\bigoplus_i H^1(D_i,V_i)$, spanned by the parameter variations.
Thus the graph space submerses onto the product of the Hurwitz
stacks and $B$. It has complex dimension $\dim S'$, since the
unmarked threefold index of each component is its Chern number zero
and the target parameters occur only once. Each graph has trivial
stabilizer by birationality. This calculation only counts distinct
normalized components. Their weights enter the tests through the
cycle sum.

Mark a critical point of a cover. This marking is finite over its
Hurwitz stack, while nonimmersion requires the two vertical first
derivatives of the graph to vanish. Mark instead two distinct source
points over the same base value, on the same or on different
components. This adds at most one complex dimension over the covers.
Coincidence of the two images then imposes two vertical equations.
Take smooth strata of these marked spaces over the Hurwitz stacks
and pull back by the submersion just proved. The fixed-cover jet
variations of Lemma~\ref{lem:shear-jets} are transverse to the
equations on every such stratum at a zero. In the second calculation
a common change of target coordinates cancels in the variation of the
difference of two equal images. In the first calculation it preserves
the vanishing of the derivative at a critical point. Hence both bad
loci have dimension strictly less than $\dim S'$. The length-two
subschemes in the lemma include the limits of two marks approaching
one another.

Move the parameter graph by finitely many smooth sections spanning
its vertical normal directions along the compact base $B$. Move the
tests independently by smooth sections spanning $T$ there. The graph
condition and the tests together have real codimension $2\dim S'$,
and the auxiliary directions span their joint normal bundle even on
the marked strata. Sard's theorem on countably many smooth charts
gives arbitrarily small choices transverse on the unmarked spaces and
missing all the bad strata. We can impose transversality to the
parameter graph alone at the same time. A connected constrained tuple
of the present type now has one component, since distinct embedded
images in the tuple are disjoint. Its weight is one, since the collar
excludes proper multiples. The surviving zeros of this type are
finite by Lemma~\ref{lem:weighted-specialization}, compactness, and
transversality in the smooth charts. They are isolated in the full
cycle space on the parameter graph with the tests imposed, again by
embedded-cycle openness. Add them to the retained list. An enlargement
preserves regularity of the parameterized graph equations at the old
zeros, since the old parameter directions are still present at its
zero section. Openness then gives the same regularity at their
continuations. At an embedded curve the obstruction for the fixed
target is $H^1$ of its normal bundle. Since the parameterized graph
problem has zero obstruction, its target variations surject onto
this group. There are only finitely many values of $\lambda$ to
treat.

Take $S$ to be the final parameter scheme. Its algebraic zero sections
embed every preceding family in this last family. Concatenating the
smooth parameter paths gives the claimed homotopy. The reference
values of $\ell,\Omega$ can be moved before each small perturbation
inside $c=0$, where the target and tests are constant. The compact
joint support there is just transported along the reference path.
The successive collars and excision give a homotopy of the joint
Thom product, relative to the complement of its compact support.
A connected stable map of genus at most $K$ has connected support,
and its arithmetic genus bounds the sum of the genera of the distinct
image normalizations by Riemann--Hurwitz. It consequently maps to
one of the retained embedded curves with degree one.

Finally, maps to fibers of the smooth open $X/S$ form an open
substack of the relative stable-map stack for a projective closure.
The obstruction theory on this open, relative to the prestable source
and to $S$, is dual to the derived sections of the pullback of
$T_{X/S}$. It gives the vertical virtual class in Borel--Moore homology
after taking cycle classes. The universal domain and this obstruction
theory commute with base change at the regular algebraic section $i$.
Its refined Gysin pullback gives the relative virtual class of the
original family \cite[\S7]{BF}. The normal Thom class computes the
same Gysin map on cycle classes. For each connected source genus the inverse
image of the compact supporting set on cycles is compact by
properness of the image-cycle map \cite{WueProduct}. This proves the second
assertion and completes the proof.
\end{proof}

\section{Virtual degree and the universal class on the same family}\label{sec:comparison}

Let $X\to S$ be the smooth algebraic family of threefolds obtained in
Theorem~\ref{thm:separate}. At a transverse solution the image is a
smooth embedded curve $D\subset X_s$ of multiplicity one. Its normal
bundle $N_D$ has degree $2h-2$, where $h$ is the genus of $D$.
The deformations over $S$ constructed in \S\ref{sec:deformations}
surject onto the obstruction space $H^1(D,N_D)$. Indeed, the
deformations with the projection cover fixed already span the vertical
obstruction space, and the cover itself is unobstructed. The quotient
by deformations of the cover gives the obstruction space of the
embedded map. We first compute the contribution from all connected
maps with image cycle $[D]$.

\subsection{The relative virtual class}

Let $\mathcal H$ be an open neighborhood of $[D]$ in the relative Hilbert
scheme of smooth connected proper curves in $X/S$. It can be taken
inside the Hilbert scheme of a projective completion by restricting
to curves contained in $X$. Shrink $\mathcal H$ so that the normal obstruction
map remains surjective. Write $p:\mathcal C\to \mathcal H$ for the universal
curve and $N$ for its relative normal bundle in $X\times_S \mathcal H$.
Let $\mathcal M_k(X/S)$ be the stack of connected unmarked stable
maps of arithmetic genus $k$ to fibers of $X\to S$. We use its
vertical stable-map obstruction theory over the stack of prestable
curves times $S$, with the fundamental class of the smooth parameter
space. This gives a virtual class on the quasi-projective target
open of \S\ref{sec:deformations}.

\begin{proposition}\label{prop:moving-embedded-curve}
The scheme $\mathcal H$ is smooth of dimension $\dim S$. For $k\ge h$, let
\[
 q:\mathcal D=\overline{\mathcal M}_{k}(\mathcal C/\mathcal H,1)\longrightarrow \mathcal H
\]
be the proper stack of degree-one maps to its fibers. Composition
with the universal embedding identifies $\mathcal D$ with the open
part of $\mathcal M_k(X/S)$ over the multiplicity-one images
parametrized by $\mathcal H$. The virtual class of this open part, pushed
forward by $q$, is $b_{k,h}[\mathcal H]$, where
\begin{equation}\label{eq:embedded-degree-one}
 b_{k,h}=[u^{2k-2}]\bigl(2\sin(u/2)\bigr)^{2h-2}.
\end{equation}
For $k<h$ the inverse image of this image locus is empty.
\end{proposition}

\begin{proof}
For a smooth embedding the tangent complex relative to $S$ is
$Rp_*N$. This follows by canceling the tangent of the source in the
deformation complex of the map. The absolute obstruction at $[D]$
is the cokernel of the map $T_sS\to H^1(D,N_D)$. Its vanishing
gives smoothness, and $\chi(D,N_D)=0$ gives the asserted dimension.
In particular the derived space of smooth embedded maps over the
varying target is classical here. Its absolute tangent is $T_\mathcal H$.

Write $\pi:\mathcal F\to\mathcal D$ for the universal source,
$a:\mathcal F\to\mathcal C$ for the universal map, and
$\iota:\mathcal D\to\mathcal M_k(X/S)$ for composition. The cone of
the induced morphism of virtual tangent complexes is
\begin{equation}\label{eq:comparison-normal-cone}
 \operatorname{Cone}\bigl(q^*Rp_*N\longrightarrow R\pi_*a^*N\bigr).
\end{equation}
To see the arrow, work first over the stack of prestable sources.
With the target curve fixed the vertical tangent is
$R\pi_*a^*T_{\mathcal C/\mathcal H}$, and with the threefold fixed it is
$R\pi_*a^*(T_{X/S}|_{\mathcal C})$. Their quotient is
$R\pi_*a^*N$. Allowing the embedded curve to move adds the
triangle for $T_\mathcal H\to T_S|_\mathcal H$, whose fiber is $Rp_*N$.
A normal displacement of the embedded curve acts on the composite
by pullback under $a$. The two tangent triangles therefore give
\eqref{eq:comparison-normal-cone}. Adding the common motion of the
prestable source leaves the cone unchanged. Composition with the
universal embedding is a morphism of the derived mapping stacks.
Upon truncation it induces the maps of stable-map obstruction
theories, including their maps to the intrinsic cotangent complexes
\cite[\S\S1, 2.2]{STV}. The construction is
relative to the parameter scheme and local in the target along the
proper source.

At a geometric point $a:F\to D$ of $\mathcal D$, the normalization
of the unique positive component maps isomorphically to $D$.
It has no self-node, since two distinct points on the normalization
have distinct images. Each connected contracted subcurve attaches
to this component at just one point, for the same reason. Functions
on a contracted connected fiber are constant and take the value at
the attachment. Hence $a_*\mathcal O_F=\mathcal O_D$ and
$R^1a_*\mathcal O_F$ is a sheaf of finite length $k-h$, by the
Euler characteristics. In particular $k\ge h$. The unit
$\mathcal O_D\to Ra_*\mathcal O_F$ and the projection formula
identify the fiber of \eqref{eq:comparison-normal-cone} with
\[
 H^0(D,N_D\otimes R^1a_*\mathcal O_F)[-1].
\]
Both terms in \eqref{eq:comparison-normal-cone} are perfect and
commute with derived base change along the flat universal curves.
A perfect complex whose geometric fibers are concentrated in one
degree with constant rank is a vector bundle in that degree.
Thus \eqref{eq:comparison-normal-cone} is $W[-1]$ for a bundle
$W$ of rank $2(k-h)$ on $\mathcal D$.

The relative tangent complex of the derived composition morphism
restricts to $W[-2]$ on the classical truncation. For an ordinary
square-zero extension the relative infinitesimal automorphism,
tangent and obstruction groups in degrees $-1,0,1$ consequently
vanish. The classical morphism $\iota$ is formally \'etale and locally
of finite presentation, hence \'etale. Each image determines its
point of $\mathcal H$ uniquely, and the stabilizers on both sides consist of
the same automorphisms of the map from $F$. Thus $\iota$ is a
representable \'etale monomorphism, and identifies $\mathcal D$
with the asserted open substack.

On this common classical stack the vector bundle stack of the
obstruction theory of $\mathcal D$ embeds in the one for
$\mathcal M_k(X/S)$ with quotient $W$. By compatibility the
embedding of the intrinsic normal cone for the latter theory factors
through its embedding for the former. The zero-section Gysin maps
therefore give \cite[\S\S4--5]{BF}
\begin{equation}\label{eq:comparison-excess}
 \iota^*[\mathcal M_k(X/S)]^{\rm vir}
       =c_{2(k-h)}(W)\cap[\mathcal D]^{\rm vir}.
\end{equation}
Here the class on the right uses the smooth family of target curves
over $\mathcal H$. The equality holds on the stacks, including the stabilizers
of all contracted subcurves.

The right side of \eqref{eq:comparison-excess} has dimension
$\dim \mathcal H$. Its proper pushforward is a multiple of $[\mathcal H]$ on each
connected component. To compute the multiple, intersect with a point
$[D]\in \mathcal H$. Refined base change gives the usual virtual class of
$\overline{\mathcal M}_k(D,1)$ \cite[Proposition~7.2]{BF}.
On this fiber the class of $R\Gamma(D,N_D)$ in $K$-theory vanishes
because it is a constant class of rank zero. Consequently the
coefficient is
\[
 \int_{[\overline{\mathcal M}_{k}(D,1)]^{\rm vir}}
             c_{2(k-h)}(-R\pi_*a^*N_D).
\]
In rational cohomology Grothendieck--Riemann--Roch on the universal
source shows that this expression depends only on the rank and
degree of $N_D$. We may use $A\oplus K_DA^{-1}$ for a degree-$(h-1)$
line $A$ outside the theta divisor. Both summands are acyclic,
and a section of their pullback to $F$ is determined by its
restriction to the positive component. The integral is thus the
local degree-one contribution computed by
Pandharipande \cite[Theorem~2]{Pand}. Multiplying his series in
the excess genus by $u^{2h-2}$ gives \eqref{eq:embedded-degree-one}.
\end{proof}

\subsection{The local class in families of cycles}

Let $\sigma:B\to S$ be the transverse smooth section of
Theorem~\ref{thm:separate}. Give its normal bundle the orientation
of the vertical tangent of $S\to B$, and give the tests their complex
orientations. At a solution with image $D$, let $\epsilon(D)$ be
the local intersection sign on $\mathcal H$ of the parameter graph and the
tests. The total real codimension is $2\dim_{\mathbb C} \mathcal H$.
The parameter graph is transverse before the tests are imposed, so
its inverse image in $\mathcal H$ is the regular space of embedded curves in
the family over $B$. Its transverse-fiber-product orientation agrees
with the determinant orientation of the parametrized map equation,
by the tangent exact sequence for the graph. Interchanging the
normal to the graph and the tests gives no sign, as their real ranks
are even. Proposition~\ref{prop:moving-embedded-curve} and the
projection formula now give $\epsilon(D)b_{k,h}$ for the
unnormalized algebraic contribution at this solution.

\begin{lemma}\label{lem:signed-embedded-cycle}
In the smooth family colimit the local unnormalized test class at
this solution equals $\epsilon(D)\Phi(x_{h,1})$.
\end{lemma}

\begin{proof}
Embedded connected curves of multiplicity one form an open subset
of the space of cycles, with their embedded-map topology
\cite[Proposition~2.21]{DIW}. Thus the transverse solution can be
excised in cycle space. To identify its class we first deform the
complex structure of the marked target germ so that the curve
becomes regular in that target, keeping the curve and its test
values fixed.

Write $V=H^0(D,N_D)$ and $U=H^0(D,K_D\otimes N_D^\vee)$, both
of dimension $l$ by Riemann--Roch. If $l>0$, choose points
$p_1,\dots,p_l$ on unramified sheets over distinct points of the
base curve away from the test fibers. Choose covectors $\lambda_i$
and vectors $v_i$ in the respective normal fibers so that both maps
\[
 V\longrightarrow\mathbb C^l,\quad
       s\longmapsto(\lambda_i s(p_i))_i,\qquad
 U\longrightarrow\mathbb C^l,\quad
       \eta\longmapsto(\eta(p_i)(v_i)/dx)_i
\]
are isomorphisms. Here $x$ is a local coordinate on the base.
Inductively, nonzero sections in the two remaining evaluation
kernels are simultaneously nonzero on some unramified sheet. At
such a point choose the vector and covector to reduce each kernel
dimension by one. Extend these data in local frames. Multiply the
endomorphisms $v_i\lambda_i$ by smooth $(0,1)$-forms supported in
small disks about $p_i$, normalized to pair to one with $dx$, and
denote their sum by $A$. Under Serre duality the matrix of
$J:V\to H^1(D,N_D)$, $s\mapsto[As]$, tends to the product of the
two evaluation matrices, one transposed, as the disks shrink. Hence
it is invertible for a fixed choice of the disks. Splitting off the
invertible complement of kernel and cokernel in $\bar\partial_{N_D}$
gives the matrix $tJ+O(t^2)$ on these finite-dimensional
spaces for $\bar\partial_{N_D}+tA$. This operator is invertible for all
sufficiently small $t>0$.

Over the chosen base disks take disjoint holomorphic submersion
charts about the sheets. Use the base coordinate $x$ and vertical
coordinates $y$ in which the marked sheet is $y=y_i(x)$.
Over the support of the perturbation a sufficiently small target
tube meets only these charts. Write $A=A_i(x)d\bar x$ on the
chosen sheet and use zero on the other sheets. Replace its local
graph equation by
\[
        \partial_{\bar x}y=-tA_i(x)(y-y_i(x)).
\]
The coefficient vanishes near the boundary of the base disk, so
the change glues on the tube. Its holomorphic dependence on $y$
makes the horizontal $(0,1)$ vector field commute with the two
vertical $(0,1)$ coordinate fields. Hence the resulting almost
complex structure is integrable. Along the marked sheet the added vector
field vanishes and its normal linearization is the prescribed
change of Dolbeault operator. Extend the change smoothly over a
parameter neighborhood in $B$, with a parameter cutoff. All framed
test fibers are unchanged. For $l=0$ the fixed target was already
regular.

At $t=0$ the differential of the parametrized map equation together
with the tests is invertible in an unparametrized slice. The implicit
function theorem therefore isolates the retained solution on a small
closed interval of $t$. By the open embedding into the relative
cycle space its local Thom class extends by zero over the interval,
with the same sign. At the regular endpoint the local family of
curves projects diffeomorphically to $B$, preserving orientation
since the fixed-fiber normal operator is an invertible complex
operator. The test class is now the base Gysin image of the unit of
the isolated curve, multiplied by $\epsilon(D)$.

Deform the holomorphic embedding to its normal bundle. The curve
stays regular along this deformation, with the same normal
Dolbeault operator up to isomorphism at the positive parameters.
The implicit function theorem and embedded-cycle openness isolate
its compact section along the full interval. At the normal endpoint
a compact degree-one cycle is a section, since the vector-space
fibers have no compact positive curves. Acyclicity makes the zero
section the only such cycle. Its unit is the empty-tuple
representative of the multiplicity-one class of genus $h$ and Chern
number zero \cite[Definition~4.7]{Pardon}. The two local
deformations and the base Gysin map identify the asserted class
\cite[\S\S2--3]{WueProduct}.
\end{proof}

\subsection{Extraction by source genus}

\begin{proof}[Proof of Theorem~\ref{thm:main}]
Fix $K$ and use Theorem~\ref{thm:separate} for all connected source
genera through $K$. For each $h\le K$ let $a_h^{(K)}$ be the sum
of $\epsilon(D)$ over the finitely many connected embedded solutions
of genus $h$, divided by $W_{g,m}$. The parameter graph and the
tests form a single compactly supported intersection class on the
algebraic family, isotopic to the graph and tests at the original
algebraic section. Hence the pairing computed in
\S\ref{sec:weights} is the sum of the local pairings in
Proposition~\ref{prop:moving-embedded-curve}. For every $k\le K$
this gives
\begin{equation}\label{eq:connected-algebraic-extraction}
             N_{k,m,j}=\sum_{h\le k} a_h^{(K)} b_{k,h}.
\end{equation}
Here and below a smaller source genus cannot map to a
multiplicity-one image of larger genus.

In the smooth family over the moving graph the same test classes
give a compactly supported homotopy on cycles. For fixed $k\le K$
the image-cycle map from connected genus-$k$ stable maps is proper
\cite[\S5]{WueProduct}. Outside neighborhoods of the finitely many
embedded solutions through genus $K$, the compact constrained set
misses its closed image. The pullback of that part of the test
class is therefore zero. Over each retained solution
Lemma~\ref{lem:signed-embedded-cycle} applies before this
pullback, including when the inverse image contains contracted
components. Finally, the product terms subtracted in the logarithm
require positive images in disjoint targets and have zero connected
pullback. Thus
\begin{equation}\label{eq:connected-universal-extraction}
          T_k(P_{m,j})=\sum_{h\le k} a_h^{(K)}T_k(x_{h,1}).
\end{equation}

In the range of weights of the theorem, the primitive-division
lemma gives a finite expansion $P_{m,j}=\sum_h a_{h,j}x_{h,1}$
in the complex cycle algebra, as in \S\ref{sec:weights}. Since
$T_k(x_{h,1})=0$ for $h>k$ and $T_k(x_{k,1})\ne0$
\cite[\S\S5--6]{WueProduct}, induction on $k$ in
\eqref{eq:connected-universal-extraction} identifies
$a_{h,j}=a_h^{(K)}$ for $h\le K$. Equation
\eqref{eq:connected-algebraic-extraction} therefore holds with the
coefficients of the complex universal class for every $k$, since
$K$ was arbitrary. Substitution of \eqref{eq:embedded-degree-one}
gives the asserted sine expansion. The fixed-genus residue of
\S\ref{sec:weights} vanishes for $k<1+m(g-1)+|j|$, and
$b_{h,h}=1$. The same induction in
\eqref{eq:connected-algebraic-extraction} gives the stated lower
bound on the genus of a nonzero coefficient.
\end{proof}

\section{Cubic evaluations on a general normal bundle}\label{sec:evaluations}

Let $D$ be a smooth projective curve of genus $g\geq2$, and put
$k=\lfloor(g-1)/2\rfloor$. For a rank-two bundle $N$ over $D$, a point
$p\in D$ and a covector $\ell\in N_p^\vee$, let $F_{p,\ell}$ be the
product of the three values of $\ell$ on the fiber of an effective
degree-three cycle in $\operatorname{Tot}_D N$, counted with multiplicity.
More generally, Pardon allows any homogeneous symmetric polynomial of
positive degree in these values \cite[(4.15)]{Pardon}. Its degree is its
weight under scalar multiplication on $N$.

\begin{theorem}\label{thm:cubic-evaluations}
For a general stable $N$ with determinant $K_D$, the common zero of $k$
general point products is $3[D]$. For every rank-two $N$ of degree $2g-2$,
any tuple of the above evaluations with compact common zero set has at
least $k$ weights divisible by three. In particular the least possible
product of weights is $3^k$, and $3^k x_{g,3}$ is integral.
\end{theorem}

\begin{proof}
Write $K=K_D$. The stable fixed-determinant moduli space $M$ is irreducible
of dimension $3g-3$. On a dense open its bundles are $2$-rigid
\cite[Theorem~2.2 and \S3]{BP06}: the pullback to each stable map
of degree two has no nonzero section. Thus $H^0(N)=0$, and for every line
subbundle $S\subset N$ one has $H^0(S^2)=0$. Indeed a rational square
root of a regular section is regular, and the normalization of the spectral
double of a nonsquare carries a nonzero section of the pullback of $N$.
There are no compact curves in a fiber of $N$. Consequently any cubic
cycle distinct from $3[D]$ is integral with multiplicity one. Its
normalization gives a smooth connected triple cover $f:F\to D$. The
normal coordinate on $F$ gives a nonzero map $T\to N$, where
$T$ is the Tschirnhausen bundle defined by the trace splitting $f_*\cO_F=\cO_D\oplus T^\vee$.

We first exclude a generically injective map $T\to N$. Write $b=\deg T$. The discriminant divisor $\Delta$
has degree $2b$ and satisfies $\cO(\Delta)=(\det T)^2$ \cite[\S\S3--4]{Miranda}.
An injection $T\to N$ gives an effective determinant divisor $E$ of
degree $n=2g-2-b$, with $\Delta+2E\in|K^2|$. There are finitely many
covers for fixed $\Delta$, by monodromy on the punctured curve. For fixed
$T,E$ the upper modifications have at most $n$ parameters, by successive
elementary modifications of the dual at the specified points. Stratify
the divisors by the multiplicities at their $t$ distinct support points.
Unless these points lie on a canonical divisor, the Abel map for
$\Delta+2E$ has differential of rank $g$, dual to weighted evaluation of
canonical forms. Thus the dimension, including modifications, is at most
$t-g+n\leq2b+2n-g=3g-4$. On a canonical divisor the support has only
finitely many choices once the divisor is given, yielding the bound
$g-1+n=3g-3-b$. This suffices for $b>0$. If $b=0$ the cover is
\'etale. For a noncyclic triple, $\det T$ is its nontrivial sign
character, so $E\in|K(\det T)^{-1}|$ has $g-2$ parameters and the total
is at most $3g-4$. For a cyclic triple $T$ splits as two lines of order
three. These estimates involve finite-type Hurwitz and modification
spaces since $0\leq b\leq2g-2$.

For a saturated subline $S$ of degree $d\leq g-2$ the projective space
of extensions by $S$ with quotient $KS^{-1}$ has dimension $3g-4-2d$,
by Riemann--Roch for $S^2K^{-1}$. Stability excludes split extensions
and larger degrees. Letting $S$ vary adds $g$ parameters, so a general
$N$ has no subline of degree $d>k$. A map from a prescribed degree-zero
line $R$ would saturate in $R(A)$, with $a=\deg A\geq0$. The divisor
and extension have at most $a+3g-4-2a<\dim M$ parameters. This excludes
all lines of order dividing three and finishes the rank-two case.

For rank one, let $S$ be the saturation of the image. The normal
coordinate in $f^*S$ has trace zero and generates the cubic field since
the degree is prime. Its characteristic equation reduces to $y^3=s$,
with $s\in H^0(S^3)$ nonzero, since its coefficient in $H^0(S^2)$
vanishes. Hence $1\leq d\leq k$, degree zero having been excluded.
Conversely each such $s$ defines an integral cubic, since a rational
cube root would give a section of $S$. Its generic fiber recovers $S$
by saturation and then $s$ by the equation.

The cubics modulo dilation thus form a projective algebraic set
$Z\subset\PP H^0(D,\operatorname{Sym}^3N)$. Indeed its affine
cone consists of the tensors whose value at every point lies in the
closed cone of pure cubes. A nonzero such tensor determines a rational
line direction by the inverse Veronese map. In a frame adapted to its
saturation $S$, the tensor belongs to $S^3$ across its zeros as well.
To bound $\dim Z$, the pairs $(S,\operatorname{div}s)$ of degree $d$
are parametrized by the finite \'etale pullback of $\operatorname{Sym}^{3d}D$
under cubing on Picard varieties. The projective extension family over
these pairs has dimension $3d+3g-4-2d$. The stable part maps to $M$.
Over general $N$ the fibers have dimension at most $d-1$ and parametrize
the degree-$d$ locus in $Z$, since inclusions up to scalar specify the
subbundle. It follows that $\dim Z\leq k-1$, with $Z$ empty for $k=0$.

On this cone the product $F_{p,\ell}$ is the linear form
$q\mapsto\ell^{\otimes3}(q(p))$. For fixed $[q]\in Z$ its zeros in the
surface $\PP_D(N^\vee)$ form a proper closed subset. The incidence
of $k$ simultaneous zeros above $Z$ therefore has dimension at most
$k-1+k<2k$. Its projection is closed by projectivity of $Z$, so general
products cut out $3[D]$. Their joint map is proper. Indeed its norm is
bounded below on the unit sphere of the coefficient cone. Bounded
coefficients bound the cubic roots in a compact disk bundle. In
$\PP_D(\cO\oplus N)$ all these cycles have the degree of
$3[D]$ by dilation, and their projective Chow limits remain in the disk.
Over the Borel base $\PP^k$ their joint
integral Thom class represents $3^kx_{g,3}$ by
\cite[Definition~4.7 and (4.16)]{Pardon}. The empty tuple over a point
applies when $k=0$.

For the lower bound suppose $k\geq1$, take any $N$ of degree $2g-2$,
and prescribe an effective divisor $P$ of degree $k-1$. We claim that
some $S\in A=\operatorname{Pic}^kD$ admits both a nonzero map $S\to N$
and a nonzero section of $S^3(-P)$. Let $p:D\times A\to A$ be the
projection, $\mathcal P$ the normalized Poincar\'e bundle and $\theta$
the principal polarization. For $V=-Rp_*(N\otimes\mathcal P^{-1})$,
Grothendieck--Riemann--Roch gives
\[
 \operatorname{ch}(V)=2k+2\theta,\qquad c_j(V)=(2\theta)^j/j!.
\]
Indeed $\operatorname{ch}(N)\operatorname{td}(D)=2$, and the square of
the mixed Poincar\'e class pushes forward to $-2\theta$. On a closed
locus where all maps $S\to N$ vanish, the differential in a two-term
presentation is fiberwise injective. Its cokernel represents $V$ there
by a rank-$2k$ bundle, so $c_{2k+1}(V)$ restricts to zero.

If $g=2k+1$, the integral of this class on $A$ is $2^g$, and
$\deg S^3(-P)=g$ gives a section by Riemann--Roch. If $g=2k+2$, let
$T$ be the inverse image of the theta divisor $W_{g-1}(D)$ under
$S\mapsto S^3(-P)$. Translation and cubing give $[T]=9\theta$, whence
$\int_T c_{2k+1}(V)=9g\,2^{g-1}>0$. Some $S\in T$ therefore admits
a map, and $S^3(-P)$ has a section by definition. In either case push
the resulting spectral cubic into $N$ by $S\to N$. It is a compact
effective cycle whose generic fiber lies off the zero section, even if the line map has
zeros. Over each point of $P$ the entire cubic lies on the zero section.

The three values of any fiber covector on this cubic are
$z,\zeta z,\zeta^2z$, where $\zeta$ is a primitive cube root of unity.
A homogeneous symmetric polynomial of weight prime to three
vanishes on them, since multiplication by $\zeta$ permutes them. If
fewer than $k$ weights are divisible by three, include their points in
$P$, completing it to degree $k-1$. Every evaluation then vanishes on
the entire dilation orbit of this cubic. At a base point where the three normal coordinates are nonzero
some point product is unbounded on the orbit, contradicting compactness
of the common zero set.
\end{proof}

\begin{corollary}\label{cor:cubic-fiber-count}
Fix $1\leq d\leq k$ and $d-1$ distinct points with sum $P$. For general
stable $N$ of determinant $K$ the number of pairs $(S,[s])$ with
$S\subset N$ a degree-$d$ subline and $0\ne s\in H^0(S^3(-P))$ is
\[
        2^{2d+1}9^{g-2d-1}\binom{g}{2d+1}.
\]
\end{corollary}

\begin{proof}
Put $r=2d+1$. The pairs $(S,B)$ with $S^3(-P)=\cO(B)$,
$B\in\operatorname{Sym}^rD$, form a smooth finite \'etale pullback $W$
of that symmetric power. Its fundamental class pushes to
$9^{g-r}\theta^{g-r}/(g-r)!$ on $\operatorname{Pic}^dD$ by the
Poincar\'e formula \cite[Introduction]{Debarre}. The preceding index
calculation with $d$ in place of $k$ gives $\int_W c_r(V)$ equal to the
displayed number. Positivity forces a nonzero map from some $S$.
For general $N$ each such map is saturated, since a saturation $S(A)$
with $a=\deg A>0$ has at most $r+a+(3g-4-2d-2a)=3g-3-a$ parameters
including the extension. The saturated extension incidence is smooth
of dimension $3g-3$, so its general fibers over $M$ are finite and
reduced by generic smoothness. Thus every kernel at a rank drop has
dimension one, since its projectivization parametrizes the maps up to
scalar. An invertible minor in the two-term presentation identifies
the projective kernel incidence locally with the rank-drop scheme.
Its codimension is $r$ and the Thom--Porteous formula gives its length
as $\int_W c_r(V)$.
\end{proof}

\begin{corollary}\label{cor:any-triple}
For every smooth curve $D$ of genus $g\ge3$ and every holomorphic
rank-two bundle $N$ of degree $2g-2$, there is a connected smooth
degree-three cover $f$ with $H^0(f^*N)\ne0$ and genus at most $3g+1$.
If $H^0(N)=0$, the corresponding cyclic images in the normal bundle
are integral and their nonzero dilates approach $3[D]$.
\end{corollary}
\begin{proof}
On $\Pic^1D$ use the degree-one Picard complex in the preceding proof.
The inverse image $W$ of the effective locus $W_3$ by cubing has
$\int_Wc_3(V)=8\cdot9^{g-3}\binom g3>0$. If the kernel were
everywhere zero, the rank-two actual cokernel bundle on $W$ would
have zero third Chern class. A line $S$ with both $S\to N$ nonzero
and a nonzero section $s$ of $S^3$ therefore exists, without any
saturation or generality assumption. If $s$ is a meromorphic cube
its cube root is holomorphic and gives a section of $N$, which
pulls back to any connected \'etale cyclic triple. Otherwise normalize
the cyclic cubic $y^3=s$. At most three zeros have order nonzero
modulo three and each contributes two to ramification. Its genus
is at most $3g-2+3$. The generic injection from $S$ to $N$ embeds
the generic fiber multiset, and dilation then gives the claimed
cycle limit.
\end{proof}

\section{The isolated triple in genera three and four}\label{sec:covers}
Let $D$ be a smooth projective curve of genus $h$ and $K=K_D$.
\begin{theorem}\label{thm:classification}
Suppose $h=3$ or $4$, and let $N$ be a general stable bundle of
rank two with determinant $K$. The smooth connected degree-three
covers $f:F\to D$ for which $H^0(F,f^*N)\ne0$, up to isomorphism
over $D$, are in bijection with the degree-one subbundles
$S\subset N$ such that $H^0(D,S^3)\ne0$. There are eight if $h=3$
and $288$ if $h=4$. Each cover is cyclic, has source genus $3h+1$,
and satisfies $h^0(F,f^*N)=1$.

The effective degree-three cycles in $\operatorname{Tot}_D N$ are
$3[D]$ and the smooth embedded curves
\[
                   y^3=a s,\qquad a\in\mathbb C^*,
                   \qquad y\in\operatorname{Tot}_D S ,
\]
where $s$ is a generator of $H^0(D,S^3)$ for one of these subbundles.
The resulting parameter lines meet only at $a=0$, and scalar
multiplication by $\lambda$ acts on them by $a\mapsto\lambda^3a$.
For $h=2$ a general $N$ with the same determinant condition is
$3$-rigid and its only effective degree-three cycle is $3[D]$.
\end{theorem}


\begin{proof}
The spectral description in the proof of Theorem~\ref{thm:cubic-evaluations}
leaves only degree-one lines, since $\lfloor(h-1)/2\rfloor=1$ for
$h=3,4$. Corollary~\ref{cor:cubic-fiber-count} gives the two counts.
On the finite \'etale pullback of $\operatorname{Sym}^3D$, the loci of
multiple zeros and of nonunique effective divisors in a complete
linear system have dimension at most two. Indeed three general points
impose independent conditions on differentials. Adding the extensions
of $KS^{-1}$ by $S$ leaves dimension at most $3h-4$ over these loci.
Thus $h^0(S^3)=1$ and the divisor is reduced. Riemann--Hurwitz gives
$3h+1$. Each point of this finite incidence has a unique subbundle map
up to a constant. The quotient of its Tschirnhausen bundle $S\oplus S^2$ by the unique
degree-two subline recovers $S$, and its inclusion in $N$ is unique up
to scalar. Thus distinct sublines give distinct covers over $D$.
To identify the normal sections on a given cover, write $f_*\cO_F=\cO_D\oplus S^{-1}\oplus S^{-2}$. The last summand
tensored by $N$ has no sections because $N$ has no subline of degree
two or more. The middle summand gives the unique inclusion. Finally at
$h=2$ the cubic space is a point. A nonzero section on a connected
stable map would be nonzero on some normalized nonconstant component
since sections on attached complete contracted curves are constant.
The vanishing just proved on smooth covers through degree three
excludes this possibility.
\end{proof}
Put $r_3=8$, $r_4=288$ and $H=3h+1$.
An integral insertion on an open and closed compact subset of a
cycle space is the extension by zero of its unit. Products use disjoint
targets. In genus $h$ let $x_{h,d}$ be the rational universal class
associated to degree $d$ in the normal bundle with scalar action on
both directions \cite[Definition 4.7]{Pardon}. For $d=1$ this is
the integral class of an isolated unobstructed embedded curve.
If $D$ is embedded in a smooth quasi-projective threefold $X$ and
$d[D]$ is isolated in a tube about $D$, denote its integral unit
by $J_d$. The local degree here is the coefficient of $[D]$ under
retraction of the tube to $D$. Excision identifies the unit in
every smaller tube.

In source arithmetic genus $k$, let $G_k(X,3D)$ be the ordinary
connected Gromov--Witten degree on the maps with this image cycle,
once it is isolated. For comparison let $S^{\rm sc}_{h,3}(u)$ be the
connected degree-three residue of $\Tot N$ for scalar multiplication
on the fibers, with $u^{2k-2}$ recording the source genus. It can be
evaluated on the proper stack of stable maps to $D$ by equivariant
localization \cite{BP}. The distinguished rational universal class
$x_{h,3}$ is the coefficient before taking a logarithm, whereas the
numerical series uses connected domains throughout.

\begin{theorem}\label{thm:isolated}
Let $h=3$ or $4$, let $N$ be general with determinant $K_D$, and let
the embedding $D\subset X$ in a smooth quasi-projective threefold
have normal bundle $N$ and a second jet whose obstruction at each of the finitely many dilation lines is nonzero, as defined in Section~\ref{sec:neighborhood}. Then $3[D]$ is isolated in a neighborhood of $D$. Such
neighborhoods exist in smooth algebraic threefolds. In the rational
complex universal algebra,
\begin{equation}\label{eq:main-class}
           J_3=x_{h,3}+\frac{r_h}{3}x_{H,1}.
\end{equation}
In the fixed-target Gromov--Witten theory,
\begin{equation}\label{eq:main-gw}
   \sum_k G_k(X,3D)u^{2k-2}
      =S^{\rm sc}_{h,3}(u)+\frac{r_h}{3}(2\sin(u/2))^{2H-2}.
\end{equation}
In degrees one and two one has $J_d=x_{h,d}$.
\end{theorem}

\section{Isolation of the tripled curve}\label{sec:neighborhood}

Let $D$ have genus $h=3$ or $4$, and let $N$ be a general stable
bundle of rank two with determinant $K_D$. Write
$Y=\operatorname{Tot}_D N$. By Theorem~\ref{thm:classification},
the degree-three cycles in $Y$ away from $3[D]$ are smooth embedded
curves in finitely many orbits of fiber dilation. Each orbit comes
from a line subbundle $S\subset N$ of degree one. For a generator
$s\in H^0(D,S^3)$ its members are the curves
\[
                 F_a=\{y^3=a s\}\subset\operatorname{Tot}_D S,
                      \qquad a\in\mathbb C^\times.
\]
The three zeros of $s$ are simple, and $H^0(S)=H^0(S^2)=0$.
We first compute the obstruction to deforming $F_a$ when the
threefold varies.

\subsection{The normal bundle of a cubic}

It suffices to take $a=1$. Put $F=F_1$, write $f:F\to D$ for the
projection and $Q=K_DS^{-1}$, and denote the class of the extension
of $Q$ by $S$ defining $N$ by $e\in H^1(SQ^{-1})$.
The equation of $F$ in $\operatorname{Tot}_D S$ gives the normal
sequence
\begin{equation}\label{eq:cubic-normal-sequence}
 0\longrightarrow f^*S^3\longrightarrow N_{F/Y}
               \longrightarrow f^*Q\longrightarrow0 .
\end{equation}
Its connecting map on global sections is cup product with
$3y^2f^*e$, up to sign. Indeed, local splittings of
$N\to Q$ lift a $Q$-displacement to vertical vector fields whose
differences are the corresponding $S$-displacements given by $e$.
The derivative of the cubic equation on these differences is
multiplication by $3y^2$.

\begin{lemma}\label{lem:cubic-normal-cohomology}
For all the cubics of Theorem~\ref{thm:classification}, after shrinking
the dense open of normal bundles if necessary, the tangent
$H^0(F,N_{F/Y})$ is generated by varying $a$, and projection in
\eqref{eq:cubic-normal-sequence} induces an isomorphism
\[
        H^1(F,N_{F/Y})\ \xrightarrow{\ \sim\ }\ H^1(F,f^*Q)
           = H^1(D,QS^{-2}) \cong H^0(D,S^3)^\vee .
\]
\end{lemma}

\begin{proof}
Use $f_*\mathcal O_F=\bigoplus_{i=0}^2 S^{-i}$, with the $i$th
summand represented by $y^i$ in a local frame. The three spaces
$H^0(QS^{-i})$ have dimensions $h-2,h-3,h-3$ by Riemann--Roch.
The first maps isomorphically to $H^1(S)$ in the $y^2$ summand of
$H^1(f^*S^3)$. This is three times the connecting map for $N$,
which is acyclic by degree-two rigidity
\cite[Theorem~2.2]{BP06}. In genus three these are the only nonzero
blocks.

In genus four take generators $q_1$ and $q_2$ of $H^0(QS^{-1})$
and $H^0(QS^{-2})$. Their blocks map to $H^1(S^3)$ and $H^1(S^2)$,
respectively. Since $y^3=s$, both Serre pairings are, up to the
common factor and sign, $\langle e,sq_1q_2\rangle$. The product is
a nonzero section of $Q^2$, the Serre dual of the extension space.
Its annihilator is a hyperplane in each projective space
$\mathbb P H^1(SQ^{-1})$ of dimension six. The possible $S$ with
a cubic section range over a three-dimensional locus in
$\operatorname{Pic}^1D$. These exceptional extensions therefore
have image of dimension at most eight in the nine-dimensional
moduli space of bundles with fixed determinant.

The connecting map in \eqref{eq:cubic-normal-sequence} is consequently
an isomorphism in both genera. The remaining tangent space is
$H^0(f^*S^3)=H^0(S^3)$, spanned by the derivative of $F_a$.
Serre duality gives the asserted obstruction space, since its other
two summands are dual to $H^0(S)$ and $H^0(S^2)$.
\end{proof}

\subsection{The second infinitesimal neighborhood}

Let $D\subset X$ be an embedding in a smooth quasi-projective
threefold, with an identification $N_{D/X}=N$. Write
$\mathcal X\to\mathbb A^1$ for its deformation to the normal cone, obtained by
blowing up $D\times0$ in $X\times\mathbb A^1$ and removing the proper
transform of $X\times0$. The central fiber is $Y$ and every
noncentral fiber identifies with $X$. For an embedded cubic
$F\subset Y$, let
\begin{equation}\label{eq:cubic-unfolding}
        \kappa_F(X)\in H^1(F,N_{F/Y})
\end{equation}
be the restriction of the Kodaira--Spencer class of this deformation
at zero, projected to the normal bundle of $F$. This is the
obstruction to lifting $F$ over the tangent vector $\partial_t$ of
the parameter line.

The class depends only on the second infinitesimal neighborhood
defined by $\mathcal I_D^3$. Indeed, in adapted coordinates $(z,n)$
along $D$, write a change of charts as
\[
  z'=h_0(z)+A(z)n+O(n^2),\qquad
  n'=G(z)n+B(z)(n,n)+O(n^3).
\]
In normal-cone coordinates replace $n$ by $tn$ and divide $n'$ by
$t$. The derivative at zero has horizontal component $A(z)n$
and vertical component $B(z)(n,n)$. Projecting this cocycle to
$N_{F/Y}$ gives \eqref{eq:cubic-unfolding}. The definition thus
applies to embeddings without a retraction onto $D$. Dilation
lifts to the normal-cone family, acting on its parameter by the
inverse character. Nonvanishing on one member of a dilation orbit
is equivalent to nonvanishing on every member.

\begin{proposition}\label{prop:isolation}
For $h=3,4$, suppose that $\kappa_F(X)\ne0$ for one cubic in
each orbit of Theorem~\ref{thm:classification}. Then $3[D]$ is
isolated among all compact effective cycles of local degree three
in a tube about $D$ in $X$. This condition on the second
infinitesimal neighborhood is Zariski open in algebraic families
with identified $D$ and $N$. It is attained by smooth
quasi-projective threefolds with normal bundle $N$ and
determinant-preserving changes of the fiber coordinates.
\end{proposition}

\begin{proof}
In a family with fixed $D$ and $N$, the restricted
Kodaira--Spencer classes in \eqref{eq:cubic-unfolding} vary
algebraically in the fixed obstruction lines. There are only
finitely many conditions. We exhibit an algebraic family where
each vanishing condition cuts out a proper hyperplane.

Start with $Y$. Fix a point $p\in D$, a frame of $N$ on an affine
neighborhood $U$, and a local parameter $u$ regular on $U$ with
no other zero there. Choose a covector $\ell$ on the framed plane
which is nonzero on every one of the finitely many lines $S_p$,
and a nonzero $v_\ell\in\ker\ell$. Reglue the plane bundle over
$U$ to its restriction over $D\setminus\{p\}$ by composing the
transition on the overlap with
\begin{equation}\label{eq:normal-quadratic-shear}
        n\longmapsto n+v_\ell\,\ell(n)^2
                         \sum_{j=1}^m c_j u^{-j}.
\end{equation}
This polynomial automorphism has determinant one and its inverse
negates the coefficients. It fixes the zero section to first order.
The resulting smooth threefold is affine of finite type over $D$,
hence quasi-projective. Its projection to $D$ has the same vertical
tangent determinant as the original bundle, since the added
transition has determinant one.

For a given $S$ the restriction of $v_\ell\ell^2$ to $S^2$,
followed by projection to $Q$, generates $QS^{-2}$ at $p$.
The Kodaira--Spencer class on the cubic therefore projects to the
principal-parts boundary in $H^1(QS^{-2})$, in the $y^2$ summand of
$H^1(f^*Q)$. For sufficiently large $m$, these boundaries span
that space by $H^1(QS^{-2}(mp))=0$. Multiplying the principal
parts by the local generator preserves their span. Choose $m$
large enough for all the $S$. By Lemma~\ref{lem:cubic-normal-cohomology}
each $\kappa_F$ is then a nonzero linear functional of the $c_j$,
with values in a one-dimensional space. Arbitrarily small choices
outside the union of their kernels satisfy the hypothesis.

Now fix an $X$ satisfying this hypothesis. At a nonzero cubic $F$
the tangent sequence for the relative Hilbert scheme of
$\mathcal X/\mathbb A^1$ is exact through
\[
 0\longrightarrow H^0(N_{F/Y})\longrightarrow T_{[F]}\operatorname{Hilb}
   \longrightarrow T_0\mathbb A^1
   \xrightarrow{\ \kappa_F\ }H^1(N_{F/Y}).
\]
Indeed, local lifts of the embedding over a first-order parameter
differ on overlaps by precisely the normal cocycle
\eqref{eq:cubic-unfolding}. The Hilbert tangent has dimension one.
The actual orbit $F_a$ has dimension one in the central fiber
and its derivative generates that tangent. Equality of the local
dimension and embedding dimension gives a regular local ring.
Thus the relative Hilbert scheme near $[F]$ consists entirely of
the central orbit. In particular this conclusion concerns every
deformation of the embedded curve.

Suppose extra cycles converged to $3[D]$ in $X$. In a smooth
tubular identification let $t_i>0$ be the ratio of the maximum
normal distance of the $i$th cycle to a fixed small radius. Then
$t_i\to0$. Regard the cycles in the fibers $\mathcal X_{t_i}$,
using the divided normal coordinates. They lie in a common
closed disk and reach its prescribed radius. The smooth structures
after this rescaling converge on a larger disk to the complex
structure of $Y$. A closed K\"ahler form on the central disk tames
these structures for large $i$. Its integral on every cycle is
three times its integral on $D$ by retraction to the zero section.
It follows that the masses are uniformly bounded. Cycle compactness,
including Hausdorff convergence of support, gives a limit of degree
three meeting the prescribed positive radius
\cite[Theorem~2.11 and Remark~2.13]{DIW}. Theorem~\ref{thm:classification}
identifies the limit with a single smooth cubic of multiplicity one.

Openness of the embedded locus and its topology in the cycle space
give smooth embedded nearby curves of the same genus
\cite[Proposition~2.21]{DIW}. To pass to the relative Hilbert
scheme, use a projective completion of the normal-cone family.
The ample degree of the cycles is locally constant, so the Hilbert
polynomial of the nearby embedded curves agrees with that of $F$.
Any flat Hilbert limit has support $F$ and generic multiplicity
one by convergence of supports and cycles. The kernel of its
structure sheaf onto $\mathcal O_F$ has finite length, hence
vanishes by equality of Hilbert polynomials. The resulting Hilbert
limit $[F]$ contradicts the central-orbit description above.
\end{proof}

\section{Comparison with the scalar local curve}\label{sec:triplecomparison}

Let $D\subset X$ satisfy Proposition~\ref{prop:isolation}, with $h=g(D)$
equal to three or four. Write $N=N_{D/X}$, and put
\[
                 H=3h+1,\qquad r_3=8,\qquad r_4=288 .
\]
Theorem~\ref{thm:classification} gives $r_h$ families of embedded
genus-$H$ curves in the total space of $N$. Write $\Phi$ for the map from complex to smooth families of
threefolds in the universal cycle algebra. We compare the scalar
insertion at $3[D]$ in this bundle with the compact open class $J_3$
of the isolated cycle in $X$. Write $G_k(X,3D)$ for the local virtual
degree of connected unmarked stable maps of arithmetic genus $k$ with
image cycle $3[D]$. It includes maps with contracted subcurves.
The inverse image of the isolated cycle is open and closed
in the degree-three stable-map space over a sufficiently small tube.
It is proper by fixed-genus compactness. The ordinary algebraic
stable-map virtual class restricts to it, using a projective completion
and the open set of maps with image in $X$.

\begin{proposition}\label{prop:borel}
Let $w=(2\sin(u/2))^2$, and let $S^{\rm sc}_{h,3}(u)$ be the connected
scalar residue in degree three. In the smooth cycle algebra, and in the
fixed-target connected Gromov--Witten theory, respectively, one has
\begin{align}
 \Phi(J_3)&=\Phi(x_{h,3})+\frac{r_h}{3}\Phi(x_{H,1}),
           \label{eq:cubic-smooth-comparison}\\
 \sum_k G_k(X,3D)u^{2k-2}
     &=S^{\rm sc}_{h,3}(u)+\frac{r_h}{3}w^{H-1}.
           \label{eq:cubic-residue-comparison}
\end{align}
\end{proposition}

\begin{proof}
\emph{The scalar insertion.}
Choose $p\in D$ outside the branch divisors of the $r_h$ families and
a linear form $\ell$ on $N_p$ nonzero on each of their defining
sublines. For a degree-three cycle let $F$ be the product of the
three values of $\ell$ over $p$, with multiplicities. On a family
$y^3=a s_i$ of Theorem~\ref{thm:classification}, this product is
$c_i a$ with $c_i\ne0$. Thus its unique zero is $3[D]$. Under
scalar dilation $F$ has weight three.

Take $B=\mathbb P^1$ and $L=\mathcal O_B(1)$. In the diagonal Borel
quotient a function of weight one is a section of $L$. Indeed a
pair $(v,z)$ defines the functional on $\mathbb Cv$ taking $v$ to
$z$, which is unchanged under simultaneous scaling. The function
$F$ therefore gives a section of $L^3$ on the relative cycle space.
The pullback of its Thom class divided by three, over the oriented
base $B$, represents $x_{h,3}$ by the scalar construction
\cite[\S4.3]{Pardon}.

We identify its algebraic virtual pairing before deforming the
insertion. Compactify the normal bundle by
$\mathbb P_D(\mathcal O\oplus N)$ and extend the supported class by
zero. The only fixed maps meeting its support are maps to $D$. On
their proper connected moduli stack
$\overline{\mathcal M}_k(D,3)$, with universal curve $\pi$ and map
$f$, the moving virtual normal complex is $R\pi_*f^*N$ with scalar
character one. If $\eta$ is the equivariant parameter, the
normalized test restricts to $3\eta/3=\eta$. The virtual localization
formula \cite{GP} gives the equivariant pushforward
\[
              \eta\,[u^{2k-2}]S^{\rm sc}_{h,3}(u).
\]
Here the inverse moving Euler class is exactly the fixed-target
residue of \cite[\S2.2]{BP}. It applies to $N$ since its
$K$-class on $D$ is determined by rank and determinant and its
scalar character is common to both normal directions. Passing to the
associated bundle over $B$ substitutes $c_1(L)$ for $\eta$, and
$\int_Bc_1(L)=1$. This computes the pairing of the relative
algebraic virtual class with the test in source genus $k$: on a
trivialization of the frame torsor the universal source, universal
map and their relative obstruction complex are pulled back from the
fixed-target equivariant data, and these identifications descend
with the scalar character.

\emph{The deformation.}
Form the deformation to the normal cone of $D\subset X$ and delete
the proper transform of $X\times\{0\}$. Denote its parameter by $t$.
The fibers over $0$ and over $t\ne0$ are $N$ and $X$, respectively.
Give $t$ weight $-1$ and fix the original $X$ coordinate. In normal
coordinates $n=y/t$ the action at $t=0$ is scalar of weight one.
The associated family of threefolds thus lies over
$E=\operatorname{Tot}_B(L^{-1})$. Its restriction to the zero
section is the preceding Borel family. Over the complement of that
section the original $X$ coordinate trivializes the family.

The test extends over a neighborhood of the marked cycle in this
family. Choose adapted local coordinates $(x,y_1,y_2)$ near $p$ with
$D=\{y_1=y_2=0\}$ and take the transverse surface $x=x(p)$.
For a nearby cycle its three intersections with this surface near $p$,
counted with multiplicity, give a continuous unordered triple. In
the deformation coordinates use the product of their values under
$\ell(y/t)$. The divided coordinates extend over $t=0$, and the
product has weight three. It restricts there to $F$ and vanishes on
the marked triple in every fiber. On smooth embedded curves transverse
to the surface the evaluation is holomorphic.

Choose a small section $v\in H^0(B,L^3)$ with three simple zeros.
Choose a small smooth section $\tau$ of $L^{-1}$ with a single
transverse zero $b_*$ disjoint from those of $v$. The metric dual of
a holomorphic section of $L$ with a simple zero gives such a section
of local index $-1$. On the zero section first deform $F$ to $F-v$.
Keep this translated test fixed and move the parameter over the
graphs $t=s\tau(b)$, $0\le s\le1$.

These intersections have compact support through both homotopies.
Use a closed normal disk of small radius over $D$. At $t=0$ the
maximum normal radius on the $i$th curve is a positive constant
times $|a|^{1/3}$. Hence $|F|$ is bounded below away from zero
where a cycle reaches the boundary of this disk, uniformly over
$B$ in unitary frames and over the finitely many families. Choose
$\|v\|$ smaller than this bound. If $t\ne0$, the disk maps into
an $O(|t|)$ tube about $D$ in the original target. For $|t|$
sufficiently small its only degree-three cycle is $3[D]$ by
Proposition~\ref{prop:isolation}. Within the closed disk, a common
closed taming form in each local smooth trivialization bounds the
masses by the fixed homology class $3[D]$. Relative cycle compactness
on a slightly larger disk therefore includes any limit across
$t=0$ \cite[Theorem~2.11]{DIW}. These choices can be made uniform
on the compact base. In particular the compactly supported
cohomology class on cycles is transported in the smooth colimit
\cite[\S2]{WueProduct}. For the algebraic pairing use the relative
virtual class over the smooth algebraic base $E$ and intersect it
with the Thom class of the oriented parameter graph and the test.
The inverse images of the compact supports are proper at each fixed
$k$. The homotopic graphs give equal pairings. At the zero section
the refined fiber restriction is the virtual class of the Borel
family by base change \cite[Proposition~7.2]{BF}.

\emph{Intersections on the final graph.}
Over each of the three zeros of $v$, the parameter $\tau$ is nonzero.
The only supported cycle is the isolated triple in $X$. Its local
intersection is the simple zero of $-v$, of positive index. It gives
$\Phi(J_3)$ in the smooth colimit and $G_k(X,3D)$ in the relative
virtual pairing, since the family is a product there.

Over $b_*$ there is one solution $a=v(b_*)/c_i\ne0$ on each
normal family. The local relative Hilbert scheme is smooth with
coordinates $(b,a)$ and maps to $(b,0)\in E$. Indeed the only
fixed-fiber tangent direction is the variation of $a$, and the
$t$ direction surjects onto the one-dimensional fixed-fiber
obstruction space by Proposition~\ref{prop:isolation}. The
normal curves themselves supply the smooth two-dimensional family
in $t=0$. On this chart the graph equation is $-\tau(b)=0$,
and the remaining equation has nonzero complex derivative $c_i$
in $a$. Their intersection sign is $-1$. The graph has the
orientation of the vertical complex line in $E$, and the test
has its complex orientation. Interchanging their normal factors
preserves orientation.

The Hilbert chart here satisfies the hypotheses of
Proposition~\ref{prop:moving-embedded-curve}. In particular the
motion of $t$ surjects onto the fixed-curve obstruction and the
graph is transverse before testing the amplitude. The shear argument
of Lemma~\ref{lem:signed-embedded-cycle} is local along the embedded
image. At the finitely many supporting points choose disjoint smooth
submersion coordinates away from the evaluation and ramification
fibers. These coordinates extend over a smaller parameter neighborhood
by smooth local trivialization. The required normal variations are
compactly supported there; a projection from the whole threefold to
the original curve is unnecessary. On the cycle space the lemma gives
$-\Phi(x_{H,1})$ for each such image. The proper degree-one map stack
over the same chart accounts for the virtual contribution in every
source genus, equal by \eqref{eq:embedded-degree-one} to
$-[u^{2k-2}]w^{H-1}$. The embedded image has trivial automorphisms
as a map to the threefold; contracted components of maps to it retain
their ordinary stack multiplicities.

Summing the two sets of local intersections before the scalar
normalization gives
\[
        3\Phi(J_3)-r_h\Phi(x_{H,1}),\qquad
        3G_k(X,3D)-r_h[u^{2k-2}]w^{H-1}.
\]
After division by three these are respectively $\Phi(x_{h,3})$
and the coefficient of $u^{2k-2}$ in $S^{\rm sc}_{h,3}$ computed
above. This proves both identities.
\end{proof}
Let $\Phi$ denote the change from complex families to smooth families
of almost-complex threefolds in the Chern-number-zero universal
algebra. We recall the detection results used here
\cite[\S\S5--7]{WueProduct}. In the complex algebra a primitive class
annihilated by every proper division of cycles belongs to the finite
rational span of the $x_{j,1}$. Their images under $\Phi$ are linearly
independent. The latter assertion uses proper pullback to connected
stable maps of increasing arithmetic genus. In genus below $j$ that
pullback is zero for $x_{j,1}$, and in genus $j$ the regular embedding
gives a nonzero class. Neither assertion requires an evaluation by a
virtual fundamental class for arbitrary analytic families.

\begin{proof}[Proof of the universal identity in Theorem~\ref{thm:isolated}]
Let $J_0=x_{h,0}=1$. On the tube provided by
Proposition~\ref{prop:isolation}, the inverse image of the supported
point $d[D]$, $1\le d\le3$, under cycle addition is the union of the
pairs $(a[D],b[D])$ for $a+b=d$. The preimage under multiplication
by $m$ is the divided point if $m\mid d$, and is empty otherwise.
The maps pull back unit cochains, so there are no stabilizer factors.
Thus the $J_d$ and the scalar $x_{h,d}$ have the same coproduct and
division laws through degree three. In particular
\[
 \Delta J_d=\sum_{a+b=d}J_a\otimes J_b,\qquad
 \rho_m J_d=
 \begin{cases}
 J_{d/m}&m\mid d,\\
 0&m\nmid d .
 \end{cases}
\]
For the scalar classes these formulas follow from
\cite[\S4]{Pardon} in Chern number zero.

For $d=1,2$ deformation to the normal bundle preserves the isolated
section $d[D]$. Indeed a sequence of additional cycles would, by normal
dilation and compactness as in \S\ref{sec:neighborhood}, give an
off-section cycle of degree at most two in the linear normal model.
The entire degree-$d$ cycle space of that model is a point, and the
empty tuple represents $x_{h,d}$. Consequently
$\Phi(J_d)=\Phi(x_{h,d})$. At the first degree of a possible difference
the displayed operations make $J_d-x_{h,d}$ primitive and killed by
proper divisions. Detection therefore gives $J_d=x_{h,d}$ successively
in degrees one and two. For $d=3$ subtract also $(r_h/3)x_{H,1}$.
This correction is primitive and killed by all proper divisions.
The difference lies in the same multiplicity-one span, and has zero
smooth image by Proposition~\ref{prop:borel}. This proves
\eqref{eq:main-class}. The numerical identity is the separate
algebraic pairing in that proposition.
\end{proof}


In the compact example, change just the cubic coefficient of the
genus-three scalar factor to $J_3$.
Its logarithm gains $(8/3)x_{10,1}$. The new exponent at $3\alpha$
in genus ten is therefore $c_{3\alpha,10}-(4720/3+8/3)c_{\alpha,3}
=c_{3\alpha,10}-1576c_{\alpha,3}$. Thus the isolated local unit
accounts for the fractional part even though the complete change
from equal to opposite weights also has an integral part.

In genus five the quotient by dilation contains pencils. A single
cubic evaluation no longer suffices; two independent point products
give the least tuple in Theorem~\ref{thm:cubic-evaluations}. Theorem~\ref{thm:pencil}
in Appendix~\ref{app:pencils} proves the comparison in this case,
including the virtual multiplicity of a cuspidal normalization.

\section{Composite degrees and the contributions of multiple images}
\label{sec:composite}
Write $p_{h,d}^\epsilon=[q^d]\log F_h^\epsilon(q)$ for a universal
connected class and $S_{h,d}^\epsilon$ for its numerical connected
residue, $\epsilon=\mathrm{sc},\mathrm{ad}$. The power-locus bound in
Theorem~\ref{thm:main} excludes cycles that are proper integral multiples of another entire cycle. A cycle can still have an individually repeated component. We now retain the double of an entire cycle and determine its contribution.
In the formulas below $j$ may have either sign.

\begin{theorem}\label{thm:doubleimage}
Suppose $m\ge2$, $|j|\le n_m$ and $|j|>n_{m/e}$ for every divisor
$e\ge3$ of $m$. For even $m$ write
$P_{m/2,j}=\sum_h b_h x_{h,1}$, interpreting a weight outside the
Laurent range as zero. For odd $m$ set $b_h=0$. Then there is a
finite polynomial $R$ such that
\[
 P_{m,j}=\sum_h b_hp_{h,2}^{\rm sc}+\iota(R),\qquad
 R=\sum_kN_{k,m,j}u^{2k-2}-\sum_h b_h S_{h,2}^{\rm sc}.
\]
The coefficient of $w^{h-1}$ in $R$ vanishes if
$h<1+m(g-1)+|j|$. In particular this computes every nonconstant
quartic weight in every core genus and the sextic weights $|j|>g-1$.
\end{theorem}

\begin{proof}
On the projective base for the $j$th coefficient the translated
spectral tests avoid $e$th powers for each $e\ge3$ by the
codimension computation of Section~\ref{sec:weights}. Avoidance
persists in a collar of the compact incidence. If $m$ is even, all
proper multiples in degree $m/2$ are similarly excluded, so
Theorem~\ref{thm:main} supplies the $b_h$ when $m/2>1$.
In degree one the only weight is zero. Proper division by two
sends $P_{m,j}$ to $P_{m/2,j}$ and $p_{h,2}^{\rm sc}$ to $x_{h,1}$.
All remaining proper divisions vanish on both sides. The logarithms
are primitive. It follows from multiplicity-one detection that the
remainder has the form $\iota(R)$ for a unique finite Laurent
polynomial with exponents at least $-1$.

Fix a bound $K$ on source genus. In the separation argument of
Theorem~\ref{thm:separate} keep every tuple of distinct components
with normalization genus sum at most $K$, including the root of a
double whose own degree-two covers have minimum genus above $K$.
At each step order by total reduced degree, the genus sum and then
minus the number of distinct images. Lemma~\ref{lem:weighted-specialization}
makes the first remaining stratum compact after the neighborhoods
already treated have been excised. The shear and jet argument makes
its normalized images embedded and pairwise disjoint. A connected
support is therefore a simple embedded curve or a double of one.
The equations on a reduced Hilbert chart are those of its
\emph{weighted} cycle. Its transverse solutions are finite.

Before excising a doubled solution, fix its reduced curve $D'$ and
vary the normal bundle. Here $\det N_{D'}=K_{D'}$, after trivializing
a constant line. At an unramified fiber away from the tests choose
a framed chart, a uniformizer $t$, a linear form $\ell$ separating
all sheets $y_i$, and the monic polynomial
$S(T)=\operatorname{Norm}(T-\ell(y))$ of their union on the chart.
For a constant area bivector choose $v=\iota_\ell\Omega$, so
$\ell(v)=0$. The transition
\[
 y\longmapsto y+v S(\ell(y))
          \sum_{a=1}^{B}\sum_{l=0}^{r-1}c_{al}t^{-a}\ell(y)^l
\]
is an automorphism of the punctured chart with determinant one.
It fixes each sheet. Its first normal derivative on sheet $i$ is
given by $v\otimes\ell$ times $S'(\ell(y_i))$ times the displayed
polynomial. The Vandermonde determinant is a unit modulo $t^B$ and
three generic forms span $\mathfrak{sl}_2$ in this way.
For large $B$ these principal parts map onto
$H^1(D',\operatorname{End}_0 N_{D'})$; using the union of disjoint
sheet values makes the jets independent for several marked curves.
The same framed parameter bundles and base cutoffs as in
Lemma~\ref{lem:shear-jets} realize these jets in one algebraic
target family, with a smooth section over the compact parameter base.

A good normal bundle is accessible even from an unstable one.
For an ample enough line $Q$, the initial bundle and any prescribed
good bundle of this determinant both admit quotients $Q$ (a general
section of the globally generated $N^*Q$ has no zeros). They belong
to one irreducible extension space
$\operatorname{Ext}^1(Q,K_{D'}Q^{-1})$. The open set of stable
two-rigid bundles there is dense by \cite[Theorem~2.2]{BP06}.
The preceding Kodaira--Spencer map is onto and deformations of
fixed-determinant bundles on a curve have no second obstructions.
Thus the parameter family is smooth over the bundle stack locally
and can realize an arbitrarily small change into this open.
The regular cut and its sign persist. Each root then has acyclic
normal bundle and extends uniquely on a relative Hilbert chart.
Any other nearby cycle of multiplicity at most two, recentered
and dilated to unit normal height, would limit to a compact
off-section cycle in the normal bundle. The disk mass bound follows
from the fixed projected class and a common closed taming form.
The normalization of a positive component would supply a nonzero
normal section on a cover of degree at most two. Two-rigidity excludes
this case. This proves isolation on the full cycle space, uniformly
on a smaller parameter neighborhood. The separation induction can
now proceed with these compact neighborhoods excised.

Over the isolated chart the doubled point is the scalar class with
empty evaluation tuple. Its connected pullback tests give those of
$p_{h,2}^{\rm sc}$: the product in the logarithmic subtraction
vanishes on connected domains. The local intersection degree for the
divided insertion is the same sign, since doubling maps the reduced
section homeomorphically to the doubled section. Testing the divided
insertion successively in genera through $K$ identifies its normalized
signed totals with $b_h$. The same triangular test on the residual
simple intersections identifies their totals with $[w^{h-1}]R$
through $K$.

For the virtual pairing the difference of tangent complexes at a
map $f$ to the moving reduced curve is
$\operatorname{Cone}(q^*Rp_*N\to R\pi_*f^*N)$, with the notation of
Proposition~\ref{prop:moving-embedded-curve}. The first term vanishes
by acyclicity. Two-rigidity and cohomology and base change make the
second a bundle of rank $2(k-2h+1)$ in degree one on the whole
proper stack of degree-two maps, including its nodal boundary.
Its top Chern class pushes to the scalar degree-two connected residue
times the base class: on a single fiber the equality follows by
equivariant localization, where only the top Chern number survives
in virtual dimension zero. Simple intersections contribute
$w^{h-1}$ by the same proposition. The fixed-genus pairings thus
give the asserted equality for arbitrary $K$. A nonzero double
coefficient has $h\ge1+(m/2)(g-1)+|j|$; its covers start at
$1+m(g-1)+2|j|$. Together with the spectral genus bound for the
original cut, this proves the support assertion.
\end{proof}

\subsection{The exceptional sextic coefficient}
For a genus-two core $n_6=15$. At $j=1$ the compact base is
$\PP^{16}\times\PP^{14}$. A cube of a trace-free quadratic
$T^2+b$ has coefficients in orders three through six equal to
$0,a_2^2/3,0,a_2^3/27$ where $a_2=3b$. The independent sections
enforcing these four conditions on the second factor have ranks
$2,3,4,5$, total $14$, and positive top Chern number. They cut
transversely for general translations and avoid the further condition
$a_2=0$. For general $M$ the sole divisor of $M^2$ has two distinct
points, so each reduced quadratic is smooth of genus four. Properness
of multiplication of cycles and openness near a smooth embedded
cycle show that throughout a small collar every pure triple in the
cut has such a root. Pure sixth powers are absent.

Repeat the preceding proof, retaining doubles and triples separately.
At a tripled root, first vary the normal bundle to the general stable
locus on the genus-four curve. Theorem~\ref{thm:isolated} shows that
288 cyclic families remain in the normal bundle. At a fresh unramified
fiber the same interpolation with $S(\ell(y))^2$ fixes the first
normal jets while prescribing quadratic terms on every sheet. In
particular it realizes the determinant-one shears of
Proposition~\ref{prop:isolation} simultaneously at the finitely many
curves. A general quadratic change now isolates each tripled root
while preserving the reduced transverse cut. These changes precede
excision. The full relative germ near the multiple is its moving
section by the uniform isolation in that proposition. Thus no
remaining connected support through the chosen bound is lost.

The two proper divisions give
\[
 \rho_2P_{6,1}=30x_{5,1}-5x_{6,1}+\tfrac49x_{7,1},\qquad
 \rho_3P_{6,1}=\tfrac12 x_{4,1}.
\]
Indeed, subtract $P_{3,3}=x_{7,1}/18$ from one quarter of the
genus-two cubic row, and apply the quadratic row in
\eqref{eq:arith-two}. Let $b_5=30,b_6=-5,b_7=4/9$ and put the other
$b_h$ equal to zero. The two divided tests measure the root signs
separately without automorphism factors. Fix $K\ge13$ and apply the
connected pullback tests up to genus $K$. The undivided test at
each tripled root uses $J_3=x_{4,3}+96x_{13,1}$, and its
fixed-target test has the same extra $96w^{12}$. Consequently
\[
 P_{6,1}=\tfrac12p_{4,3}^{\rm sc}+\sum_h b_hp_{h,2}^{\rm sc}
           +\iota\!\left( [z]S_{2,6}(z)-\tfrac12 S_{4,3}^{\rm sc}
                                      -\sum_h b_hS_{h,2}^{\rm sc}\right).
\]
Here $[z]$ means Laurent weight one; the remainder in parentheses
includes $48w^{12}$ in addition to the normalized simple-image counts.
This use of the nonlinear germ is essential: tests on a moving
reduced chart need not isolate the cubic amplitude in the linear
bundle, even when they do detect the amplitude at a fixed base point.
Indeed the derivative in the moving base coordinates at the root is
invertible. The implicit function theorem solves the test equations
as the amplitude moves.

\subsection{Conversion to opposite-weight factors}\label{sec:rows}
Define $t_{g,d,h}$ by
$F_g^{\rm sc}(q)=\prod_{d,h}F_h^{\rm ad}(q^d)^{t_{g,d,h}}$.
Let $D_{g,d}=S_{g,d}^{\rm sc}-S_{g,d}^{\rm ad}$. Subtracting the
proper divisors in the logarithm gives
\begin{align*}
 \sum_h t_{g,4,h}w^{h-1}
   &=D_{g,4}-\sum_h t_{g,2,h}S_{h,2}^{\rm ad},\\
 \sum_h t_{2,6,h}w^{h-1}
   &=D_{2,6}-2S_{4,3}^{\rm ad}
              -120S_{5,2}^{\rm ad}+20S_{6,2}^{\rm ad}-2S_{7,2}^{\rm ad}.
\end{align*}
The quartic formula follows since the sum over the doubled images
in Theorem~\ref{thm:doubleimage} is the entire quadratic change.
For the sextic the positive odd weights three and one contribute
$1/18$ and $4/9$, respectively, to the genus-seven double, giving
$4(1/18+4/9)=2$. The other doubles and the triple come from weight
one alone. The remaining proper factors convert scalar lower
logarithms into opposite-weight logarithms by Theorem~\ref{thm:main}.

For a genus-two core the quartic row in genera six through eleven is
\[
     (6860,-1470,1302,-317,40,-2).
\]
Its first positive Laurent coefficient starts at $1715w^5$.
The sextic row in genera eight through twenty-two is
\[
\begin{array}{c|rrrrr}
 h&8&9&10&11&12\\ \hline
 t_{2,6,h}&34606024&61350344&125323022&-36417600&62126672\\[2pt]
 h&13&14&15&16&17\\ \hline
 t_{2,6,h}&-41880636&22703802&-9395788&2975620&-715462\\[2pt]
 h&18&19&20&21&22\\ \hline
 t_{2,6,h}&128364&-16652&1476&-80&2
\end{array}
\]
For reproducibility, in each degree $d$ form the disconnected
partition trace $Z_d$ of \cite[Sections~4--8]{BP} and recover
$S_d=Z_d-d^{-1}\sum_{i=1}^{d-1}iS_iZ_{d-i}$. With $r=e^{iu}$ this
calculation takes place over $\QQ(r)$. Multiply the transposition
matrix by the unit and its iterates, express the other partition
operators in the resulting cyclic basis, and form the handle and
inverse caps as in Appendix~\ref{app:handle}. At $z=-1$ use
the hook sum, where the transposition operator need not be cyclic.
For odd weights $1,3,\ldots,2l-1$ the columns
$z^j+z^{-j}$ are $y=z+z^{-1}$ times polynomials of successive
degrees in $y^2$. Interpolate $S_d(z)-S_d(-z)$ at $z=1,\ldots,l$,
with $l=3$ for the quartic and $l=8$ for the sextic. An additional
pair of weights verifies each row. Substitution $w=2-r-r^{-1}$
gives the displayed polynomials. For instance the sole nonintegral
coefficient of $D_{2,6}$ is $376348226/3$ at $w^9$. The cubic
subtraction there is $563480/3$ and the sum of the double terms
to be subtracted is $-61440$, giving $125323022$.

In a compact class $m\alpha$ with $\alpha$ primitive in the free
quotient, substitution of the factors gives
\[
 c_{m\alpha,h}-e_{m\alpha,h}
 =\sum_{d\mid m,\ d>1}\ \sum_{d\gamma=m\alpha}\sum_g
       e_{\gamma,g}t_{g,d,h}
\]
over effective integral roots. For example, at $(m,h)=(6,8)$ one gets
\[
192\sum_{2\delta=6\alpha}c_{\delta,4}
 +14448\sum_{3\delta=6\alpha}c_{\delta,3}
 +34606024\sum_{6\gamma=6\alpha}c_{\gamma,2}.
\]
Indeed the further balanced double of a cubic change starts at
$w^8$ and the further triple of a quadratic change starts at $w^9$.
A genus at least three at a sixth root contributes no term of degree
seven in $w$, by Riemann--Hurwitz. For $m=4$ the corresponding identity at output genus six is
$24\sum_{2\delta=4\alpha}c_{\delta,3}
+6860\sum_{4\gamma=4\alpha}c_{\gamma,2}$. Here
$J_2=24w^5-2w^6$ and the further double of a genus-four correction
starts at $w^6$. No formula here identifies the inner modes on
arbitrary higher power loci.

\appendix
\section{A compact hypersurface containing a Hirzebruch surface}\label{sec:geometry}

We specify the compact example by its fan. Let $N=\ZZ^4$, and let
$v_0,\ldots,v_{10}$ be the rows of
\begin{equation}\label{eq:rays}
 \begin{pmatrix}
 1&0&0&0\\0&1&0&0\\-1&-1&0&0\\0&0&1&0\\0&0&0&1\\
 -4&-2&-1&0\\-4&-2&0&-1\\3&1&1&1\\2&1&1&1\\
 -2&-1&0&0\\-1&0&0&0
 \end{pmatrix}.
\end{equation}
The maximal cones of $\Sigma$ are generated by the following
four-tuples of rays:
\begin{equation}\label{eq:cones}
\begin{gathered}
 (0,1,3,6),\ (0,1,3,7),\ (0,1,4,5),\ (0,1,4,7),\ (0,1,5,6),\\
 (0,2,4,5),\ (0,2,4,7),\ (0,2,5,6),\ (0,2,6,7),\ (0,3,6,7),\\
 (1,3,6,10),\ (1,3,7,8),\ (1,3,8,10),\ (1,4,5,10),\\
 (1,4,7,8),\ (1,4,8,10),\ (1,5,6,10),\ (2,4,5,9),\\
 (2,4,7,9),\ (2,5,6,9),\ (2,6,7,9),\ (3,6,7,9),\\
 (3,6,9,10),\ (3,7,8,9),\ (3,8,9,10),\ (4,5,9,10),\\
 (4,7,8,9),\ (4,8,9,10),\ (5,6,9,10).
\end{gathered}
\end{equation}
Write $V_\Sigma$ for the associated toric fourfold and $D_i$ for
its invariant divisors. Let $Y$ be a general anticanonical
hypersurface, nondegenerate with respect to its Newton polytope,
with nonzero vertex coefficients. This is the resolution chamber used in \cite[Section~4]{WueCounter}.

\subsection{The smooth projective hypersurface}
The polytope $P=\conv(v_0,\ldots,v_8)$ is reflexive, and
\eqref{eq:cones} is a subdivision of its boundary fan. Each maximal
cone has determinant $\pm1$; the subdivision is complete.
For an explicit strictly convex support function take divisor
coefficients
\begin{equation}\label{eq:ample}
 A_i=1009(10000+h_i)+(i+1),\qquad
 h=(10,11,3,9,2,7,1,4,8,-5,-3).
\end{equation}
For each maximal cone $\sigma$, solve
$\langle m_\sigma,v_i\rangle=-A_i$ for $i\in\sigma$.
The off-cone values
$\langle m_\sigma,v_i\rangle+A_i$ are all positive, with minimum
$5047$. This proves projectivity. Repeating the same calculation
with all coefficients equal to one gives integral slopes and
nonnegative values. Thus the anticanonical divisor is basepoint-free
and the fan is a crepant refinement of the reflexive fan. Bertini
and adjunction give smoothness and $K_Y\simeq\mathcal O_Y$.
The anticanonical exact sequence, toric structure-sheaf vanishing,
and Serre duality give $H^1(\mathcal O_Y)=H^2(\mathcal O_Y)=0$.

The four vertices $v_0,v_1,v_3,v_4$ of $P$ form the standard
basis of $N$. In particular, the lattice $N''$ generated by the
lattice points on vertices and edges of $P$ is all of $N$.
The Brauer-group formula of \cite[Corollary 3.9]{BK} therefore gives
\[
 \operatorname{Br}(Y)
 \simeq\operatorname{Hom}\bigl(\Lambda^2N/(N\wedge N''),\QQ/\ZZ\bigr)=0.
\]
For this Calabi--Yau threefold the Brauer group is the torsion in
$H^3(Y;\ZZ)$. The universal coefficient theorem then implies
\begin{equation}\label{eq:torsion}
 \operatorname{Tors}H_2(Y;\ZZ)=0.
\end{equation}
This uses the full lattice, rather than a rational cohomology rank.

\subsection{Nef divisors and the contracted locus}
The divisor $D_{10}$ misses $Y$: its dual Newton face is a vertex
whose coefficient is nonzero. Define integral divisors $H_1,\ldots,H_6$
on $Y$ by the columns of
\begin{equation}\label{eq:nef-matrix}
 \begin{pmatrix}
 0&0&1&1&2&2\\
 1&1&2&4&5&5\\
 0&1&2&4&4&5\\
 0&-1&-1&-3&-3&-3\\
 0&-1&-1&-2&-2&-2\\
 0&0&1&2&2&2
 \end{pmatrix}
\end{equation}
in the ordered divisor list $(D_2,D_5,D_6,D_7,D_8,D_9)|_Y$.
Each has a nef lift to $V_\Sigma$: use the displayed coefficients,
zero on the other rays except $10$, and add coefficient zero on
$D_{10}$ for $H_1,H_2,H_3$, or coefficient one for $H_4,H_5,H_6$.
Solving the support-function equalities as above gives nonnegative
values of the support-function inequalities for every cone and every ray.
Since $D_{10}|_Y=0$, all six
$H_i$ are nef on $Y$. Toric linear equivalence identifies $H_4$
with $D_0|_Y$.

The support function of $D_0$ defines a semiample contraction.
Its positive-dimensional orbit fibers have the following minimal
cones in their closures:
\begin{equation}\label{eq:contracted-cones}
 (8),\quad(9),\quad(10),\quad(2,4),\quad(6,7).
\end{equation}
Here is a direct finite way to obtain the list. Form
$P_0=\{m:\langle m,v_i\rangle+\delta_{i0}\ge0\}$ and, for
each cone $\sigma$, impose its equalities to get the face
$F_\sigma\subset P_0$. The orbit fiber is positive-dimensional
exactly when $4-\dim F_\sigma>\dim\sigma$; taking minimal cones
in this set gives \eqref{eq:contracted-cones}.
Consequently every curve of $H_4$-degree zero on $Y$ lies in
\begin{equation}\label{eq:contracted-locus}
 E\cup F\cup C_1\cup C_2,
 \qquad E=D_9|_Y,\quad F=D_8|_Y,
\end{equation}
where $C_1,C_2$ are the intersections with the orbit closures
of $67,24$. Their restricted Newton polytopes are the primitive
segments joining $(0,1,-1,-1)$ to $(-1,2,-1,1)$ and
$(-1,2,1,-1)$, respectively. Each binomial cuts the dense torus
in a translate of $\CC^*$. Basepoint-freeness excludes boundary
components for a general equation, and nondegeneracy makes the
closures smooth; hence both curves are $\PP^1$. They miss
both surfaces: there are no cones containing $678$ or $248$,
and the Newton faces at $679,249$ are singletons. There is no
cone containing $2467$, so the two curves are disjoint.

\begin{lemma}\label{lem:surfaces}
The surfaces are $E\simeq\Bl_{p_1,p_2}\PP^2$ and $F\simeq\FF_1$.
Their intersection is the negative section $s$ on $F$, and the
strict transform of the line through $p_1,p_2$ on $E$.
\end{lemma}
\begin{proof}
The anticanonical face at ray $9$ is the primitive edge from
$(0,1,-1,-1)$ to $(1,-1,-1,-1)$. Its binomial equation fixes a
nonzero torus parameter. Rays $2,7$ are excluded on $E$ because
their intersections have singleton Newton faces; ray $10$ is absent.
The remaining rays $3,4,5,6,8$ project to
\[
 (1,0),\quad(0,1),\quad(-1,0),\quad(0,-1),\quad(1,1),
\]
with cones $38,84,45,56,63$. This is the fan of the plane blown
up at two distinct torus fixed points.

For $F$, the star of ray $8$ away from $D_{10}$ has cones
$137,147,379,479$. Quotient by the fiber ray $7$ gives the base
rays $(1,0),(0,1),(-1,-1),(-1,0)$, the fan of $\FF_1$.
The Newton vertex $(1,-1,-1,-1)$ has exponents
$(0,0,0,1,0,0)$ on rays $1,3,4,7,9,10$; all other vertices
have fiber exponent zero. The equation is monic linear in that
fiber coordinate and its zero locus is a graph over $\FF_1$.
The common divisor in these two surface fans is the stated curve.
\end{proof}

Let $e_1,e_2$ denote the exceptional curves on $E$, and $f$ a
ruling fiber of $F$. Restriction of \eqref{eq:nef-matrix} to the
surface fans and to the two orbit curves gives
\begin{equation}\label{eq:degrees}
\begin{array}{c|rrrrrr}
 &H_1&H_2&H_3&H_4&H_5&H_6\\\hline
 s&0&1&0&0&0&0\\
 f&0&0&1&0&0&0\\
 e_1&1&0&0&0&1&1\\
 e_2&0&0&0&0&0&1\\
 C_1&0&0&0&0&1&0\\
 C_2&1&0&0&0&0&0.
\end{array}
\end{equation}
These pairings can be computed with the usual toric divisor
relations: the linear relations are
$\sum_i\langle m,v_i\rangle D_i=0$, and a divisor product vanishes
when its rays span no cone. The monic and binomial equations above
identify the actual surfaces on which the restrictions are taken.

\begin{proposition}\label{prop:support}
For every positive integer $a$, the class $\alpha=as+(a+1)f$ is primitive in $H_2(Y;\ZZ)$.
Every effective cycle of class $m\alpha$, $m\ge1$, is supported on
$F$. The same holds for every effective additive summand of $\alpha$.
\end{proposition}
\begin{proof}
By \eqref{eq:degrees}, $(H_3-H_2)\cdot\alpha=(a+1)-a=1$,
which proves primitivity. A cycle in class $m\alpha$ has zero
$H_1,H_4,H_5,H_6$ degrees. Nefness forces the same for each
component. The $H_4$ contraction confines it to
\eqref{eq:contracted-locus}; $H_1,H_5$ exclude $C_1,C_2$.
The effective cone of $E$ is generated by $e_1,e_2,s$.
The $H_1,H_6$ degrees exclude any $e_1$ or $e_2$ coefficient,
so a remaining component in $E$ is supported on $s=E\cap F$.
All components therefore lie in $F$.

Every effective summand of $\alpha$ has the same four vanishing
degrees, so the argument applies to it as well. On $F$, the
$H_2,H_3$ pairings recover the coefficients of $s,f$. Hence a
decomposition of $\alpha$ into effective classes corresponds to
an actual effective decomposition of $as+(a+1)f$ on $F$.
\end{proof}

\section{The scalar handle in the partition basis}\label{app:handle}
We use the multiplication operator of
\cite[Appendix, equations (29) and (31)]{BP}, the cap of its Lemma 6.3,
and the starred normalization of its Section 6.4.1.
Put $r=e^{iu}$. Initially keep general nonzero weights $t_1,t_2$, and put
\[
 k=t_1t_2,\qquad
 a=-(t_1+t_2)\frac{r-1}{r+1},\qquad
 b=-3(t_1+t_2)\frac{r^2-1}{r^2+r+1}.
\]
In the partition basis $1=e_{111}$, $x=e_{21}$, $y=e_3$,
multiplication by $x$ is
\begin{equation}\label{eq:M}
 M=\begin{pmatrix}0&-3k&0\\1&a&-2k\\0&3&b\end{pmatrix}.
\end{equation}
Since $M1=x$ and $M^21=-3k+ax+3y$, the unit is cyclic. Consequently
\begin{equation}\label{eq:algebra}
 \begin{gathered}
 L_1=I,\qquad L_x=M,\qquad L_y=\frac{M^2-aM+3kI}{3},\\
 \mathcal A=\QQ(t_1,t_2,r)[x]/\bigl((x-b)(x^2-ax+3k)+6kx\bigr).
 \end{gathered}
\end{equation}
Indeed $My=-2kx+by$ gives the stated relation. Multiplication also gives
\[
 y^2=2k^2-\frac{2kb}{3}x+\frac{b^2-ab-3k}{3}y.
\]

The starred Frobenius pairing in this basis is
\begin{equation}\label{eq:metric}
 D=\operatorname{diag}\left(-\frac1{6k^3},\frac1{2k^2},-\frac1{3k}\right).
\end{equation}
In general its partition entry is
$1/(\mathfrak z(\lambda)(-k)^{\ell(\lambda)})$, where
$\mathfrak z(\lambda)=\prod_j j^{m_j}m_j!$ and $m_j$ is the
multiplicity of $j$ in $\lambda$. The Appendix Fock pairing differs
by $(-1)^d$ in degree $d$; its multiplication-operator convention
has the same sign. Thus \eqref{eq:M} is the multiplication matrix
in the starred normalization,
and the handle operator is
\begin{equation}\label{eq:handle}
 G=-6k^3I+2k^2M^2-3kL_y^2.
\end{equation}
This is multiplication by the Euler element, the sum of each basis
element times its dual. In the quotient algebra it equals
\[
 -18k^3+2k^2(a+b)x+(9k^2+kab-kb^2)y.
\]

Raising the level-$(-1,0)$ cap covector with \eqref{eq:metric}
and removing its half-integral phase gives the vector
\begin{equation}\label{eq:cap}
 \widetilde\eta_\lambda
 =(-1)^d t_1^{\ell(\lambda)}
   \prod_{j=1}^{\ell(\lambda)}(1-r^{-\lambda_j})^{-1},
 \qquad d=3.
\end{equation}
Let $\widetilde A=\sum_\lambda\widetilde\eta_\lambda L_\lambda$
and let $\widetilde{\bar A}$ interchange $t_1,t_2$. If $A^*$ is
the unscaled starred level operator, then
$\widetilde A=r^{d/2}A^*$. The balanced-level operator is therefore
\begin{equation}\label{eq:B}
 B=r^3G\widetilde A^{-1}\widetilde{\bar A}^{-1},\qquad
 [q^3]F_g=\tr(B^{g-1})\quad(t_1=t_2=1).
\end{equation}
All operators commute because they act by multiplication.
The closed starred normalization has exponent
$d(2-2g+(g-1)+(g-1))=0$, so no further phase multiplies the
ordinary GW coefficient in \eqref{eq:B}.

\begin{proposition}\label{prop:trace}
At scalar weights the characteristic polynomial of $B$ is
\begin{equation}\label{eq:charpoly}
 \begin{split}
 \det(TI-B)={}&T^3-w^3(w^2+50w+81)T^2\\
 &+\frac43w^6(w-6)(w^2+54w-243)T\\
 &-\frac43w^9(w-4)(w-3)^2(w^2+54w-243).
 \end{split}
\end{equation}
In particular,
\begin{equation}\label{eq:scalar-traces}
 \begin{split}
 [q^3]F_2&=81w^3+50w^4+w^5,\\
 [q^3]F_3&=2673w^6+9612w^7+2534w^8+\frac{292}{3}w^9+w^{10}.
 \end{split}
\end{equation}
\end{proposition}
\begin{proof}
Substitute $t_1=t_2=1$ into \eqref{eq:M}--\eqref{eq:B} and reduce
using $r^2-(2-w)r+1=0$. The three elementary symmetric functions
of $B$ are respectively
\[
 \begin{gathered}
 w^3(w^2+50w+81),\qquad
 \frac43w^6(w-6)(w^2+54w-243),\\
 \frac43w^9(w-4)(w-3)^2(w^2+54w-243).
 \end{gathered}
\]
These give \eqref{eq:charpoly}; the trace and the identity
$\tr(B^2)=(\tr B)^2-2\sigma_2(B)$ give \eqref{eq:scalar-traces}.
The calculation can also be made entirely in \eqref{eq:algebra}:
invert $\widetilde\eta^2$, multiply by $r^3$ times the Euler
element, and take the resultant of its characteristic equation with
the monic cubic defining $\mathcal A$. This uses the same printed
algebraic data and avoids a matrix inverse.
\end{proof}

At antidiagonal weights the same construction has eigenvalues $U,U,V$,
where
\begin{equation}\label{eq:UV}
 U=w^3(4-w)(3-w)^2,\qquad V=w^3(3-w)^2.
\end{equation}
These are the three hook-length contributions of \cite[Corollary~7.3]{BP}. The leading scalar coefficient
$2673=6^4+3^4+6^4$ also agrees with the unramified degree-three
cover count for a genus-three base.

\subsection{A comparison in degree five}
For a genus-two core, the first prime beyond the cubic gives
\begin{align*}
 D_{2,5}(w)=w^6\big(&444088+142232w+340864w^2-156542w^3+75396w^4\\
                    &-23737w^5+5054w^6-696w^7+56w^8-2w^9\big).
\end{align*}
Here the subscripts denote genus and covering degree as in \eqref{eq:compact-prime}. For example, at $5\alpha$ in output genus
seven the scalar exponent is the elementary exponent minus $444088$
times the sum of the genus-two elementary exponents at its effective
fifth roots. Larger core genera are excluded by Theorem~\ref{thm:main}.

For the calculation let $Z_a$ denote the disconnected degree-$a$
residue and $L_a$ the coefficient of its logarithm. Apply
$aL_a=aZ_a-\sum_{b=1}^{a-1}b L_b Z_{a-b}$ at each pair of weights.
For a genus-two core $Z_a$ is the trace of the balanced handle
of \cite[\S\S5--8]{BP}. In degrees $1$ through $5$ the transposition
operator is cyclic on the unit over $\QQ(r)$, $r=\exp(iu)$, at both
weight specializations. Using the powers of that operator on the unit
as basis, express each partition class as a polynomial in the operator.
Multiply the handle and the inverse level operators modulo its
characteristic polynomial and take the trace. At antidiagonal weights
one may equivalently sum $\prod_{b\in\lambda}(2-r^{\ell(b)}-
r^{-\ell(b)})$ over partitions, with $\ell(b)$ the hook length.
Subtract the two values of $L_5$ and substitute $w=2-r-r^{-1}$.
These are rational function identities over $\QQ$.

\section{Individual coefficients on a split genus-three curve}\label{app:split}
Let $L$ be general of degree two on a smooth curve $C$ of genus three,
and let $M=K_C L^{-1}$. For this appendix $g=m=3$ unless otherwise
specified. The split bundle is useful for distinguishing the
individual weight coefficients. The eight lines in
Theorem~\ref{thm:classification} instead occur on a general stable
bundle with the same determinant.

\subsection{A proper incidence space for a weight coefficient}
For arbitrary $g,m$ put $a=k-1-m(g-1)$ and let $P_A=R\pi_*f^*A$ on
the proper stack of connected genus-$k$ maps to $C$ of degree $m$.
Localization gives
\[
 N_{k,m,j}=\int_{[\overline{\mathcal M}_k(C,m)]^{\rm vir}}
           c_{a-j}(-P_L)c_{a+j}(-P_M).
\]
For $1\le j\le a$ represent $P_M$ as $[U\to V]$. A global
presentation is obtained by the dual evaluation into a sufficiently
positive relatively generated line on the universal curve tensored
with its relative section-space dual: tensor by $f^*M$, take the
quotient and push forward. Source flatness, relative Serre vanishing
and the sequence give two locally free terms. On the projectivization
of lines $\PP(U)$ the tautological map to $V$ cuts out the proper
stack $Y$ of sections up to nonzero scalar. Let $\alpha_Y$ be its
refined zero-section class on the pullback of the stable-map virtual
class. With $\xi=c_1(\cO(1))$ the projective-bundle formula gives
\[
 N_{k,m,j}=\int_{\alpha_Y}\xi^{j-1}c_{a-j}(-P_L).
\]
Indeed the pushforward of $\xi^{j-1}c_{\mathrm{rk}V}(V(1))$
is $c_{a+j}(V-U)$, since $\mathrm{rk}V-\mathrm{rk}U=a$.

At a genus-three cubic the first possible source genus of mode
$j>0$ is $7+j$ by Theorem~\ref{thm:main}. Here the fixed-genus
integral is also the leading multiplicity-one coefficient. For
$j=1,\ldots,6$ these values are
\[
 3612,\quad 1228,\quad 506/3,\quad 21/2,\quad 8/27,\quad 1/324 .
\]
To recover them directly, at $z=1,\ldots,7$ evaluate
$\tr B_3(z)^2-w^2\tr B_2(z)^2+w^6/3$ in the matrices of
Appendix~\ref{app:handle} and their degree-two analogues. A symmetric
Laurent polynomial through exponent six is a polynomial of degree at
most six in $z+z^{-1}$. Interpolation at these seven distinct values,
followed by extraction at $w^{6+j}z^j$, gives the list. Alternatively
interpolate the odd and even parts separately at pairs $\pm z$.
On one normal axis the six coefficient weights have product
$2^2 3^4=324$, so for $j=3,4,5,6$ the corresponding degrees before
this division are $54648,3402,96,1$. Mode five has no contribution
through genus eleven.

In genus eight the degree $3612$ is integral on an automorphism-free
space of smooth ramified covers. For a general $M$ no connected
\'etale double or triple admits a section of its pullback. Indeed
there are finitely many such covers, and their trace-free bundles
are semistable of degree zero and rank at most two, being trivialized
by a finite \'etale cover. After a general degree-two twist there
are no sections by \cite[Proposition~1.6.2]{Raynaud}. A nonzero
height on a connected stable map must be nonzero on a nonconstant
component. A degree-three component then has normalization genus
at least eight, while a double plus a degree-one component has genus
sum at least nine. Thus at genus eight the source is smooth.
A cyclic triple with total ramification two would have one branch
value, contrary to the abelian monodromy relation. The transitive
monodromy is $S_3$ and its centralizer is trivial. The deformation
complex of a finite smooth-source map to $C$ is its torsion
ramification cokernel, so the Hurwitz chart is smooth. The degree
comes from an integral refined Euler class and permits excess.

At genus ten it matters whether one extracts $u^{18}$ or $w^9$.
For positive weights one and three the change of variable gives the
following decomposition. Each entry is a contribution to $[w^9]$.
\[
\begin{array}{c|rrr}
 \text{source genus}&j=1&j=3&4\text{ times the sum}\\ \hline
 8&3612(161/720)&0&48461/15\\
 9&-2391(2/3)&0&-6376\\
 10&60659/60&506/3&23593/5\\ \hline
 \text{total}&674/3&506/3&4720/3
\end{array}
\]
Indeed $u=2\arcsin(\sqrt{w}/2)$ gives the indicated source multipliers;
the raw numbers are obtained from the same handle polynomial by
$w=2-2\cos u$. Opposite weights exchange under inversion, explaining
the factor four. The genus-ten row is integral at three after summing
the modes. The pole at three in the final column enters through
the genus-eight term.

The cancellation of denominators at fixed genus ten is also visible
in the restrictions of the two integrands, although it does not imply
descent of their sum to the coarse incidence space. Consider an
order-three automorphism of a point of $Y$, and put $t=c_1(\chi)$ for
a generating character. If it acts nontrivially on the section line,
then in $A^2(B\mu_3)=(\ZZ/3)t^2$ one has
\[
 c_2(-P_L)=-t^2,\qquad \xi^2=t^2.
\]
Indeed invariance of the fiber multiset under cubic dilation forces a
cyclic spectral triple, with normalization of genus nine or ten (as
classified below). In genus ten the two summands of
$-R\Gamma(f^*L)$ of nontrivial character give $\chi+2\chi^2$ up
to inversion: the covering eigenlines have degrees $-1,-2$.
In genus nine their degrees are both $-1$, since for local exponents
one and two the square eigensheaf is normalized by division by
a local parameter at the second point. Their contributions are then
$\chi+\chi^2$, together with one additional character from the nodal
or contracted part. In fact for the normalization of all positive
components the correction to $-R\Gamma$ is the sum of the sections
of the finite sheaves $(f_{\rm norm,*}\cO)/f_*\cO$ and
$R^1 f_*\cO$, tensored with $L$. In either case the pairwise products
of the character exponents sum to $-1$ modulo three.

When the automorphism fixes the section its $\xi^2$ restricts to
zero, and it fixes the positive normalizations pointwise: a nonzero
section parametrizes its spectral image birationally in prime degree,
and the only possible additional positive component has degree one.
For $c$ positive components of genera $h_i$ the finite-sheaf
correction has dimension $10-\sum h_i+c-1$, with at least $c-1$
trivial characters coming from independent constants. A nonzero second
Chern restriction thus requires $\sum h_i\le8$. The preceding
genus bounds leave a single genus-eight component. Such a point
exists since its virtual degree is 3612. Attach to it the contracted
genus-two tail $y^2=x^6-1$ at $x=0$; the action $x\mapsto\zeta x$
fixes the attachment and has characters $\chi,\chi^2$ on the two
differentials $dx/y,x\,dx/y$. It gives
$c_2(-P_L)=2t^2\ne0$ and $\xi^2=0$. This excludes descent of
$c_2(-P_L)+\xi^2$ as a three-local operational Chow class. The
restriction to a stabilizer does not by itself compute a stratum's
virtual degree.

\subsection{The support and multiplicity of the third mode}
In genus ten two independent quadratic coefficient tests on the
projective section have class $4\xi^2$. Their common zero on $Y$
has spectral polynomial $y^3=b$ with $0\ne b\in H^0(M^3)$, up to
scale. Indeed the trace vanishes and there are two independent
coordinates on $H^0(M^2)$. A meromorphic cube root of $b$ would be
holomorphic, contrary to $H^0(M)=0$. The sole nonconstant component
is the normalization of the integral spectral cubic. Its genus is
$7+r$ with $r$ the number of distinct zeros whose multiplicities
are nonzero modulo three. Thus $r\le3$; the case $r=2$ gives an
actual genus-ten map on attaching a contracted elliptic tail.

Put $A=M^3$. For every $p$, $h^0(A(-3p))=1$: a higher dimension
requires $A=K_C(3p-q)$, a locus of dimension at most two. Two
triple zeros would require $A\in3W_2$, also excluded. The sections
with a triple zero therefore form a curve in $|A|$ with normalization
$C$, and the determinant of the order-two jet bundle gives its
degree $3\deg A+3\deg K_C=30$. On this curve a genus drop occurs
on divisors $3p+2q+s$. The weighted Abel map has degree
$3!(3\cdot2)^2=216$: the tangent columns scale holomorphic volume
by $3\cdot2$, hence real volume by its square, relative to the
ordinary degree-six Abel map. General $A$ avoids branch and collision
loci. The other possible divisors consist of three double zeros,
one for each of the 64 square roots $B$ of $A$, all with a unique
reduced effective divisor. Their cubic images have three cusps and
smooth normalizations of genus ten.

Each isolated three-cusp polynomial contributes $2/3$ to the integral
of $\xi^2$. To see the virtual multiplicity, take $B^2=M^3$ and a
section $s$ with its three simple zeros $p_i$. The smooth normalization
has $t^3=s$ in $B/M$ and $y=t^2$. Fix the amplitude of $b_0=s^2$
by a linear form $\lambda(b_0)=1$. This removes the cyclic stabilizer
on a framing cover of degree three. Vary both quadratic tests by
$q=\epsilon^2Q$, with $Q\in H^0(M^2)$ nonzero at the $p_i$;
$M^2$ is generated there for a general $M$. In coordinates at these
points the equations are $F_i=y^3+\epsilon^2Q_i(x)y-x^2-\beta_i(x)$
where $b=b_0+\beta$ and $\beta\in\ker\lambda$. Evaluation
of $\ker\lambda$ at the three points is invertible. Thus $F_i=F_{i,x}=0$
determine $\beta$ and the three $x_i$ as functions of the three
$y_i$. The residual equations after $y_i=\epsilon z_i$ are
$3z_i^2+Q_i(x_i)=0$. There are eight simple solutions for small
$\epsilon$, and no others since the $z_i$ are bounded. Each singular
point is nodal by its Hessian. The parametric Morse lemma and cubic
Weierstrass preparation bound the local delta invariant by one.
Total delta three is necessary, so all three nodes must occur, and
each resulting normalization is the unique lift of its curve.

At these immersed maps an infinitesimal solution fixing the quadratic
tests has $\dot\beta$ zero at all three nodes by differentiating the
image equation, hence identically zero. The normal sheaf of an
immersion is a line, so a normal section vanishing on the smooth image
is zero. These are eight reduced intersections on the framing cover.
Near their limit the Hurwitz base is smooth of dimension six,
and in $\PP(U)$ the rank of the defining section plus the two
quadratic tests equals the ambient dimension. Conservation of the
isolated refined intersection degree gives length eight also at
the limit. Divide first by three for the framing, then by four for
the class of the tests. This gives $2/3$. Consequently the moving
degree-thirty family together with its boundary contributes
$506/3-64(2/3)=126$. This computation leaves the separate degrees
over the smooth part and the 216 limiting values undetermined.

\subsection{Axis support after the lower coefficient tests}\label{sec:matching}
Consider the lower elementary symmetric coefficients of each projection
in a balanced split bundle. Impose vanishing of all orders less than
$m$ on both axes. The pure-axis families are the vector spaces of
sections of $L^m$ and $M^m$, parametrizing the monic spectral divisors.
For general $L$, the following table gives exactly the degrees in
which these two spaces, with their origins identified, exhaust the cut:
\[
\begin{array}{c|l}
 g&\text{covering degrees with only axis cycles}\\ \hline
 2&1,2,3,4,5,6,7,9\\
 3&1,2,3,4,5,6\\
 4,5&1,2,3\\
 g\ge6&1,2 .
\end{array}
\]
Here generality is for each fixed $m$. The subsequent constructions
of mixed cycles work for every $L$ in the asserted existence ranges.
This criterion concerns the separate projected equations; a fixed
finite subgroup action specifies how the two root sets must be matched.
Put $K=K_C$ and $M=K/L$ in the proof.
\subsubsection{Reduction to a connected root cover}

Away from the zero divisors of $b$ and $c$ each projection has $m$ distinct simple values. A fiber of a mixed cycle therefore defines a bijection of the two root torsors, compatible with continuation along the punctured curve. Choose base roots and label the torsors by $\mathbb Z/m$. Each monodromy acts by a translation on either set. The two translation representations have the same kernel, because a permutation of either set determines the permutation of the other via the bijection. Thus the common image is cyclic of order $d\mid m$. Its orbits on each set have size $d$. A single orbit of the matching gives a connected component over the punctured curve, and its closure gives an irreducible degree-$d$ cycle. Both of its projected $d$-th powers are single valued, extending holomorphically over the missing points by boundedness. The two projected lower coefficients vanish on that degree-$d$ cycle. On this orbit choose the generator acting on the first coordinate by a primitive $d$-th root $\zeta$. It acts on the other by $\zeta^k$ for a unit $k\bmod d$. Equivalently
\[
 v=\phi u^k,\qquad \phi\in H^0_{\rm mer}(C,M/L^k)
\]
on this component. Here $k$ is any positive representative of the slope. The relation follows by invariance of $v/u^k$ under the cyclic action on the function field. The case $d=1$ requires a nonzero global section of $L$, excluded for general $L$.

Conversely every mixed cut cycle of degree $d$ gives one of degree $m$ whenever $d\mid m$. Take the sum of its rotations under $(u,v)\mapsto(\zeta_m^j u,\zeta_m^j v)$ for $0\le j<m/d$. At a general fiber the translated sets partition the roots of $u^m=b_d^{m/d}$ and of $v^m=c_d^{m/d}$. This argument needs only the two projected polynomials. In prime degree, disconnectedness of a mixed cycle would force $d=1$, which explains the irreducibility assertion.

For genus three the slopes $1$ and $-1$ are excluded for generic $L$. For slope $1$ the curve lies in a saturated line subbundle $J\subset L\oplus M$ mapping nontrivially to both factors. The nonzero invariant $d$-th power of its fiber coordinate gives a section of $J^d$. A positive $\deg J$ would force $L=M$ or $L/M\in C-C$, both excluded for general $L$. Hence $J$ is torsion of order dividing $d$ and $L$ belongs to one of finitely many translates of $W_2$, also excluded. For slope $-1$ the invariant $uv$ is a nonzero canonical differential $t$. Thus $\operatorname{div}(b_d)$ is a degree-$2d$ effective subdivisor of $d\operatorname{div}(t)$. Its possible line bundles range over a closed set of dimension at most two, by the finite-subdivisor incidence over $\mathbb P H^0(K)$, whereas $L\mapsto L^d$ is finite surjective onto a Picard variety of dimension three.  

All unit slopes for $d\mid m$, with $m=1,2,3,4,6$, are covered by these cases. When either $b$ or $c$ vanishes the cycle lies on the corresponding axis and is determined by the other monic polynomial, including its multiplicities.

\subsubsection{Exclusion of the remaining slopes at degree five}

Interchanging the two coordinates replaces $k$ by its inverse, so the only additional slope to consider is $k=2$. Write $P=(\phi)_\infty$ and $Q=(\phi)_0$ with disjoint support. Then $\deg P-\deg Q=2$ and $\mathcal O(P-Q)=L^2/M$. Since the latter line of degree two is general as $L$ varies, $\deg P\ge3$. Holomorphicity of $v=\phi u^2$ implies
\[
 (b)\ \ge\ \sum_x \left\lceil\tfrac52\operatorname{ord}_x P\right\rceil x .
\]
The left side has degree ten, giving $\deg P\le4$. If $\deg P=4$, equality forces $P=2D$ and $(b)=5D$ with $\deg D=2$. Thus $L$ is a torsion translate of an effective degree-two line, excluded for general $L$.

If $\deg P=3$ then $Q=q$ is a point. If $P$ is reduced, $(b)=3P+x$ for some point $x$. Using $\mathcal O(P)=(L^2/M)(q)$ gives $M^3/L=\mathcal O(3q+x)$. The target line of degree four is general because $L\mapsto K^3/L^4$ is finite surjective. The locus of divisors $3q+x$ has dimension at most two, so this is impossible for general $L$.

Otherwise $P=2p+r$, allowing $p=r$. The inequality gives $(b)=5p+3r+S$ with an effective divisor $S=s_1+s_2$. This also holds for $p=r$, where the minimum required vanishing is eight. Eliminating $L$ from the two identities
\[
 L^3/K=\mathcal O(2p+r-q),\qquad L^5=\mathcal O(5p+3r+s_1+s_2)
\]
gives
\[
 K^5=\mathcal O(5p+4r+5q+3s_1+3s_2).
\]
The fiber of this weighted Abel map has dimension at most two. Stratify by the coincidences of its five marked points. At $n$ distinct points the differential has rank equal to the rank of their canonical evaluations since all combined weights are positive. Five distinct points impose rank three, since no canonical differential has five zeros. Four distinct points impose rank at least two, since a fiber of the genus-three canonical morphism contains at most two distinct points. Indeed in the hyperelliptic case that map is two-to-one onto a conic and otherwise it is the plane quartic embedding. For $n\le3$ the rank is at least one by base-point freeness. On each stratum the tangent dimension of the fiber is consequently at most two. The points determine $L^5$ and hence finitely many $L$. Properness of the weighted Abel incidence shows that its image in the Picard variety of $L$ lies in a closed subset of dimension at most two. This excludes the remaining slope for genus three.

\subsubsection{Construction of mixed cycles}

Fix a positive integer $k$ coprime to $m$ and a nonzero meromorphic section $\phi$ of $M/L^k$. If a nonzero section $b\in H^0(L^m)$ satisfies
\[
 k(b)\ \ge\ m(\phi)_\infty ,
\]
the graph $v=\phi u^k$ on the normalization of the spectral cycle $u^m=b$ gives a compact mixed cycle in $\operatorname{Tot}(L\oplus M)$. Indeed the coordinate $v$ extends holomorphically at every point of the normalization by the valuation inequality. Projection onto the $u$ coordinate identifies the generic image cycle with the original spectral cycle. Since $k$ is a unit, over a general fiber the $m$ values of $v$ are exactly the distinct roots of $v^m=\phi^m b^k$. This proves all the cut equations even when the normalization is disconnected.

For odd $m\ge9$, take $k=2$, a point $q\in C$, and $P\in|(L^2/M)(q)|$. The degree of this last line is three, so such a divisor exists. The meromorphic section with divisor $q-P$ has pole divisor bounded by $P$. Let $a=(m+1)/2$. Since $\deg L^m(-aP)=(m-3)/2\ge3$, Riemann–Roch gives a nonzero $b$ in this space. It satisfies the required inequality.

For $m=7$ use the same construction with $a=4$. The line in which to find $b$ after imposing $4P$ is
\[
 L^7(-4P)=(M^4/L)(-4q).
\]
The degree-six bundle $M^4/L$ has four independent sections. Its complete linear series has a point $q$ with a section vanishing to order at least four. For instance the Wronskian is a nonzero section of $(M^4/L)^4\otimes K^6$ in characteristic zero, of positive degree. Its zeros are precisely the points where vanishing orders differ from $0,1,2,3$, forcing some vanishing order at least four. Thus this choice of $q$ supplies $b$. Existence of $P$ still follows from its degree.

For $m=8,10$, use $k=3$, $P\in|L^3/M|$ and $\operatorname{div}(\phi)=-P$. The degree of $P$ is four. The line $L^m(-\lceil m/3\rceil P)$ has degree four for both values of $m$ and therefore a nonzero section. For $m=12$, choose an effective degree-four divisor $D$ with $\mathcal O(2D)=L^5/M$: square roots exist on the complex Picard variety and every degree-four line on $C$ is effective. With $k=5$ and $\operatorname{div}(\phi)=-2D$, any nonzero $b\in H^0(L^{12}(-5D))$ satisfies $5(b)\ge24D$, since $25D\ge24D$. The residual degree is four.

Every even integer at least eight either has an odd divisor at least seven or is divisible by $8,10$ or $12$. For the last assertion its odd part, when below seven, is $1,3$ or $5$. The constructions and the degree-multiplication observation finish the theorem. 

\subsubsection{Dependence on the genus}

The genus-three threshold depends on the degree and genus together. For example mixed solutions of degree five already exist for every $L$ if $g\ge4$. With $k=2$ choose $P\in |(L^2/M)(q)|$ as above, now of degree $g$. The line $L^5(-3P)$ identifies with $(M^3/L)(-3q)$ of degree $2g-5$. For $g\ge5$ this is effective by degree. For $g=4$ the degree-six line $M^3/L$ has at least three sections by Riemann–Roch. If it has more than three sections every $q$ works. Otherwise its three-dimensional complete series has positive Wronskian degree, and a ramification point supplies the order-three zero. In either case the valuation condition constructs the cycle. For a general $L$ its normalization cover is connected since five is prime.

\emph{Genus two.} For a genus-two curve and general $L\in\operatorname{Pic}^1 C$, the degrees with axis-only support are exactly $1,2,3,4,5,6,7,9$. In particular the first prime admitting mixed cycles here is eleven.

For slope $1$, a saturated common line of positive degree would force $L=M$. A common degree-zero line carrying a nonzero $d$-th power is a $d$-torsion translate of the trivial line, and would put $L$ in a finite union of translates of the effective degree-one locus. For slope $-1$ the canonical subdivisor argument applies with parameter dimension one. For the remaining slopes at $m=5,7,9$ the connected orbit reduction still applies. Lower-degree divisors of these $m$ have only the excluded slopes. In each case use the inverse slope as well, by interchanging the two normal lines.

For slope $2$ in any odd degree, write $P$ for the pole divisor and $Q$ for the zero divisor of $\phi$. Here $\deg P-\deg Q=1$, and degree one for $P$ would force $L^2/M$ effective, excluded generically. Thus $\deg P\ge2$. The valuation bound $2(b)\ge mP$ forces $\deg P=2$ and equality throughout. Since $m$ is odd, $P=2p$ and $(b)=mp$, excluded up to $m$-torsion. This handles both remaining unit slopes for $m=5$.

At $m=7$ one must also consider slope $3$ and its inverse $5$. Now $\deg P-\deg Q=2$, and $3(b)\ge7P$ bounds $\deg P\le3$. Degree three forces $P=3p$ and $(b)=7p$ by equality and coprimality. At degree two $Q=0$ and $\mathcal O(P)=L^3/M$ is general of degree two. Its only effective representative is generically reduced, since divisors $2p$ vary in dimension one. Thus $(b)\ge3P$, forcing the degree-one line $M^3/L^2=L^7(-3P)$ to be effective, again excluded for general $L$.

At $m=9$ the additional inverse pair comprises slopes $4,7$. For $k=4$ we have $\deg P-\deg Q=3$ and $\deg P\le4$. Degree four forces $P=4p$ and $(b)=9p$. At degree three, $Q=0$ and $\mathcal O(P)=L^4/M$. If $P$ is reduced, $(b)=3P$ so the degree-zero line $M^3/L^3$ must be trivial. If $P=3p$ then $L^4/M$ belongs to a locus of dimension one. If $P=2p+r$ with $p\ne r$, write $(b)=5p+3r+x$. Elimination gives
\[
 K^9=\mathcal O(7p+6r+5x).
\]
The corresponding fiber of the weighted Abel map has dimension at most one. Three distinct points impose canonical evaluation rank two since a genus-two differential has degree two, and two or fewer distinct points impose rank at least one. These points determine $L^9$, so the excluded $L$ again form a subset of dimension at most one. This proves the axis-only assertions.

For the complementary degrees the valuation construction uses that every degree-two line on $C$ is effective. For $m\ge8$ with $3\nmid m$ take $k=3$, $\operatorname{div}(\phi)=-P$ with $P\in|L^3/M|$ of degree two. The residual degree $m-2\lceil m/3\rceil$ is at least two (including $8$ and $10$). For odd $m\ge15$ divisible by three, take $k=4$, $\operatorname{div}(\phi)=-P$, $P\in|L^4/M|$ of degree three. The residual degree $m-3\lceil m/4\rceil$ is at least two: the first cases are $15,21$, and it holds for all $m\ge17$. For $m=12,18$ use $k=5$ and $\operatorname{div}(\phi)=-2D$ where $\mathcal O(2D)=L^5/M$ and $\deg D=2$. Imposing $(b)\ge\lceil2m/5\rceil D$ leaves degree two in both cases. An even multiple of three is divisible by twelve or is twice an odd multiple of three. In the latter case the odd half is $3,9$ or at least $15$. Degree multiplication completes the classification.

\emph{All higher genera.} For $g=4,5$ the only axis-only degrees are $1,2,3$. For every $g\ge6$ they are $1,2$. Each exclusion holds for general $L$, while each existence statement below holds for every $L$.

Write $a=g-1$. To construct a slope-one solution, choose $J\in\operatorname{Pic}^j C$ with $J^m$, $L/J$, and $M/J$ all effective. Nonzero sections of the last two bundles map the line $J$ into $N$. The map is injective on the generic fiber. The image of the compact spectral cycle $w^m=b_J$, for a nonzero $b_J\in H^0(J^m)$, therefore has both projected $m$-th powers nonzero and satisfies the lower coefficient cuts. Simultaneous vanishing of the two line maps at an exceptional point of the base leaves the proper image of the spectral curve well defined.

In $\operatorname{Pic}^j C$ the three effectivity loci are an isogeny pullback of $W_{mj}$ and two translated inversions of $W_{a-j}$. By the Poincaré formula $[W_e]=\theta^{g-e}/(g-e)!$ for $0\le e<g$ and $W_e=\operatorname{Pic}^e C$ for $e\ge g$. Their classes thus have nonzero product whenever
\[
 \max(g-mj,0)+2(j+1)\le g,\qquad 0\le j\le a .
\]
Indeed the isogeny pullback multiplies the polarization class by $m^2$, and every power of the principal polarization of exponent at most $g$ is nonzero. The nonzero product has support in the triple intersection, which is therefore nonempty, even for the linked translates specified by $L$ and $K/L$. With $j=1$ the inequality holds for $g\ge4$, $m\ge4$. With $j=2$ it holds for $g\ge6$, $m=3$. \cite[Introduction]{Debarre}

In degree two the only unit slope is $-1\bmod 2$ and the canonical subdivisor argument applies for every $g$. For $m=3$, $g=4,5$, a remaining mixed cycle needs slope $1$ and a saturated common line $J$ as before. If $j=\deg J>0$, the ratio $L/M$ lies in $W_{a-j}-W_{a-j}$. For general $L$ this requires $2(a-j)\ge g$, hence $j=1$ in these genera. Degree zero is excluded by the finite torsion translates. The possible $J$ with $J^3$ effective range over a locus of dimension at most three. For each fixed $J$ the sum of the divisors of the two nonzero maps to $L,M$ belongs to $|K/J^2|$, of dimension $g-4+h^0(J^2)$. Each divisor in this system has finitely many bipartitions of the prescribed degrees and each choice determines $L$. Where $h^0(J^2)=0$ the dimension of this incidence is at most $g-1$. Where $h^0(J^2)=1$ the possible $J$ lie in the finite pullback of $W_2$, restoring the same bound. If $h^0(J^2)=2$ then $C$ is hyperelliptic and $J^2$ is its unique degree-two pencil, so $J$ ranges over a finite set. No higher value occurs. These constructible strata therefore miss general $L$ and complete the classification.

For a given order $m\ge2$, the antidiagonal action of $\mu_m$ on a
nonzero fiber point is free. An invariant multiset of length $m$
meeting the complement of the origin is a single orbit. A mixed orbit
would have slope $-1$ and is excluded for general $L$ by the
canonical-divisor argument above. Thus in degree $m$ the finite
subgroup has only the two pure axes for general $L$, including degrees
where the lower projected equations alone admit mixed curves.

In the genus-three cubic the intermediate quadratic cut is consequently
the union $E_+\cup_BE_-$ of two vector bundles of rank four over
$B=\PP^i\times\PP^{12-i}$. On the first branch the product of the
eight cubic Thom classes restricts, with support in the zero section,
to $(3y)^4\operatorname{Th}(E_+)$, and on the second to
$(3x)^4\operatorname{Th}(E_-)$, where $x,y$ are the hyperplanes.
The Thom class on the moving axis comes from its own fiber coordinates;
the other four tests vanish identically on that axis. The two
restrictions divided by $3^4$ agree on $B$, with value $3^4x^4y^4$.
Mayer--Vietoris with support in $B$ identifies the degree-sixteen
classes on the union with the kernel of the difference of restrictions
(the preceding $H^{15}(B)$ is zero). Thus their common class on the
union is divisible by $3^4$. Cup product with the four quadratic
Thom classes, supported on the union, maps this divisibility into the
cohomology of the ambient cycle space with support in $B$. Formally
this is the module action on $i^!\ZZ$ for the closed inclusion $i$ of
the union, followed by composition of supports. The original
denominator $2^4 3^8$ consequently bounds the 3-part of every weight
coefficient by $3^4$ after cancellation. The outer coefficient
$x_{13,1}/324$ in Lemma~\ref{lem:outer} attains this bound.

\subsection{The genus of a mixed integral cubic}\label{sec:szego}
On the same genus-three base a stronger statement about irreducible
curves holds even before imposing any coefficient tests. For general
$L$, every cubic projecting nontrivially to both $L$ and $M$ has
normalization genus at most nine.

Let $F$ denote its smooth normalization, $h=g(F)$ and
$T=\ker(f_*\cO_F\to\cO_C)$, of degree $7-h$. The two trace-free
coordinates give $N^*\to T$, $N=L\oplus M$. If its generic rank is
one, the fiber values belong to a common line of degree at most zero,
since $L\ne M$ and $L/M\notin C-C$. Its spectral divisor in this
line has arithmetic genus at most seven by adjunction. Otherwise
the quotient has length $11-h$. Set $A=L^3/K_C$ and use a hat for
the complementary line $\widehat P=K_C/P$. The four summands of
$V=\operatorname{Sym}^3N\otimes K_C^{-1}=A\oplus L\oplus M\oplus\widehat A$
are acyclic for general $L$. In genus eleven the cubic specifying
the cover would be a section of $V$ by the triple-cover
correspondence \cite[Sections~3--4]{Miranda}. In genus ten its
Tschirnhausen bundle is $W=\ker(N\to k(p))$ at one point. The
binary cubic is now a section of $V(p)$ uniquely determined by its
principal part, which is pure cube in the kernel direction.
An axis direction would give a generic pure cube of zero discriminant,
impossible for a separable field extension.

For an acyclic degree-two line $P$ the residue-one Szeg\H{o} kernel
in $(P\boxtimes\widehat P)(\Delta)$ exists uniquely by the diagonal
sequence and K\"unneth \cite{BB}. In local complementary frames its
scalar expansion is
$(z-w)^{-1}+a_P(w)+b_P(w)(z-w)+\cdots$.
Interchanging the two arguments gives
$a_{\widehat P}=-a_P$, $b_{\widehat P}=b_P-a'_P$.
In frames $X$ of $L$ and $Y=dz/X$ of $M$ write a mixed direction
as $X+qY$, $q\ne0$. The four coefficients of the unique meromorphic
binary cubic with principal part $(X+qY)^3$ are the kernels
$s_A,3qs_L,3q^2s_M,q^3s_{\widehat A}$, with the common determinant
frame suppressed. Set $\ell=X+qY$. The coefficient of $\ell Y^2$
has constant term $3q^2(a_A-3a_L)$; that of $Y^3$ has constant term
$-2q^3(a_A-3a_L)$ and first-order term $-q^3(a_A-3a_L)'$.
Since $W$ is generated locally by $\ell,(z-w)Y$, membership in
$\operatorname{Sym}^3W\otimes K_C^{-1}(p)$ forces a multiple zero
of $\eta=(a_A-3a_L)dz$ at $p$. This is a global differential: the
changes in $a_Pdz$ are $-d\log\lambda_P+
\tfrac12d\log\phi'$ for a frame change $\lambda_P$ and coordinate
change $\phi$, and $\lambda_A=\lambda_L^3/\phi'$.

For completeness the general differential so obtained has only simple
zeros, including on a hyperelliptic base. On the smooth locus of
$\Theta=W_2$, write $L_0=\cO(x+y)$ with unique generators
$v_0,w_0$ of $L_0,\widehat L_0$. The product $v_0w_0$ annihilates
the tangent of $\Theta$ by differentiating the Abel map. The Gauss
map is generically finite since every differential has only finitely
many degree-two subdivisors; its image is therefore dense. A general
such product has four simple zeros. The map $P:L\mapsto L^3/K_C$
does not preserve $\Theta$. If it did, multiplication by three on
tangents would leave the Gauss map invariant and force
$\deg(P|_\Theta)=1$. The projection formula on the surface $\Theta$
gives instead $9^2$ using the action on the polarization. We may thus
choose $L_0$ with $P(L_0)$ acyclic. On an algebraic transverse arc
through $L_0$, multiply the normalized kernel by the smallest
power $t^e$ giving a regular specialization. Here $e>0$, since
$h^0(L_0(w))=1$ for general $w$ and no section has the requisite
simple pole. The nonzero limit has zero residue and is therefore
a nonzero multiple of $v_0\boxtimes w_0$. The kernel for $P(L_0)$
stays regular by acyclicity. On taking finite parts along the diagonal
$t^e\eta$ consequently tends to a nonzero multiple of $v_0w_0$.
Algebraicity of the kernels on the acyclic locus and closedness of the
multiple-zero locus prove the assertion for general $L$.

In particular a reduced integral cubic on an axis of normalization
genus at least ten has no nearby mixed cubics. Indeed nearby cycles
are integral of multiplicity one by properness of Chow addition,
and a fixed normalization genus below ten in a convergent sequence
would give a limiting stable map with a positive component birational
to the central curve, which requires genus at least ten. The two axes
are closed and meet only in the zero-section multiple.

\section{Cubic pencils in genus five}\label{app:pencils}

Fix a smooth genus-five curve $D$. A genus-five cubic can move in a
pencil even after quotienting by dilation. At a cusp its normalization
has one fewer handle than the general member. Both genera occur in the
universal comparison.

\begin{theorem}\label{thm:pencil}
Let $N$ be a general stable bundle of rank two and determinant $K_D$.
In $\Tot N$ the reduced space of degree-three cycles is the union of
$6480$ lines and $32$ planes meeting only at $3[D]$. Dilation has
weight three on their parameters. A general determinant-one quadratic
shear along $D$ isolates $3[D]$ in a threefold $X$. For its unit $J_3$,
\[
 J_3-x_{5,3}=2160x_{16,1}+\tfrac{1280}{3}x_{18,1}
                            -\tfrac{32}{9}x_{19,1}.
\]
The difference of the connected fixed-target degree-three series in
the same order is $2160w^{15}+(1280/3)w^{17}-(32/9)w^{18}$.
\end{theorem}

\subsection{The pencils and their limits}
By Section~\ref{sec:evaluations} all nonzero cubics are $y^3=s$ on
saturated sublines of degrees one and two. For degree one
Corollary~\ref{cor:cubic-fiber-count} gives $6480$, with unique
reduced divisors by the dimension argument in Section~\ref{sec:covers}.
Each smooth cubic has genus 16. There are $2^5$ maximal degree-two
sublines, reduced for general $N$ \cite[Section~4]{Gillam}. One may
fix the determinant here since tensoring by a line reaches every
determinant of the same degree. At a degree-two subline $S$ the
projective extensions of $Q=K_DS^{-1}$ by $S$ have dimension seven,
so a proper closed locus of $S$ has image of dimension at most eleven
in the twelve-dimensional moduli of bundles. Thus $S^3$ is a general
degree-six line. Its two sections form a base-point-free pencil:
specialness requires $S^3=K_D(-B_2)$, a base point requires
$S^3=K_D(p-B_3)$, and a triple zero or two double zeros require
$S^3=\cO(3p+B_3)$ or $\cO(2p+2q+B_2)$, where $B_i$ is effective
of degree $i$. Each locus has dimension at most four. The pencil
therefore has twenty simple branch values by Riemann--Hurwitz. A
general cubic has six totally ramified points and genus 19. Each
exceptional member has one cusp $x^2=y^3$ and four other simple
zeros, so its normalization has genus 18. Saturation and the spectral
equation recover $(S,s)$ from the cycle; two general evaluations
bound all coefficients on each of the finitely many spaces. Hence
the bijection from their union is proper and gives its topology.

For the embedded cubic $F$ in $S\subset N$ the boundary in its
normal sequence is multiplication by $3y^2 e$, where
$e\in H^1(SQ^{-1})$ is the extension class. For a degree-two $S$
the nonzero connecting block is $H^0(Q)\to H^1(S)$, invertible by
acyclicity of $N$. Here $H^0(S^2)=0$ and $H^1(S^3)=0$. The finite
flat algebra $\cO\oplus S^{-1}\oplus S^{-2}$ also applies to the
embedded cusp. Thus throughout the punctured plane
\[
 H^0(N_{F/\Tot N})=H^0(S^3),\qquad H^1(N_{F/\Tot N})=
 H^1(QS^{-2})=H^0(S^3)^* .
\]
For a degree-one line $R$ the two additional connecting blocks are
dual to the pairing $(a,b)\mapsto\langle e,sab\rangle$ on
$H^0(K_DR^{-2})\times H^0(K_DR^{-3})$. Both spaces have dimension
two. Two general points off $\operatorname{div}s$ give independent
evaluations on both pencils, so a sum of their weighted evaluations
makes this form nondegenerate. Bad forms occupy a proper subset of
each extension fiber and the finite incidence excludes them for
general $N$. The remaining tangent and obstruction each have
dimension one.

Choose a determinant-one shear as in Section~\ref{sec:neighborhood}.
Its quadratic term on $S^2$, projected to $Q$, gives a fixed functional
in $H^0(S^3)^*$ over the entire punctured plane. On the normalization
of a cusp over $p$ the covering algebra has last eigenline
$S^{-2}(p)$; its obstruction is the restriction of the functional
to $H^0(S^3(-p))$. This is the line generated by the cusp section.
Principal parts span $H^1(QS^{-2})$, so finitely many open conditions
give a shear nonzero on all $640$ cusp lines, all 32 plane obstruction
spaces and all $6480$ isolated-line obstructions. At a smooth
embedded limit the relative Hilbert tangent space now has precisely
the coefficient directions. The actual coefficient space has that
dimension and hence fills its regular germ in the central target.

At a cuspidal limit the Hilbert tangent test alone would miss arcs
whose generic curves have smaller arithmetic genus. Suppose such a
noncentral arc exists. Curve selection in a projective Chow space
and properness of stable maps extend the normalizations over a DVR,
after finite base change. Write $t=cu^k+O(u^{k+1})$ for the target
parameter. The sole positive special component is birational to the
cusp; its normalization has the above eigenline decomposition.
Since $H^0(K_DS^{-3}(p))=0$ by Riemann--Roch and
$h^0(S^3(-p))=1$, and $H^0(K_DS^{-2})=0$ by $2$-rigidity and
Riemann--Roch, all sections of the pulled-back $Q$ on that component
come from $D$. Matching gives the same assertion on contracted pieces.
Over $A=\CC[[u]]/(u^k)$ all sections of $g^*Q$ for the map
$g:\mathcal C_A\to D_A$ still come from $D_A$: subtract the lift of
each successive special-fiber section and use source flatness.
The target is linear over $A$. The normal height in $g^*N$ has
zero trace since $H^0(D,N)=0$. This trace is defined even with
contracted components: $Rg_*\cO_{\mathcal C_A}$ is perfect by the
regular graph into $\mathcal C_A\times_A D_A$ and proper flat
perfect projection. Its rank on $D_A$ is three. Hence a pulled-back
quotient height has trace three times its base value, so the quotient
height vanishes modulo $u^k$. Modulo $u^{k+1}$ its first possible
coefficient would have to trivialize the projected quadratic
cocycle times $c$ on the special source. Its restriction to the
normalization is nonzero in $H^1(QS^{-2}(p))$, a contradiction.
This excludes noncentral arcs. Rescaling any sequence of displaced
cubic cycles to a fixed height, as in Proposition~\ref{prop:isolation},
would give one of the smooth or cuspidal limits just excluded.
Thus the sheared triple is isolated.

\subsection{The normalization of one cusp}
Near a nonzero cuspidal member the connected genus-18 map is its
smooth normalization, with trivial stabilizer. Any node, positive-genus
contracted component or graph loop would raise the arithmetic genus.
A contracted rational tree with one attachment would be unstable.
On an Artin base the map to $D$ stays finite flat. The same trace
and projected obstruction, now applied to consecutive powers of the
maximal ideal, force $t=0$ and the height into $S$. Its characteristic
equation is $y^3=s_A$ by $H^0(S)=H^0(S^2)=0$.
In local coordinates near the double zero the equation is
$X^2=Y^3+b$, with $b$ the critical value and $a$ the independent
parameter along the cusp ray. On the map the parametric Morse lemma
gives $Y=z^2+r$, with $r$ in the maximal ideal and the coordinate
sign fixed on the central fiber. The equality $X(z)^2=X(-z)^2$ and
regularity of $X(z)-X(-z)$ (reduce first to $2z^3$) give
$X=z h(z^2)$. Then
\[
 b=-r^3,\qquad h(v)^2=v^2+3rv+3r^2 .
\]
Weierstrass preparation forces $h=v+d$ with $2d=3r$ and $d^2=3r^2$.
Hence $r^2=b=0$. Conversely $(Y,X)=(z^2+r,z^3+3rz/2)$ for any
$r^2=0$ gives the required local normalization family. Its finite
free rank-three algebra agrees with the spectral algebra off the
marked disk by reduction and Nakayama. The two height functions
determine the gluing on the punctured disk uniquely. Thus the completed
map ring in a fiber frame is $\CC[[a,r]]/(r^2)$.
With two further smooth base directions this is an lci of the absolute
virtual dimension three. Its two-term obstruction theory has the same
rank as the lci cotangent complex. At any residue field the degree-zero
map is an isomorphism and the degree-minus-one map is surjective;
equality of ranks and derived Nakayama identify the two complexes.
The local virtual multiplicity is two.

\subsection{Moving the two evaluations}
Work over $B=\PP^2$ with tautological dual $M=\cO(1)$. Two general
fiber products give an injective map on each line, an isomorphism on
every plane and a Thom product representing $9x_{5,3}$. Translate by
two general small cubic forms in $M^3$. Move the shear parameter
graph to the metric dual $\tau$ of a linear section. Its zero line
$T$ has oriented normal index $-1$. The translations have nine common
zeros off $T$, accounting for $9J_3$. Their ratio on $T$ has degree
three, unramified over the finitely many line and cusp directions.
Each isolated line gives three negative genus-16 intersections.
Over each plane there remains a compact family above the whole of
$T$, with three incidences per cusp direction. At each of these the
two translations cut the amplitude and the ratio transversely, so the
genus-18 virtual degree is $-2$. Compactness of this homotopy follows
from the uniform lower bound on the pair of linear evaluations on
the unit link of the finite union of planes and lines, together with
isolation on the small noncentral fibers and the disk mass estimate.

To extend the cusp calculation to all source genera we deform each
whole pencil, since a genus-19 image can approach a cusp. Add a second
quadratic parameter spanning the two-dimensional embedded obstruction
with the first at every plane. On the graph take a second small
metric-dual linear section, cut off away from $T$ near the nine
marked points. Its zero on $T$ is a single general point. At this
point the two anti-complex normal factors have positive product,
yielding one positive smooth genus-19 intersection per plane.
For control of the other images take a relative projective completion.
The uniform mass bound bounds their ample degrees on the disk.
Nearby cycles at a nonzero plane member are integral with multiplicity
one, by proper Chow addition. In bounded degree their arithmetic
genera are bounded by the generic plane-projection genus. A Hilbert
limit with fundamental cycle $C$ surjects onto $\cO_C$ with a
finite-length kernel, giving $p_a(\mathrm{generic})\le p_a(C)$
with equality only for the reduced Hilbert limit. Stable-map limits
give the opposite lower bound on geometric genus by the normalization
of $C$. Thus near smooth members the arithmetic genus equals 19 and
the reduced Hilbert germ forces both parameters to vanish. Near a cusp
away from the common parameter zero the same argument excludes
arithmetic genus 19, while geometric genus is at least 18.
Every remaining small-parameter image there is therefore smooth
embedded of genus 18. Compactness gives uniform neighborhoods.

At these images the shears at fresh unramified sheets of
Lemma~\ref{lem:shear-jets} make the Hilbert chart regular; one may use
powers two through four of the separating covector, whose values are
nonzero and distinct, to preserve the first jet at the core. The new
three columns interpolate on the three sheets by the Vandermonde
determinant. A small smooth change of graph and translations gives
transverse isolated curves. Excision in source genus 18 preserves the
degree $-2$ in each cusp box. Proposition~\ref{prop:moving-embedded-curve}
now assigns the entire series $w^{17}$ to each signed genus-18 image,
including all contracted components. The same signed multiplicity-one
units arise by proper connected pullback. The complete unnormalized
cut is consequently
\[
 9J_3-3\cdot6480x_{16,1}-2\cdot60\cdot32x_{18,1}+32x_{19,1}.
\]
Dividing by nine gives the two identities in Theorem~\ref{thm:pencil}.
For the complex class, the two lower degrees agree by the empty-tuple
argument; hence the difference after subtraction is primitive and
killed by proper division. Its smooth image vanishes by the signed
calculation, so it vanishes by multiplicity-one detection.
Independently the handle recurrence gives
$9D_4\equiv3w^{17}+5w^{18}\pmod9$, consistent with the displayed signs.

\begin{thebibliography}{Wue26d}
\bibitem[BK05]{BK} V.~Batyrev and M.~Kreuzer,
\emph{Integral cohomology and mirror symmetry for Calabi--Yau 3-folds},
\href{https://arxiv.org/abs/math/0505432v1}{arXiv:math/0505432v1} (2005).
\bibitem[BF97]{BF} K.~Behrend and B.~Fantechi,
\emph{The intrinsic normal cone}, Invent. Math. \textbf{128} (1997), 45--88,
\href{https://arxiv.org/abs/alg-geom/9601010}{arXiv:alg-geom/9601010}.
\bibitem[BB03]{BB} D.~Ben-Zvi and I.~Biswas, \emph{Theta functions and
Szeg\H{o} kernels}, Int. Math. Res. Not. \textbf{2003} (2003), 1305--1340,
\href{https://doi.org/10.1155/S1073792803211181}{doi:10.1155/S1073792803211181}.
\bibitem[BP06]{BP06} J.~Bryan and R.~Pandharipande,
\emph{On the rigidity of stable maps to Calabi--Yau threefolds},
Geom. Topol. Monogr. \textbf{8} (2006), 97--104,
\href{https://doi.org/10.2140/gtm.2006.8.97}{doi:10.2140/gtm.2006.8.97}.
\bibitem[BP08]{BP} J.~Bryan and R.~Pandharipande,
\emph{The local Gromov--Witten theory of curves},
J. Amer. Math. Soc. \textbf{21} (2008), 101--136,
\href{https://arxiv.org/abs/math/0411037v3}{arXiv:math/0411037v3}.
\bibitem[Deb95]{Debarre} O.~Debarre, \emph{Minimal cohomology classes and Jacobians},
J. Algebraic Geom. \textbf{4} (1995), 321--335,
\href{https://arxiv.org/abs/alg-geom/9301002}{arXiv:alg-geom/9301002}.
\bibitem[DIW26]{DIW} A.~Doan, E.-N.~Ionel and T.~Walpuski,
\emph{The Gopakumar--Vafa finiteness conjecture},
Ann. of Math. (2) \textbf{204} (2026), 819--871,
\href{https://doi.org/10.4007/annals.2026.204.3.1}{doi:10.4007/annals.2026.204.3.1}.
Numbering follows \href{https://arxiv.org/abs/2103.08221v2}{arXiv:2103.08221v2}.
\bibitem[Gil11]{Gillam} W.~D.~Gillam,
\emph{Maximal subbundles, Quot schemes, and curve counting},
\href{https://arxiv.org/abs/1103.2169}{arXiv:1103.2169} (2011).
\bibitem[GP99]{GP} T.~Graber and R.~Pandharipande,
\emph{Localization of virtual classes}, Invent. Math. \textbf{135} (1999), 487--518,
\href{https://arxiv.org/abs/alg-geom/9708001}{arXiv:alg-geom/9708001}.
\bibitem[IP18]{IP} E.-N.~Ionel and T.~H.~Parker,
\emph{The Gopakumar--Vafa formula for symplectic manifolds},
Ann. of Math. (2) \textbf{187} (2018), 1--64,
\href{https://arxiv.org/abs/1306.1516v6}{arXiv:1306.1516v6}.
\bibitem[Mir85]{Miranda} R.~Miranda, \emph{Triple covers in algebraic geometry},
Amer. J. Math. \textbf{107} (1985), 1123--1158,
\href{https://www.math.colostate.edu/~miranda/preprints/TripleCoversInAG.pdf}{author's copy}.
\bibitem[Pan99]{Pand} R.~Pandharipande, \emph{Hodge integrals and
degenerate contributions}, Comm. Math. Phys. \textbf{208} (1999), 489--506,
\href{https://doi.org/10.1007/s002200050766}{doi:10.1007/s002200050766}.
\bibitem[PP14]{PP} R.~Pandharipande and A.~Pixton,
\emph{Gromov--Witten/pairs descendent correspondence for toric 3-folds},
Geom. Topol. \textbf{18} (2014), 2747--2821,
\href{https://doi.org/10.2140/gt.2014.18.2747}{doi:10.2140/gt.2014.18.2747}.
\bibitem[PT10]{PT} R.~Pandharipande and R.~P.~Thomas,
\emph{Stable pairs and BPS invariants}, J. Amer. Math. Soc. \textbf{23}
(2010), 267--297, \href{https://arxiv.org/abs/0711.3899}{arXiv:0711.3899}.
\bibitem[Par25]{Pardon} J.~Pardon,
\emph{Universally counting curves in Calabi--Yau threefolds},
\href{https://arxiv.org/abs/2308.02948v3}{arXiv:2308.02948v3} (2025).
\bibitem[Ray82]{Raynaud} M.~Raynaud, \emph{Sections des fibr\'es vectoriels sur une courbe}, Bull. Soc. Math. France \textbf{110} (1982), 103--125, \href{https://doi.org/10.24033/bsmf.1955}{doi:10.24033/bsmf.1955}.
\bibitem[STV15]{STV} T.~Sch\"urg, B.~To\"en and G.~Vezzosi,
\emph{Derived algebraic geometry, determinants of perfect complexes,
and applications to obstruction theories for maps and complexes},
J. Reine Angew. Math. \textbf{702} (2015), 1--40,
\href{https://doi.org/10.1515/crelle-2013-0037}{doi:10.1515/crelle-2013-0037}.
\bibitem[Sie99]{Siebert} B.~Siebert, \emph{Algebraic and symplectic
Gromov--Witten invariants coincide}, Ann. Inst. Fourier (Grenoble)
\textbf{49} (1999), 1743--1795,
\href{https://arxiv.org/abs/math/9804108}{arXiv:math/9804108}.
\bibitem[Wue26a]{WueCounter} B.~J.~Wuebben,
\emph{A counterexample to integral product exponents in universal curve counting}, 28 October 2026, revised 21 November 2026,
\href{https://doi.org/10.5281/zenodo.23750048}{doi:10.5281/zenodo.23750048}.

\bibitem[Wue26b]{WueProduct} B.~J.~Wuebben,
\emph{A geometric product formula for universal curve counting on Calabi--Yau threefolds}, 23 November 2026,
\href{https://doi.org/10.5281/zenodo.24455518}{doi:10.5281/zenodo.24455518}.

\bibitem[Wue26c]{WueBalanced} B.~J.~Wuebben,
\emph{Antidiagonal local curves and integral exponents in universal curve counting}, 11 November 2026, revised 22 November 2026,
\href{https://doi.org/10.5281/zenodo.24432700}{doi:10.5281/zenodo.24432700}.

\bibitem[Wue26e]{WueGenerators} B.~J.~Wuebben,
\emph{Integral generators for curve counting on threefolds}, 24 November 2026,
\href{https://doi.org/10.5281/zenodo.24510962}{doi:10.5281/zenodo.24510962}.

\end{thebibliography}

\end{document}
